%============================================================
%  Bd. 21 - Folgen
%============================================================

\documentclass[main.tex]{subfiles}

\ifSubfilesClassLoaded{
  \BandSetupStandalone{B21}
}{
}

\title{Bd. 21 - Folgen}
\author{Martin Kunze}
\date{}
\setcounter{file}{21}

\hypersetup{
  pdftitle={Bd. 21 - Folgen},
  pdfauthor={Martin Kunze}
}

\providecommand{\PosNatSeg}[1]{\mathbb{N}_{1,\ldots,#1}}
% Die folgenden Symbole werden in Band 19 lokal eingeführt. Die Fallbacks
% halten diesen Band auch im Standalone-Build lesbar und kollidieren wegen
% \providecommand nicht mit den dort bereits vorhandenen Definitionen.
\ifSubfilesClassLoaded{%
  \providecommand{\RealLe}{\mathrel{\leq_{\mathbb R}}}
  \providecommand{\RealLt}{\mathrel{<_{\mathbb R}}}
  \providecommand{\RealZero}{0_{\mathbb R}}
  \providecommand{\RealOne}{1_{\mathbb R}}
  \providecommand{\RealAdd}{\mathbin{\oplus}}
  \providecommand{\RealMul}{\mathbin{\odot}}
  \providecommand{\RealNeg}[1]{\mathop{\ominus}\!#1}
  \providecommand{\RealInv}[1]{\operatorname{inv}_{\mathbb R}\!\left(#1\right)}
  \providecommand{\RealAbs}[1]{\left\lvert #1\right\rvert_{\mathbb R}}
  \providecommand{\RealDist}[2]{d_{\mathbb R}\!\left(#1,#2\right)}
}{%
}

\begin{document}

\title{Bd. 21 - Folgen}
\author{Martin Kunze}
\date{}
\hypersetup{pageanchor=false}
\maketitle
\hypersetup{pageanchor=true}
\setcounter{tocdepth}{3}
\tableofcontents
\setcounter{file}{21}
\renewcommand{\FormulaBandID}{21}

% Band 21 verwendet dieselben festen Beweisspalten wie Band 27.
% Feste, oben ausgerichtete Absatzspalten halten die Begründungen dennoch
% durchgehend an derselben Stelle; sowohl Formel als auch Begründung dürfen
% innerhalb ihrer Spalte regulär umbrechen.
\let\BandTwentyOneOriginalProofStepWideStar\proofstepwidestar
\let\BandTwentyOneOriginalRuleRefWithArg\RuleRefWithArg
\ExplSyntaxOn
\cs_set_eq:NN
  \mx_band_twentyone_original_formula_ref_suffix:n
  \mx_formula_ref_suffix:n
\NewDocumentCommand{\BandTwentyOneBreakArgs}{m}
  {
    \tl_set:Nn \l_tmpa_tl {#1}
    \tl_replace_all:Nnn \l_tmpa_tl {,} {,\allowbreak}
    \tl_use:N \l_tmpa_tl
  }
\ExplSyntaxOff
\makeatletter
\newcol@{G}[0]{>{\raggedleft\arraybackslash\scriptsize}p{.10\linewidth}}
\makeatother
\RenewDocumentCommand{\proofstepwidestar}{O{} m m}{%
  \stepcounter{proofstepnr}%
  \BandTwentyOneBreakArgs{#1} & \theproofstepnr
     & \multicolumn{3}{@{}>{\raggedright\arraybackslash}p{.59\linewidth}@{}}{%
         \ensuremath{#2}}%
     & \multicolumn{1}{@{\hspace{.4em}}>{\raggedright\arraybackslash}p{.23\linewidth}@{}}{%
         #3}\\%
}
\ExplSyntaxOn
\cs_set_protected:Npn \mx_formula_ref_suffix:n #1
  {
    \IfNoValueF{#1}
      {\tl_if_blank:nF {#1}
        {\allowbreak (\BandTwentyOneBreakArgs{#1})}}
  }
\ExplSyntaxOff
\renewcommand{\RuleRefWithArg}[3]{%
  \ensuremath{\hyperref[rule:#1]{#2}\allowbreak
    (\BandTwentyOneBreakArgs{#3})}}

\chapter{Einleitung}

Eine Folge ist keine neue Art mathematischer Gegenstand, sondern eine
Funktion, deren Definitionsmenge als Indexmenge gelesen wird. Diese Sichtweise
trennt drei Dinge sauber voneinander: die Indexmenge, die Zielmenge und die
tatsächlich angenommene Bildmenge. Insbesondere darf aus einer Schreibweise
\(a\colon I\to M\) nicht ohne einen Surjektivitätsbeweis geschlossen werden,
dass jedes Element von \(M\) als Folgenglied vorkommt.

Der Band beginnt deshalb mit beliebig indizierten Familien. Unendliche Folgen
\(\mathbb N\to M\), endliche Folgen und die traditionelle Schreibweise
\((a_i)_{i\in I}\) erscheinen anschließend als Spezialfälle. Für reelle
Folgen werden Konvergenz, Grenzwerte, Beschränktheit, die elementaren
Grenzwertregeln und das Cauchy-Kriterium entwickelt. Die Bände~22 bis~26
verwenden endliche Folgen anschließend als Wege in gerichteten und
ungerichteten Graphen und bauen daraus den axiomatischen Baumbegriff auf.
Band~43 löst die Abstandsargumente von der besonderen Ordnung und Arithmetik
der reellen Zahlen und entwickelt daraus metrische Räume und metrische
Vollständigkeit. Das letzte Kapitel dieses Bandes stellt die rein folgen- und
mengentheoretischen Familien bereit, die in der Mogiljanskaja-Konstruktion der
späteren Halbgruppentheorie verwendet werden.

\chapter{Indizierte Familien}

\section{Familien als Funktionen}

\FormulaDefDeltaK[Indizierte Familie]{%
  a\colon I\to M\vdash
  (a_i)_{i\in I}\coloneqq a\dsep
  i\in I\vdash a_i\coloneqq a(i)%
}{IndexedFamilyDef}{%
  \DeltaRow{Mengen}{I\dsep M\dsep i}
  \DeltaRow{Funktionen}{a}
  \DeltaRow{Neue Schreibweisen}{(a_i)_{i\in I}\dsep a_i}
}

Die Menge \(I\) heißt \textbf{Indexmenge}, \(M\) heißt \textbf{Zielmenge}
der Familie, und \(a_i\) heißt das zum Index \(i\) gehörende
\textbf{Familienglied}. Die Klammer \((a_i)_{i\in I}\) bezeichnet die ganze
Funktion und nicht die ungeordnete Menge ihrer Werte. Verschiedene Indizes
dürfen daher dasselbe Familienglied tragen.

Für die folgenden Definitionen gilt die direkte Familienkonvention. Sei
\(y\) ein bereits gebildeter Term mit ausgezeichnetem Index \(i\). Sobald
\(i\in I\vdash y\in M\) bewiesen ist, darf ein frisches Familiensymbol \(a\)
unmittelbar durch
\[
  a\colon I\to M\dsep i\in I\vdash a_i\coloneqq y
\]
eingeführt werden. Der Satz
\FormulaRefAuto{TypedFunctionDefinitionPrinciple}[thm] aus Band~05 liefert die
eindeutig bestimmte Funktion \(a\colon I\to M\) mit \(a(i)=y\), und nach
\FormulaRefAuto{IndexedFamilyDef}[def] ist dies genau die Gleichung \(a_i=y\).
In einem solchen Definitionsblock ist daher nur noch die Typisierung des
rechten Terms nachzuweisen; ein eigener Existenzsatz wird nicht wiederholt.

\FormulaThmDeltaK[Automatische Typisierung eines Familiengliedes]{%
  a\colon I\to M\dsep i\in I\vdash a_i\in M%
}{IndexedFamilyCoordinateTyping}{%
  \DeltaRow{Mengen}{I\dsep M\dsep i}
  \DeltaRow{Funktionen}{a}
}
\begin{tabproofsplitwide}
  \proofpartwide{Typisierung}
    \proofstepwidestar[1]{%
      a\colon I\to M}{%
      \rA}
    \proofstepwidestar[2]{%
      i\in I}{%
      \rA}
    \proofstepwidestar[1,2]{%
      a(i)\in M}{%
      \FormulaRefAuto{F\colon A\to B\dsep x\in A\vdash F(x)\in B}[thm]{1,2}}
    \proofstepwidestar[1,2]{%
      a_i=a(i)}{%
      \FormulaRefAuto{IndexedFamilyDef}[def]{1,2}}
    \proofstepwidestar[1,2]{%
      a_i\in M}{%
      \rIE{4,3}}
  \closeproofpartwide
\end{tabproofsplitwide}

\FormulaThmDeltaK[Erweiterung der Zielmenge einer Familie]{%
  a\colon I\to M\dsep M\subseteq N\vdash a\colon I\to N%
}{IndexedFamilyCodomainExtension}{%
  \DeltaRow{Mengen}{I\dsep M\dsep N}
  \DeltaRow{Funktionen}{a}
}
\begin{tabproofsplitwide}
  \proofpartwide{Zielmengenerweiterung}
    \proofstepwidestar[1]{%
      a\colon I\to M}{%
      \rA}
    \proofstepwidestar[2]{%
      M\subseteq N}{%
      \rA}
    \proofstepwidestar[1,2]{%
      a\colon I\to N}{%
      \FormulaRefAuto{F\colon A\to B\dsep B\subseteq C\vdash F\colon A\to C}[thm]{1,2}}
  \closeproofpartwide
\end{tabproofsplitwide}

\section{Bildmenge und Wertemenge}

\FormulaDefDeltaK[Wertemenge einer Familie]{%
  a\colon I\to M\vdash
  \operatorname{Bild}(a)\coloneqq a[I]%
}{IndexedFamilyImageDef}{%
  \DeltaRow{Mengen}{I\dsep M\dsep a[I]}
  \DeltaRow{Funktionen}{a}
  \DeltaRow{Neue Schreibweise}{\operatorname{Bild}(a)}
}

\FormulaThmDeltaK[Elementkriterium der Wertemenge]{%
  a\colon I\to M\vdash
  x\in\operatorname{Bild}(a)
  \leftrightarrow
  \exists i\in I\,(x=a_i)%
}{IndexedFamilyImageMembership}{%
  \DeltaRow{Mengen}{I\dsep M\dsep x\dsep i}
  \DeltaRow{Funktionen}{a}
}
\begin{tabproofsplitwide}
  \proofpartwide{Elementkriterium}
    \proofstepwidestar[1]{%
      a\colon I\to M}{%
      \rA}
    \proofstepwidestar[1]{%
      \operatorname{Bild}(a)=a[I]}{%
      \FormulaRefAuto{IndexedFamilyImageDef}[def]{1}}
    \proofstepwidestar[1]{%
      x\in a[I]\leftrightarrow\exists i\in I\,(x=a(i))}{%
      \FormulaRefAuto{F\colon A\to B\vdash y\in F[A]\leftrightarrow\exists x\in A\,y=F(x)}[thm]{1}}
    \proofstepwidestar[1]{%
      \forall i\in I\,(a_i=a(i))}{%
      \FormulaRefAuto{IndexedFamilyDef}[def]{1}}
    \proofstepwidestar[1]{%
      x\in\operatorname{Bild}(a)\leftrightarrow\exists i\in I\,(x=a_i)}{%
      \rIE{2,3},\ \rUE{4}}
  \closeproofpartwide
\end{tabproofsplitwide}

\FormulaThmDeltaK[Wertemenge liegt in der Zielmenge]{%
  a\colon I\to M\vdash \operatorname{Bild}(a)\subseteq M%
}{IndexedFamilyImageSubsetCodomain}{%
  \DeltaRow{Mengen}{I\dsep M}
  \DeltaRow{Funktionen}{a}
}
\begin{tabproofsplitwide}
  \proofpartwide{Inklusion}
    \proofstepwidestar[1]{%
      a\colon I\to M}{%
      \rA}
    \proofstepwidestar[1]{%
      a[I]\subseteq M}{%
      \FormulaRefAuto{F\colon A\to B\vdash F[A]\subseteq B}[thm]{1}}
    \proofstepwidestar[1]{%
      \operatorname{Bild}(a)=a[I]}{%
      \FormulaRefAuto{IndexedFamilyImageDef}[def]{1}}
    \proofstepwidestar[1]{%
      \operatorname{Bild}(a)\subseteq M}{%
      \rIE{3,2}}
  \closeproofpartwide
\end{tabproofsplitwide}

\FormulaThmDeltaK[Surjektive Familien und ihre Wertemenge]{%
  a\colon I\to M\vdash
  \bigl(a\colon I\sur M\bigr)
  \leftrightarrow
  \operatorname{Bild}(a)=M%
}{IndexedFamilySurjectiveIffImage}{%
  \DeltaRow{Mengen}{I\dsep M}
  \DeltaRow{Funktionen}{a}
}
\begin{tabproofsplitwide}
  \proofpartwide{Surjektivitätskriterium}
    \proofstepwidestar[1]{%
      a\colon I\to M}{%
      \rA}
    \proofstepwidestar[1]{%
      (a\colon I\sur M)\leftrightarrow a[I]=M}{%
      \FormulaRefAuto{SurjectiveIffImageEqualsCodomain}[thm]{1}}
    \proofstepwidestar[1]{%
      \operatorname{Bild}(a)=a[I]}{%
      \FormulaRefAuto{IndexedFamilyImageDef}[def]{1}}
    \proofstepwidestar[1]{%
      (a\colon I\sur M)\leftrightarrow\operatorname{Bild}(a)=M}{%
      \rIE{3,2}}
  \closeproofpartwide
\end{tabproofsplitwide}

Die Zielmenge gehört zur Typisierung, die Wertemenge wird dagegen durch die
Glieder bestimmt. So kann dieselbe Familie zunächst als \(a\colon I\to M\)
und nach einer Zielmengenerweiterung als \(a\colon I\to N\) gelesen werden,
ohne dass sich \(\operatorname{Bild}(a)\) ändert.

\section{Gleichheit, Einschränkung und Umindizierung}

\FormulaThmDeltaK[Extensionalität indizierter Familien]{%
  a,b\colon I\to M\vdash
  a=b\leftrightarrow\forall i\in I\,(a_i=b_i)%
}{IndexedFamilyExtensionality}{%
  \DeltaRow{Mengen}{I\dsep M\dsep i}
  \DeltaRow{Funktionen}{a\dsep b}
}
\begin{tabproofsplitwide}
  \proofpartwide{Beide Richtungen}
    \proofstepwidestar[1]{%
      a,b\colon I\to M}{%
      \rA}
    \proofstepwidestar[2]{%
      a=b}{%
      \rA}
    \proofstepwidestar[3]{%
      i\in I}{%
      \rA}
    \proofstepwidestar[1,2,3]{%
      a_i=b_i}{%
      \rIE{2,\rII},\ \FormulaRefAuto{IndexedFamilyDef}[def]{1,3}}
    \proofstepwidestar[1,2]{%
      \forall i\in I\,(a_i=b_i)}{%
      \rUI{\rRI{3,4}}}
    \proofstepwidestar[1]{%
      a=b\rightarrow\forall i\in I\,(a_i=b_i)}{%
      \rRI{2,5}}
    \proofstepwidestar[7]{%
      \forall i\in I\,(a_i=b_i)}{%
      \rA}
    \proofstepwidestar[1,7]{%
      \forall i\in I\,(a(i)=b(i))}{%
      \FormulaRefAuto{IndexedFamilyDef}[def]{1,7}}
    \proofstepwidestar[1,7]{%
      a=b}{%
      \FormulaRefAuto{FunctionExtensionality}[thm]{1,8}}
    \proofstepwidestar[1]{%
      \forall i\in I\,(a_i=b_i)\rightarrow a=b}{%
      \rRI{7,9}}
    \proofstepwidestar[1]{%
      a=b\leftrightarrow\forall i\in I\,(a_i=b_i)}{%
      \rLRI{6,10}}
  \closeproofpartwide
\end{tabproofsplitwide}

\FormulaDefDeltaK[Einschränkung einer Familie]{%
  a\colon I\to M\dsep J\subseteq I\vdash
  (a_i)_{i\in J}\coloneqq a\restriction_J%
}{IndexedFamilyRestrictionDef}{%
  \DeltaRow{Mengen}{I\dsep J\dsep M\dsep i}
  \DeltaRow{Funktionen}{a\dsep a\restriction_J}
}

\FormulaThmDeltaK[Eigenschaften der eingeschränkten Familie]{%
  \begin{aligned}[t]
  &a\colon I\to M\dsep J\subseteq I\vdash
    a\restriction_J\colon J\to M\dsep{}\\[-2pt]
  &\quad\forall i\in J\,\bigl((a\restriction_J)_i=a_i\bigr)\dsep{}\\[-2pt]
  &\quad\operatorname{Bild}(a\restriction_J)\subseteq
    \operatorname{Bild}(a)
  \end{aligned}%
}{IndexedFamilyRestrictionFacts}{%
  \DeltaRow{Mengen}{I\dsep J\dsep M\dsep i}
  \DeltaRow{Funktionen}{a\dsep a\restriction_J}
}
\begin{tabproofsplitwide}
  \proofpartwide{Typisierung, Koordinaten und Bild}
    \proofstepwidestar[1]{%
      a\colon I\to M}{%
      \rA}
    \proofstepwidestar[2]{%
      J\subseteq I}{%
      \rA}
    \proofstepwidestar[1,2]{%
      a\restriction_J\colon J\to M}{%
      \FormulaRefAuto{RestrictionDef}[def]{1,2}}
    \proofstepwidestar[4]{%
      i\in J}{%
      \rA}
    \proofstepwidestar[1,2,4]{%
      (a\restriction_J)(i)=a(i)}{%
      \FormulaRefAuto{RestrictionDef}[def]{1,2,4}}
    \proofstepwidestar[1,2,4]{%
      (a\restriction_J)_i=a_i}{%
      \FormulaRefAuto{IndexedFamilyDef}[def]{1,3,4,5}}
    \proofstepwidestar[1,2]{%
      \forall i\in J\,((a\restriction_J)_i=a_i)}{%
      \rUI{\rRI{4,6}}}
    \proofstepwidestar[1,2]{%
      (a\restriction_J)[J]=a[J]}{%
      \FormulaRefAuto{C\subseteq A\dsep F\colon A\to B\vdash(F\restriction_C)[C]=F[C]}[thm]{2,1}}
    \proofstepwidestar[1,2]{%
      a[J]\subseteq a[I]}{%
      \FormulaRefAuto{F\colon M\to N\dsep A\subseteq B\dsep B\subseteq M\vdash F[A]\subseteq F[B]}[thm]{1,2,\FormulaRefAuto{A\subseteq A}[thm]}}
    \proofstepwidestar[1,2]{%
      \operatorname{Bild}(a\restriction_J)\subseteq\operatorname{Bild}(a)}{%
      \FormulaRefAuto{IndexedFamilyImageDef}[def]{1,3},\ \rIE{8,9}}
  \closeproofpartwide
\end{tabproofsplitwide}

\FormulaDefDeltaK[Umindizierte Familie]{%
  a\colon I\to M\dsep \varphi\colon J\to I\vdash
  (a_{\varphi(j)})_{j\in J}\coloneqq a\circ\varphi%
}{IndexedFamilyReindexingDef}{%
  \DeltaRow{Mengen}{I\dsep J\dsep M\dsep j}
  \DeltaRow{Funktionen}{a\dsep\varphi\dsep a\circ\varphi}
}

\FormulaThmDeltaK[Bildmenge nach Umindizierung]{%
  a\colon I\to M\dsep \varphi\colon J\to I\vdash
  \operatorname{Bild}(a\circ\varphi)=a[\varphi[J]]
  \subseteq\operatorname{Bild}(a)%
}{IndexedFamilyReindexingImage}{%
  \DeltaRow{Mengen}{I\dsep J\dsep M}
  \DeltaRow{Funktionen}{a\dsep\varphi}
}
\begin{tabproofsplitwide}
  \proofpartwide{Bild der Komposition}
    \proofstepwidestar[1]{%
      a\colon I\to M}{%
      \rA}
    \proofstepwidestar[2]{%
      \varphi\colon J\to I}{%
      \rA}
    \proofstepwidestar[1,2]{%
      a\circ\varphi\colon J\to M}{%
      \FormulaRefAuto{CompositionDef}[def]{2,1}}
    \proofstepwidestar[2]{%
      \varphi[J]\subseteq I}{%
      \FormulaRefAuto{F\colon A\to B\vdash F[A]\subseteq B}[thm]{2}}
    \proofstepwidestar[1,2]{%
      x\in\operatorname{Bild}(a\circ\varphi)\leftrightarrow\exists j\in J\,(x=a(\varphi(j)))}{%
      \FormulaRefAuto{IndexedFamilyImageMembership}[thm]{3},\ \FormulaRefAuto{CompositionDef}[def]{2,1}}
    \proofstepwidestar[1,2]{%
      x\in a[\varphi[J]]\leftrightarrow\exists j\in J\,(x=a(\varphi(j)))}{%
      \FormulaRefAuto{C\subseteq A\dsep F\colon A\to B\vdash y\in F[C]\leftrightarrow\exists x\in C\,y=F(x)}[thm]{4,1},\ \FormulaRefAuto{IndexedFamilyImageMembership}[thm]{2}}
    \proofstepwidestar[1,2]{%
      x\in\operatorname{Bild}(a\circ\varphi)\leftrightarrow x\in a[\varphi[J]]}{%
      \rLRI{\rRE{\rLREb{6},\rLREa{5}},\rRE{\rLREb{5},\rLREa{6}}}}
    \proofstepwidestar[1,2]{%
      \operatorname{Bild}(a\circ\varphi)=a[\varphi[J]]}{%
      \text{Extensionalität}{\rUI{7}}}
    \proofstepwidestar[1,2]{%
      a[\varphi[J]]\subseteq a[I]}{%
      \FormulaRefAuto{F\colon M\to N\dsep A\subseteq B\dsep B\subseteq M\vdash F[A]\subseteq F[B]}[thm]{1,4,\FormulaRefAuto{A\subseteq A}[thm]}}
    \proofstepwidestar[1,2]{%
      \operatorname{Bild}(a\circ\varphi)=a[\varphi[J]]\subseteq\operatorname{Bild}(a)}{%
      8,\ \FormulaRefAuto{IndexedFamilyImageDef}[def]{1,9}}
  \closeproofpartwide
\end{tabproofsplitwide}

Ist \(\varphi\) injektiv, so heißt die umindizierte Familie eine
\textbf{Teilfamilie}. Eine bloße Einschränkung ist der Spezialfall, in dem
\(\varphi\) die Inklusion \(J\hookrightarrow I\) ist. Ist \(\varphi\)
surjektiv, dann hat die umindizierte Familie dieselbe Wertemenge wie die
Ausgangsfamilie.

\section{Injektive und surjektive Familien}

\FormulaDefDeltaK[Injektive Familie]{%
  a\colon I\to M\vdash
  \bigl((a_i)_{i\in I}\text{ ist injektiv}\bigr)
  \coloneqq \bigl(a\colon I\inj M\bigr)%
}{InjectiveIndexedFamilyDef}{%
  \DeltaRow{Mengen}{I\dsep M}
  \DeltaRow{Funktionen}{a}
}

\FormulaThmDeltaK[Indexkriterium für Injektivität]{%
  a\colon I\to M\vdash
  \bigl(a\colon I\inj M\bigr)
  \leftrightarrow
  \forall i,j\in I\,(a_i=a_j\rightarrow i=j)%
}{InjectiveIndexedFamilyCriterion}{%
  \DeltaRow{Mengen}{I\dsep M\dsep i\dsep j}
  \DeltaRow{Funktionen}{a}
}
\begin{tabproofsplitwide}
  \proofpartwide{Indexkriterium}
    \proofstepwidestar[1]{%
      a\colon I\to M}{%
      \rA}
    \proofstepwidestar[1]{%
      (a\colon I\inj M)\leftrightarrow\forall i,j\in I\,(a(i)=a(j)\rightarrow i=j)}{%
      \FormulaRefAuto{Injektive Funktion}[def]{1}}
    \proofstepwidestar[1]{%
      \forall i\in I\,(a_i=a(i))}{%
      \FormulaRefAuto{IndexedFamilyDef}[def]{1}}
    \proofstepwidestar[1]{%
      (a\colon I\inj M)\leftrightarrow\forall i,j\in I\,(a_i=a_j\rightarrow i=j)}{%
      2,\ \rUE{3}}
  \closeproofpartwide
\end{tabproofsplitwide}

\FormulaDefDeltaK[Surjektive Familie]{%
  a\colon I\to M\vdash
  \bigl((a_i)_{i\in I}\text{ durchläuft }M\bigr)
  \coloneqq \bigl(a\colon I\sur M\bigr)%
}{SurjectiveIndexedFamilyDef}{%
  \DeltaRow{Mengen}{I\dsep M}
  \DeltaRow{Funktionen}{a}
}

Eine bijektive Familie ist zugleich injektiv und surjektiv. Sie zählt jedes
Element ihrer Zielmenge genau einmal auf. Eine nichtinjektive Folge darf
dagegen Werte wiederholen, und eine nichtsurjektive Folge lässt Zielwerte aus.

\chapter{Endliche und unendliche Folgen}

\section{Unendliche Folgen}

\FormulaDefDeltaK[Unendliche Folge in einer Menge]{%
  a\colon\mathbb N\to M\vdash
  (a_n)_{n\in\mathbb N}\coloneqq a%
}{InfiniteSequenceDef}{%
  \DeltaRow{Mengen}{M\dsep n}
  \DeltaRow{Funktionen}{a}
  \DeltaRow{Neue Schreibweise}{(a_n)_{n\in\mathbb N}}
}

Unsere natürlichen Zahlen beginnen bei \(0\). Daher ist \(a_0\) das erste,
\(a_1\) das zweite und allgemein \(a_n\) das durch den Index \(n\)
bezeichnete Glied. Die drei Schreibweisen
\[
  a\colon\mathbb N\to M,
  \qquad (a_n)_{n\in\mathbb N},
  \qquad (a_0,a_1,a_2,\ldots)
\]
bezeichnen dieselbe Funktion, betonen aber unterschiedliche Aspekte.

\FormulaDefDeltaK[Verschobene Folge und Restfolge]{%
  a\colon\mathbb N\to M\dsep r\in\mathbb N\vdash
  \begin{aligned}[t]
    &a^{\langle r\rangle}\colon\mathbb N\to M,\\[-2pt]
    &\forall n\in\mathbb N\,
      \bigl(a^{\langle r\rangle}_n\coloneqq a_{n+r}\bigr)
  \end{aligned}%
}{SequenceTailDef}{%
  \DeltaRow{Mengen}{M\dsep n\dsep r}
  \DeltaRow{Funktionen}{a\dsep a^{\langle r\rangle}}
}

Die Restfolge lässt genau die ersten \(r\) Glieder aus. Ihre Wertemenge ist
im Allgemeinen nur eine Teilmenge der ursprünglichen Wertemenge.

\FormulaDefDeltaK[Teilfolge einer unendlichen Folge]{%
  a\colon\mathbb N\to M\dsep
  \varphi\colon\mathbb N\to\mathbb N\dsep
  \forall n\in\mathbb N\,
    \bigl(\varphi(n)<\varphi(\Succ(n))\bigr)
  \vdash
  (a_{\varphi(n)})_{n\in\mathbb N}\coloneqq a\circ\varphi%
}{SubsequenceDef}{%
  \DeltaRow{Mengen}{M\dsep n}
  \DeltaRow{Funktionen}{a\dsep\varphi\dsep a\circ\varphi}
}

Die strikte Zunahme der Indexfolge verhindert Wiederholungen und eine
Umkehrung der ursprünglichen Reihenfolge. Jede Restfolge ist eine Teilfolge;
die zugehörige Indexfunktion lautet \(\varphi(n)=n+r\).

\section{Endliche Folgen mit nullbasierten Indizes}

Für \(n\in\mathbb N\) ist
\(\NatSeg{n}=\{i\in\mathbb N\mid i<n\}\) der Anfangsabschnitt mit genau
\(n\) Elementen. Insbesondere ist \(\NatSeg{0}=\varnothing\), während
\(\NatSeg{1}=\{0\}\) gilt.

\FormulaDefDeltaK[Nullbasierte endliche Folge]{%
  n\in\mathbb N\dsep a\colon\NatSeg{n}\to M\vdash
  (a_i)_{i<n}\coloneqq(a_i)_{i\in\NatSeg{n}}%
}{ZeroBasedFiniteSequenceDef}{%
  \DeltaRow{Mengen}{M\dsep n\dsep i\dsep\NatSeg{n}}
  \DeltaRow{Funktionen}{a}
}

Eine nullbasierte Folge der Länge \(n\) wird für \(n>0\) auch als
\[
  (a_0,a_1,\ldots,a_{n-1})
\]
geschrieben. Der größte Index ist also \(n-1\), die Anzahl der Glieder aber
\(n\). Bei \(n=0\) gibt es keinen Ausdruck \(a_{n-1}\); die Folge ist dann
die eindeutig bestimmte leere Funktion \(\varnothing\to M\).

\subsection{Verschobene Graphen endlicher Folgen}

Bei einer Indexverschiebung wird nur die erste Koordinate des
Funktionsgraphen verändert. Die folgenden Grundlagen gelten für
beliebige endliche Folgen.

\FormulaThmDeltaK[Eindeutige Existenz eines verschobenen Folgengraphen]{%
  \exists!G\,\forall p\,\Bigl(
    p\in G\leftrightarrow
    \exists j\,(j\in J\land p=(r+j,v(j)))\Bigr)%
}{ShiftedWordGraphUniqueExists}{%
  \DeltaRow{Mengen}{G\dsep J\dsep j\dsep p\dsep q\dsep r}
  \DeltaRow{Funktionensymbole}{v}
}
\begin{tabproofsplitwide}
  \proofpartwide{Instanz des Termschemas}
    \proofstepwidestar[]{%
      \begin{aligned}[t]&\exists!G\,\forall p\,\bigl(p\in G\leftrightarrow{}\\[-2pt]&\qquad\exists j\,(j\in J\land p=(r+j,v(j)))\bigr)\end{aligned}}{%
      \FormulaRefAuto{TermImageSetUniqueExists}[thm]}
  \closeproofpartwide
\end{tabproofsplitwide}

\FormulaDefDeltaK[Verschobener Folgengraph]{%
  \{(r+j,v(j))\mid j\in J\}\coloneqq
  \iota G\,\forall p\,\Bigl(
    p\in G\leftrightarrow
    \exists j\,(j\in J\land p=(r+j,v(j)))\Bigr)%
}{ShiftedWordGraphDef}{%
  \DeltaRow{Mengen}{G\dsep J\dsep j\dsep p\dsep r}
  \DeltaRow{Funktionensymbole}{v}
}

\FormulaThmDeltaK[Funktionalität eines verschobenen Folgengraphen]{%
  \begin{aligned}[t]
    &r\in\mathbb N\dsep s\in\mathbb N\dsep
      v\colon\NatSeg s\to A\dsep{}\\[-2pt]
    &(q,a)\in\{(r+j,v(j))\mid j\in\NatSeg s\}\dsep
      (q,b)\in\{(r+j,v(j))\mid j\in\NatSeg s\}\vdash a=b
  \end{aligned}%
}{ShiftedWordGraphFunctional}{%
  \DeltaRow{Mengen}{A\dsep a\dsep b\dsep j\dsep k\dsep q\dsep r\dsep s}
  \DeltaRow{Funktionen}{v}
}
\begin{tabproofsplitwide}
  \proofpartwide{Zeugen, Indexgleichheit und Wertgleichheit}
    \proofstepwidestar[1]{%
      r\in\mathbb N}{%
      \rA}
    \proofstepwidestar[2]{%
      s\in\mathbb N}{%
      \rA}
    \proofstepwidestar[3]{%
      v\colon\NatSeg s\to A}{%
      \rA}
    \proofstepwidestar[4]{%
      (q,a)\in\{(r+j,v(j))\mid j\in\NatSeg s\}}{%
      \rA}
    \proofstepwidestar[5]{%
      (q,b)\in\{(r+j,v(j))\mid j\in\NatSeg s\}}{%
      \rA}
    \proofstepwidestar[4]{%
      \exists j\in\NatSeg s\,((q,a)=(r+j,v(j)))}{%
      \FormulaRefAuto{ShiftedWordGraphDef}[def]{4}}
    \proofstepwidestar[5]{%
      \exists k\in\NatSeg s\,((q,b)=(r+k,v(k)))}{%
      \FormulaRefAuto{ShiftedWordGraphDef}[def]{5}}
    \proofstepwidestar[8]{%
      j\in\NatSeg s\land(q,a)=(r+j,v(j))}{%
      \rA}
    \proofstepwidestar[9]{%
      k\in\NatSeg s\land(q,b)=(r+k,v(k))}{%
      \rA}
    \proofstepwidestar[2,8]{%
      j\in\mathbb N}{%
      \FormulaRefAuto{PeanoNatSegElementNatural}[thm]{2,\rAEa{8}}}
    \proofstepwidestar[2,9]{%
      k\in\mathbb N}{%
      \FormulaRefAuto{PeanoNatSegElementNatural}[thm]{2,\rAEa{9}}}
    \proofstepwidestar[8]{%
      q=r+j\land a=v(j)}{%
      \FormulaRefAuto{(a,b)=(c,d)\eqvdash a=c\land b=d}[thm]{\rAEb{8}}}
    \proofstepwidestar[9]{%
      q=r+k\land b=v(k)}{%
      \FormulaRefAuto{(a,b)=(c,d)\eqvdash a=c\land b=d}[thm]{\rAEb{9}}}
    \proofstepwidestar[8,9]{%
      r+j=r+k}{%
      \FormulaRefAuto{a=b\dsep a=c\vdash b=c}[thm]{\rAEa{12},\rAEa{13}}}
    \proofstepwidestar[1,2,8,9]{%
      j=k}{%
      \FormulaRefAuto{PeanoAddLeftCancellation}[thm]{1,10,11,14}}
    \proofstepwidestar[1,2,3,8,9]{%
      v(j)=v(k)}{%
      \rIE{15,\rII}}
    \proofstepwidestar[1,2,3,8,9]{%
      a=b}{%
      \rIE{\rAEb{12},\rIE{\rAEb{13},16}}}
    \proofstepwidestar[1,2,3,5,8]{%
      a=b}{%
      \rEE{7,9,17}}
    \proofstepwidestar[1,2,3,4,5]{%
      a=b}{%
      \rEE{6,8,18}}
  \closeproofpartwide
\end{tabproofsplitwide}

\section{Endliche Folgen mit positiven Indizes}

\FormulaDefDeltaK[Positive Indexmenge]{%
  n\in\mathbb N\vdash
  \PosNatSeg{n}\coloneqq
  \{i\in\mathbb N\mid 1\leq i\land i\leq n\}%
}{PositiveNaturalSegmentDef}{%
  \DeltaRow{Mengen}{n\dsep i\dsep\PosNatSeg{n}}
  \DeltaRow{Neue Schreibweise}{\PosNatSeg{n}}
}

\FormulaThmDeltaK[Ränder der positiven Indexmenge]{%
  n\in\mathbb N\vdash
  \PosNatSeg{0}=\varnothing\land
  \forall i\in\mathbb N\,
    \bigl(i\in\PosNatSeg{n}\leftrightarrow
      1\leq i\land i\leq n\bigr)%
}{PositiveNaturalSegmentBoundary}{%
  \DeltaRow{Mengen}{n\dsep i\dsep\PosNatSeg{n}}
}
\begin{tabproofsplitwide}
  \proofpartwideindR[Leerer Anfangsabschnitt]{%
    \PosNatSeg 0=\varnothing
  }{\PosNatSeg 0=\varnothing}[B21PositiveSegmentZero]
    \proofstepwidestar[1]{%
      i\in\PosNatSeg 0}{%
      \rA}
    \proofstepwidestar[1]{%
      i\in\mathbb N\land1\leq i\land i\leq0}{%
      \FormulaRefAuto{PositiveNaturalSegmentDef}[def]{\FormulaRefAuto{PeanoZero}[ax],1}}
    \proofstepwidestar[]{%
      1\in\mathbb N}{%
      \FormulaRefAuto{PeanoOneInN}[thm]}
    \proofstepwidestar[1]{%
      1\leq0}{%
      \FormulaRefAuto{PeanoLeqTransitive}[thm]{3,\rAEa{2},\FormulaRefAuto{PeanoZero}[ax],\rAEa{\rAEb{2}},\rAEb{\rAEb{2}}}}
    \proofstepwidestar[]{%
      0<1}{%
      \FormulaRefAuto{PeanoLtAddOne}[thm]{\FormulaRefAuto{PeanoZero}[ax]},\ \FormulaRefAuto{PeanoAddLeftZero}[thm]{3}}
    \proofstepwidestar[1]{%
      \bot}{%
      \FormulaRefAuto{PeanoOrderNoCycle}[thm]{\FormulaRefAuto{PeanoZero}[ax],3,5,4}}
    \proofstepwidestar[]{%
      \forall i\,i\notin\PosNatSeg 0}{%
      \rUI{\rCI{1,6}}}
    \proofstepwidestar[]{%
      \PosNatSeg 0=\varnothing}{%
      \FormulaRefAuto{\forall x\,x\notin A\vdash A=\varnothing}[thm]{7}}
  \closeproofpartwideindR

  \proofpartwide{Elementkriterium und Zusammenfassung}
    \proofstepwidestar[1]{%
      n\in\mathbb N}{%
      \rA}
    \proofstepwidestar[2]{%
      i\in\mathbb N}{%
      \rA}
    \proofstepwidestar[1,2]{%
      i\in\PosNatSeg n\leftrightarrow(1\leq i\land i\leq n)}{%
      \FormulaRefAuto{PositiveNaturalSegmentDef}[def]{1,2}}
    \proofstepwidestar[1]{%
      \forall i\in\mathbb N\,(i\in\PosNatSeg n\leftrightarrow1\leq i\land i\leq n)}{%
      \rUI{\rRI{2,3}}}
    \proofstepwidestar[1]{%
      \begin{aligned}[t]&\PosNatSeg 0=\varnothing\land{}\\[-2pt]&\forall i\in\mathbb N\,(i\in\PosNatSeg n\leftrightarrow1\leq i\land i\leq n)\end{aligned}}{%
      \rAI{\FormulaRefAuto{B21PositiveSegmentZero}[thm],4}}
  \closeproofpartwide
\end{tabproofsplitwide}

\FormulaDefDeltaK[Positiv indizierte endliche Folge]{%
  n\in\mathbb N\dsep a\colon\PosNatSeg{n}\to M\vdash
  (a_i)_{i=1}^{n}\coloneqq(a_i)_{i\in\PosNatSeg{n}}%
}{PositiveFiniteSequenceDef}{%
  \DeltaRow{Mengen}{M\dsep n\dsep i\dsep\PosNatSeg{n}}
  \DeltaRow{Funktionen}{a}
}

Für \(n>0\) lautet die ausgeschriebene Form
\[
  (a_1,a_2,\ldots,a_n).
\]
Auch diese Folge hat genau \(n\) Glieder. Die Schreibweise
\((a_0,\ldots,a_n)\) hätte dagegen \(n+1\) Glieder und darf nicht als Folge
der Länge \(n\) bezeichnet werden.

\FormulaDefDeltaK[Indexverschiebung zwischen den beiden Konventionen]{%
  n\in\mathbb N\vdash
  \begin{aligned}[t]
    &\lambda_n\colon\NatSeg{n}\to\PosNatSeg{n},\\[-2pt]
    &\forall i\in\NatSeg{n}\,
      \bigl(\lambda_n(i)\coloneqq\Succ(i)\bigr)
  \end{aligned}%
}{FiniteSequenceIndexShiftDef}{%
  \DeltaRow{Mengen}{n\dsep i\dsep\NatSeg{n}\dsep\PosNatSeg{n}}
  \DeltaRow{Funktionen}{\lambda_n}
}
\begin{tabproofsplitwide}
  \proofpartwide{Typisierung der Definitionsvorschrift}
    \proofstepwidestar[1]{%
      n\in\mathbb N}{%
      \rA}
    \proofstepwidestar[2]{%
      i\in\NatSeg n}{%
      \rA}
    \proofstepwidestar[1,2]{%
      i\in\mathbb N\land i<n}{%
      \FormulaRefAuto{PeanoNatSegMembership}[thm]{1,2}}
    \proofstepwidestar[1,2]{%
      \Succ(i)\in\mathbb N}{%
      \FormulaRefAuto{PeanoSuccClosure}[ax]{\rAEa{3}}}
    \proofstepwidestar[1,2]{%
      0\leq i}{%
      \FormulaRefAuto{PeanoZeroLeq}[thm]{\rAEa{3}}}
    \proofstepwidestar[1,2]{%
      0<\Succ(i)}{%
      \FormulaRefAuto{PeanoLeqSuccStrict}[thm]{\FormulaRefAuto{PeanoZero}[ax],\rAEa{3},5}}
    \proofstepwidestar[1,2]{%
      1\leq\Succ(i)}{%
      \FormulaRefAuto{PeanoLtAddOneLeq}[thm]{\FormulaRefAuto{PeanoZero}[ax],4,6},\ \FormulaRefAuto{PeanoAddLeftZero}[thm]{\FormulaRefAuto{PeanoOneInN}[thm]}}
    \proofstepwidestar[1,2]{%
      \Succ(i)\leq n}{%
      \FormulaRefAuto{PeanoLtAddOneLeq}[thm]{\rAEa{3},1,\rAEb{3}},\ \FormulaRefAuto{PeanoAddOneSucc}[thm]{\rAEa{3}}}
    \proofstepwidestar[1,2]{%
      \Succ(i)\in\PosNatSeg n}{%
      \FormulaRefAuto{PositiveNaturalSegmentDef}[def]{1,4,7,8}}
    \proofstepwidestar[1]{%
      \forall i\in\NatSeg n\,(\Succ(i)\in\PosNatSeg n)}{%
      \rUI{\rRI{2,9}}}
  \closeproofpartwide
\end{tabproofsplitwide}

\FormulaThmDeltaK[Äquivalenz der beiden endlichen Indexkonventionen]{%
  n\in\mathbb N\vdash
  \lambda_n\colon\NatSeg{n}\bij\PosNatSeg{n}\dsep
  j\in\PosNatSeg{n}\vdash\lambda_n^{-1}(j)=\Pred(j)%
}{FiniteSequenceIndexShiftBijective}{%
  \DeltaRow{Mengen}{n\dsep i\dsep j\dsep\NatSeg{n}\dsep\PosNatSeg{n}}
  \DeltaRow{Funktionen}{\lambda_n\dsep\lambda_n^{-1}}
}
\begin{tabproofsplitwide}
  \proofpartwideindR[Injektivität und Surjektivität]{%
    n\in\mathbb N\vdash\lambda_n\colon\NatSeg n\bij\PosNatSeg n
  }{n\in\mathbb N\vdash\lambda_n\colon\NatSeg n\bij\PosNatSeg n}[B21FiniteIndexShiftBijection]
    \proofstepwidestar[1]{%
      n\in\mathbb N}{%
      \rA}
    \proofstepwidestar[1]{%
      \lambda_n\colon\NatSeg n\to\PosNatSeg n}{%
      \FormulaRefAuto{FiniteSequenceIndexShiftDef}[def]{1}}
    \proofstepwidestar[3]{%
      (i\in\NatSeg n\land k\in\NatSeg n)\land\lambda_n(i)=\lambda_n(k)}{%
      \rA}
    \proofstepwidestar[1,3]{%
      \Succ(i)=\Succ(k)}{%
      \FormulaRefAuto{FiniteSequenceIndexShiftDef}[def]{1,3}}
    \proofstepwidestar[1,3]{%
      i=k}{%
      \FormulaRefAuto{PeanoSuccInjective}[ax]{\FormulaRefAuto{PeanoNatSegElementNatural}[thm]{1,\rAEa{\rAEa{3}}},\FormulaRefAuto{PeanoNatSegElementNatural}[thm]{1,\rAEb{\rAEa{3}}},4}}
    \proofstepwidestar[1]{%
      \lambda_n\colon\NatSeg n\inj\PosNatSeg n}{%
      \FormulaRefAuto{Injektive Funktion}[def]{2,\rUI{\rUI{\rRI{3,5}}}}}
    \proofstepwidestar[7]{%
      j\in\PosNatSeg n}{%
      \rA}
    \proofstepwidestar[1,7]{%
      j\in\mathbb N\land(1\leq j\land j\leq n)}{%
      \FormulaRefAuto{PositiveNaturalSegmentDef}[def]{1,7}}
    \proofstepwidestar[1,7]{%
      j\neq0}{%
      \FormulaRefAuto{PeanoLtAddOneLeq}[thm]{\FormulaRefAuto{PeanoZero}[ax],\rAEa{8},\rAEa{\rAEb{8}}},\ \FormulaRefAuto{PeanoNoLtZero}[thm]}
    \proofstepwidestar[1,7]{%
      \Pred(j)\in\mathbb N}{%
      \FormulaRefAuto{PeanoPredNatural}[thm]{\rAEa{8},9}}
    \proofstepwidestar[1,7]{%
      \Succ(\Pred(j))=j}{%
      \FormulaRefAuto{PeanoPredRightInverse}[thm]{\rAEa{8},9}}
    \proofstepwidestar[1,7]{%
      \Pred(j)<n}{%
      \FormulaRefAuto{PeanoLtAddOneLeq}[thm]{10,1},\ \FormulaRefAuto{PeanoAddOneSucc}[thm]{10},\ \rIE{11,\rAEb{\rAEb{8}}}}
    \proofstepwidestar[1,7]{%
      \Pred(j)\in\NatSeg n}{%
      \FormulaRefAuto{PeanoNatSegMembership}[thm]{1,10,12}}
    \proofstepwidestar[1,7]{%
      \lambda_n(\Pred(j))=j}{%
      \FormulaRefAuto{FiniteSequenceIndexShiftDef}[def]{1,13},\ \rIE{11}}
    \proofstepwidestar[1,7]{%
      \exists i\in\NatSeg n\,(\lambda_n(i)=j)}{%
      \rEI{\rAI{13,14}}}
    \proofstepwidestar[1]{%
      \lambda_n\colon\NatSeg n\sur\PosNatSeg n}{%
      \FormulaRefAuto{Surjektive Funktion}[def]{2,\rUI{\rRI{7,15}}}}
    \proofstepwidestar[1]{%
      \lambda_n\colon\NatSeg n\bij\PosNatSeg n}{%
      \FormulaRefAuto{Bijektive Funktion}[def]{6,16}}
  \closeproofpartwideindR

  \proofpartwide{Wert der Umkehrfunktion}
    \proofstepwidestar[1]{%
      n\in\mathbb N}{%
      \rA}
    \proofstepwidestar[2]{%
      j\in\PosNatSeg n}{%
      \rA}
    \proofstepwidestar[1]{%
      \lambda_n\colon\NatSeg n\bij\PosNatSeg n}{%
      \FormulaRefAuto{B21FiniteIndexShiftBijection}[thm]{1}}
    \proofstepwidestar[1,2]{%
      \lambda_n^{-1}(j)\in\NatSeg n}{%
      \FormulaRefAuto{InverseFunctionDef}[def]{3,2}}
    \proofstepwidestar[1,2]{%
      \Succ(\lambda_n^{-1}(j))=j}{%
      \FormulaRefAuto{F\colon A\bij B\dsep y\in B\vdash F(F^{-1}(y))=y}[thm]{3,2},\ \FormulaRefAuto{FiniteSequenceIndexShiftDef}[def]{1,4}}
    \proofstepwidestar[1,2]{%
      \Pred(\Succ(\lambda_n^{-1}(j)))=\lambda_n^{-1}(j)}{%
      \FormulaRefAuto{PeanoPredSucc}[thm]{\FormulaRefAuto{PeanoNatSegElementNatural}[thm]{1,4}}}
    \proofstepwidestar[1,2]{%
      \lambda_n^{-1}(j)=\Pred(j)}{%
      \rIE{5,6},\ \FormulaRefAuto{a=b\vdash b=a}[thm]}
  \closeproofpartwide
\end{tabproofsplitwide}

\FormulaThmDeltaK[Umrechnung endlicher Folgen]{%
  \begin{aligned}[t]
  &n\in\mathbb N\dsep a\colon\NatSeg{n}\to M\vdash
    a\circ\lambda_n^{-1}\colon\PosNatSeg{n}\to M\dsep{}\\[-2pt]
  &\quad\forall j\in\PosNatSeg{n}\,
    \bigl((a\circ\lambda_n^{-1})_j=a_{\Pred(j)}\bigr)
  \end{aligned}%
}{FiniteSequenceConventionConversion}{%
  \DeltaRow{Mengen}{M\dsep n\dsep j}
  \DeltaRow{Funktionen}{a\dsep\lambda_n}
}
\begin{tabproofsplitwide}
  \proofpartwide{Umrechnung}
    \proofstepwidestar[1]{%
      n\in\mathbb N}{%
      \rA}
    \proofstepwidestar[2]{%
      a\colon\NatSeg n\to M}{%
      \rA}
    \proofstepwidestar[1]{%
      \lambda_n\colon\NatSeg n\bij\PosNatSeg n}{%
      \FormulaRefAuto{B21FiniteIndexShiftBijection}[thm]{1}}
    \proofstepwidestar[1]{%
      \lambda_n^{-1}\colon\PosNatSeg n\to\NatSeg n}{%
      \FormulaRefAuto{InverseFunctionDef}[def]{3}}
    \proofstepwidestar[1,2]{%
      a\circ\lambda_n^{-1}\colon\PosNatSeg n\to M}{%
      \FormulaRefAuto{CompositionDef}[def]{4,2}}
    \proofstepwidestar[6]{%
      j\in\PosNatSeg n}{%
      \rA}
    \proofstepwidestar[1,6]{%
      \lambda_n^{-1}(j)=\Pred(j)}{%
      \FormulaRefAuto{FiniteSequenceIndexShiftBijective}[thm]{1,6}}
    \proofstepwidestar[1,2,6]{%
      (a\circ\lambda_n^{-1})_j=a_{\Pred(j)}}{%
      \FormulaRefAuto{CompositionDef}[def]{4,2,6},\ \FormulaRefAuto{IndexedFamilyDef}[def]{5,2},\ \rIE{7}}
    \proofstepwidestar[1,2]{%
      \forall j\in\PosNatSeg n\,((a\circ\lambda_n^{-1})_j=a_{\Pred(j)})}{%
      \rUI{\rRI{6,8}}}
  \closeproofpartwide
\end{tabproofsplitwide}

Die Umrechnung ändert nur die Namen der Indizes, weder die Reihenfolge noch
die Wertemenge. Damit können nullbasierte und positiv indizierte endliche
Folgen verwendet werden, ohne eine versteckte Verschiebung der Länge
einzuführen.

\chapter{Grundlegende Beispiele und Rekursion}

\section{Konstante und kanonische Familien}

\FormulaDefDeltaK[Konstante Familie]{%
  c\in M\vdash
  \begin{aligned}[t]
    &\operatorname{const}_{I,c}\colon I\to M,\\[-2pt]
    &\forall i\in I\,
      \bigl((\operatorname{const}_{I,c})_i\coloneqq c\bigr)
  \end{aligned}%
}{ConstantIndexedFamilyDef}{%
  \DeltaRow{Mengen}{I\dsep M\dsep c\dsep i}
  \DeltaRow{Funktionen}{\operatorname{const}_{I,c}}
}

Ist \(I\neq\varnothing\), so hat die konstante Familie die Wertemenge
\(\{c\}\); bei leerer Indexmenge ist ihre Wertemenge leer. Für
\(I=\mathbb N\) entsteht die konstante Folge \((c,c,c,\ldots)\).

\FormulaDefDeltaK[Kanonische Aufzählungsfolge der natürlichen Zahlen]{%
  \begin{aligned}[t]
    &\iota_{\mathbb N}\colon\mathbb N\to\mathbb N,\\[-2pt]
    &\forall n\in\mathbb N\,
      \bigl((\iota_{\mathbb N})_n\coloneqq n\bigr)
  \end{aligned}%
}{NaturalEnumerationSequenceDef}{%
  \DeltaRow{Mengen}{n}
  \DeltaRow{Funktionen}{\iota_{\mathbb N}}
}

Diese Folge ist bijektiv und hat Bildmenge \(\mathbb N\). Die Folge
\((0,2,4,\ldots)\) der geraden natürlichen Zahlen wird entsprechend durch
die untereinander stehende Typ- und Grundgleichung eingeführt:
\[
  \begin{aligned}[t]
    &e\colon\mathbb N\to\mathbb N,\\[-2pt]
    &\forall n\in\mathbb N\,\bigl(e_n\coloneqq 2n\bigr).
  \end{aligned}
\]
Ihre Zielmenge ist \(\mathbb N\), ihre Wertemenge aber nur die Menge der
geraden Zahlen.

\section{Rekursiv definierte Folgen}

\FormulaThmDeltaK[Rekursive Definition einer Folge durch einen Schritt]{%
  x_0\in M\dsep F\colon M\to M
  \vdash
  \exists!a\colon\mathbb N\to M\,
  \bigl(a_0=x_0\land
    \forall n\in\mathbb N\,(a_{\Succ(n)}=F(a_n))\bigr)%
}{RecursiveSequenceExistenceUnique}{%
  \DeltaRow{Mengen}{M\dsep x_0\dsep n}
  \DeltaRow{Funktionen}{F\dsep a}
}
\begin{tabproofsplitwide}
  \proofpartwide{Rekursionssatz in Koordinatenschreibweise}
    \proofstepwidestar[1]{%
      x_0\in M}{%
      \rA}
    \proofstepwidestar[2]{%
      F\colon M\to M}{%
      \rA}
    \proofstepwidestar[1,2]{%
      \begin{aligned}[t]&\exists!a\colon\mathbb N\to M\,\bigl(a(0)=x_0\land{}\\[-2pt]&\qquad\forall n\in\mathbb N\,(a(\Succ(n))=F(a(n)))\bigr)\end{aligned}}{%
      \FormulaRefAuto{DedekindRecursionTheorem}[thm]{1,2}}
    \proofstepwidestar[1,2]{%
      \begin{aligned}[t]&\exists!a\colon\mathbb N\to M\,\bigl(a_0=x_0\land{}\\[-2pt]&\qquad\forall n\in\mathbb N\,(a_{\Succ(n)}=F(a_n))\bigr)\end{aligned}}{%
      \FormulaRefAuto{IndexedFamilyDef}[def]{3}}
  \closeproofpartwide
\end{tabproofsplitwide}

\begin{example}[Iteration]
Zu \(x_0\in M\) und \(F\colon M\to M\) gehört die eindeutig bestimmte
Iterationsfolge
\[
  x_0,\quad F(x_0),\quad F(F(x_0)),\quad\ldots.
\]
Man schreibt auch \(a_n=F^n(x_0)\), wobei \(F^0=\Id_M\) ist.
\end{example}

\begin{example}[Arithmetische Folge]
Für \(u,d\in\mathbb R\) wird nach Band~19 die benötigte Schrittabbildung
durch
\[
  \begin{aligned}[t]
    &F_d\colon\mathbb R\to\mathbb R,\\[-2pt]
    &\forall x\in\mathbb R\,
      \bigl(F_d(x)\coloneqq x\RealAdd d\bigr)
  \end{aligned}
\]
eingeführt. Daher wird durch
\[
  a_0=u,
  \qquad a_{\Succ(n)}=F_d(a_n)
\]
eine eindeutig bestimmte reelle Folge definiert. In der konventionellen
Notation des nächsten Kapitels lautet der Schritt \(a_{n+1}=a_n+d\).
\end{example}

\begin{example}[Geometrische Folge]
Für \(u,q\in\mathbb R\) wird entsprechend
\[
  \begin{aligned}[t]
    &G_q\colon\mathbb R\to\mathbb R,\\[-2pt]
    &\forall x\in\mathbb R\,
      \bigl(G_q(x)\coloneqq q\RealMul x\bigr)
  \end{aligned}
\]
eingeführt. Die Gleichungen
\[
  a_0=u,
  \qquad a_{\Succ(n)}=G_q(a_n)
\]
definieren daher eine eindeutig bestimmte reelle Folge, ohne dass eine
reelle Potenznotation vorweggenommen werden muss.
\end{example}

\section{Wachsende Folgen von Mengen}

\FormulaThmDeltaK[Additive Verschachtelung einer wachsenden Mengenfolge]{%
  \begin{aligned}[t]
    &D\colon\mathbb N\to\powerset(A)\dsep
      \forall j\in\mathbb N\,D(j)\subseteq D(j+1)\dsep
      n,k\in\mathbb N\vdash{}\\[-2pt]
    &D(n)\subseteq D(n+k)
  \end{aligned}%
}{SubsetSequenceAdditiveNesting}{%
  \DeltaRow{Mengen}{A\dsep j\dsep n\dsep k}
  \DeltaRow{Funktionen}{D}
}
\begin{tabproofsplitwide}
  \proofpartwideindR[Induktionsanfang]{%
    n\in\mathbb N\vdash D(n)\subseteq D(n+0)
  }{n\in\mathbb N\vdash D(n)\subseteq D(n+0)}[B21SubsetSequenceBase]
    \proofstepwidestar[1]{%
      n\in\mathbb N}{%
      \rA}
    \proofstepwidestar[1]{%
      n+0=n}{%
      \FormulaRefAuto{PeanoAddRightZero}[thm]{1}}
    \proofstepwidestar[]{%
      D(n)\subseteq D(n)}{%
      \FormulaRefAuto{A\subseteq A}[thm]}
    \proofstepwidestar[1]{%
      D(n)\subseteq D(n+0)}{%
      \rIE{2,3}}
  \closeproofpartwideindR

  \proofpartwideindR[Induktionsschritt]{%
    \begin{aligned}[t]&\forall j\in\mathbb N\,D(j)\subseteq D(j+1)\dsep n,k\in\mathbb N\dsep{}\\[-2pt]&D(n)\subseteq D(n+k)\vdash D(n)\subseteq D(n+(k+1))\end{aligned}
  }{\begin{aligned}[t]&\forall j\in\mathbb N\,D(j)\subseteq D(j+1)\dsep n,k\in\mathbb N\dsep{}\\[-2pt]&D(n)\subseteq D(n+k)\vdash D(n)\subseteq D(n+(k+1))\end{aligned}}[B21SubsetSequenceStep]
    \proofstepwidestar[1]{%
      \forall j\in\mathbb N\,D(j)\subseteq D(j+1)}{%
      \rA}
    \proofstepwidestar[2]{%
      n,k\in\mathbb N}{%
      \rA}
    \proofstepwidestar[3]{%
      D(n)\subseteq D(n+k)}{%
      \rA}
    \proofstepwidestar[2]{%
      n+k\in\mathbb N}{%
      \FormulaRefAuto{PeanoAddClosure}[thm]{2}}
    \proofstepwidestar[1,2]{%
      D(n+k)\subseteq D((n+k)+1)}{%
      \rRE{\rUE{1},4}}
    \proofstepwidestar[1,2,3]{%
      D(n)\subseteq D((n+k)+1)}{%
      \FormulaRefAuto{A\subseteq B\dsep B\subseteq C\vdash A\subseteq C}[thm]{3,5}}
    \proofstepwidestar[2]{%
      n+(k+1)=(n+k)+1}{%
      \FormulaRefAuto{PeanoAddAddOneRight}[thm]{2}}
    \proofstepwidestar[1,2,3]{%
      D(n)\subseteq D(n+(k+1))}{%
      \rIE{7,6}}
  \closeproofpartwideindR

  \proofpartwide{Induktionsschluss}
    \proofstepwidestar[1]{%
      D\colon\mathbb N\to\powerset(A)}{%
      \rA}
    \proofstepwidestar[2]{%
      \forall j\in\mathbb N\,D(j)\subseteq D(j+1)}{%
      \rA}
    \proofstepwidestar[3]{%
      n,k\in\mathbb N}{%
      \rA}
    \proofstepwidestar[1,2,3]{%
      D(n)\subseteq D(n+k)}{%
      \FormulaRefAuto{PeanoInductionRuleAddOne}[thm]{\FormulaRefAuto{B21SubsetSequenceBase}[thm],\FormulaRefAuto{B21SubsetSequenceStep}[thm],2,3}}
  \closeproofpartwide
\end{tabproofsplitwide}

\FormulaThmDeltaK[Verschachtelung einer wachsenden Mengenfolge]{%
  \begin{aligned}[t]
    &D\colon\mathbb N\to\powerset(A)\dsep
      \forall j\in\mathbb N\,D(j)\subseteq D(j+1)\dsep
      m,n\in\mathbb N\dsep m\leq n\vdash{}\\[-2pt]
    &D(m)\subseteq D(n)
  \end{aligned}%
}{SubsetSequenceNesting}{%
  \DeltaRow{Mengen}{A\dsep j\dsep m\dsep n\dsep k}
  \DeltaRow{Funktionen}{D}
}
\begin{tabproofsplitwide}
  \proofpartwide{Abstandsdarstellung}
    \proofstepwidestar[1]{%
      D\colon\mathbb N\to\powerset(A)}{%
      \rA}
    \proofstepwidestar[2]{%
      \forall j\in\mathbb N\,D(j)\subseteq D(j+1)}{%
      \rA}
    \proofstepwidestar[3]{%
      m,n\in\mathbb N}{%
      \rA}
    \proofstepwidestar[4]{%
      m\leq n}{%
      \rA}
    \proofstepwidestar[3,4]{%
      \exists k\in\mathbb N\,(m+k=n)}{%
      \FormulaRefAuto{m\in\mathbb{N}\dsep n\in\mathbb{N}\vdash m\leq n\leftrightarrow\exists k\in\mathbb{N}\,(m+k=n)}[thm]{3,4}}
    \proofstepwidestar[6]{%
      k\in\mathbb N\land m+k=n}{%
      \rA}
    \proofstepwidestar[1,2,3,6]{%
      D(m)\subseteq D(m+k)}{%
      \FormulaRefAuto{SubsetSequenceAdditiveNesting}[thm]{1,2,3,\rAEa{6}}}
    \proofstepwidestar[1,2,3,6]{%
      D(m)\subseteq D(n)}{%
      \rIE{\rAEb{6},7}}
    \proofstepwidestar[1,2,3,4]{%
      D(m)\subseteq D(n)}{%
      \rEE{5,6,8}}
  \closeproofpartwide
\end{tabproofsplitwide}

\chapter{Reelle Folgen und Grenzwerte}

\section{Reelle Notation und punktweise Operationen}

Band~19 konstruiert die reellen Operationen zunächst unter eigenen Symbolen.
In diesem und in späteren analytischen Bänden werden dafür die üblichen
Körpersymbole als reine Abkürzungen verwendet. Rationale Zahlen und
insbesondere rationale Zahlzeichen werden dabei, wie am Ende von Band~19
vereinbart, über \(\iota_{\mathbb R}\) mit ihren reellen Bildern
identifiziert.
Die folgenden Abkürzungen gelten innerhalb dieses Analysis-Kapitels, sobald
die beteiligten Terme als reell typisiert sind. Sie ändern nicht die
natürlichen Tags \(0,1,\ldots\) in späteren rein mengentheoretischen Kapiteln.

\FormulaDefDeltaK[Konventionelle Schreibweise für reelle Operationen]{%
  \begin{aligned}[t]
  &0\coloneqq\RealZero\dsep 1\coloneqq\RealOne,\\[-2pt]
  &x,y\in\mathbb R\vdash
    \begin{gathered}[t]
      x+y\coloneqq x\RealAdd y,\qquad
      -x\coloneqq\RealNeg{x},\qquad
      x-y\coloneqq x\RealAdd\RealNeg{y},\\[-2pt]
      xy\coloneqq x\RealMul y,\qquad
      (x\leq y)\coloneqq(x\RealLe y),\qquad
      (x<y)\coloneqq(x\RealLt y),\\[-2pt]
      (x\geq y)\coloneqq(y\RealLe x),\qquad
      (x>y)\coloneqq(y\RealLt x),\qquad
      |x|\coloneqq\RealAbs{x},
    \end{gathered}\\[-2pt]
  &x,y\in\mathbb R\dsep y\neq0\vdash
      y^{-1}\coloneqq\RealInv{y}\dsep
      \frac{x}{y}\coloneqq x\RealMul\RealInv{y}.
  \end{aligned}%
}{RealSequenceConventionalNotation}{%
  \DeltaRow{Mengen}{x\dsep y\dsep\mathbb R}
  \DeltaRow{Relationsoperatoren}{\RealLe\dsep\RealLt}
  \DeltaRow{Funktionsoperatoren}{\RealAdd\dsep\RealMul
    \dsep\RealNeg{\,\cdot\,}\dsep\RealInv{\,\cdot\,}
    \dsep\RealAbs{\,\cdot\,}}
}

Alle folgenden Rechnungen mit gewöhnlichen Symbolen sind daher Aussagen über
die in Band~19 registrierten Operationen; es wird keine zweite reelle
Arithmetik eingeführt.

\FormulaDefDeltaK[Punktweise Operationen reeller Folgen]{%
  \begin{aligned}[t]
  &a,b\colon\mathbb N\to\mathbb R\vdash\\[-2pt]
  &\quad a+b,\ a-b,\ a\mathbin{\cdot}b\colon\mathbb N\to\mathbb R,\\[-2pt]
  &\quad \forall n\in\mathbb N\,
    \begin{gathered}[t]
      (a+b)_n\coloneqq a_n+b_n,\qquad
      (a-b)_n\coloneqq a_n-b_n,\\[-2pt]
      (a\mathbin{\cdot}b)_n\coloneqq a_nb_n
    \end{gathered};\\[2pt]
  &a\colon\mathbb N\to\mathbb R\vdash\\[-2pt]
  &\quad -a,\ |a|\colon\mathbb N\to\mathbb R,\\[-2pt]
  &\quad \forall n\in\mathbb N\,
      \bigl((-a)_n\coloneqq-a_n,\qquad(|a|)_n\coloneqq|a_n|\bigr);\\[2pt]
  &a\colon\mathbb N\to\mathbb R\dsep c\in\mathbb R\vdash\\[-2pt]
  &\quad ca\colon\mathbb N\to\mathbb R,\\[-2pt]
  &\quad
    \forall n\in\mathbb N\,((ca)_n\coloneqq c\,a_n),\\[-2pt]
  &a\colon\mathbb N\to\mathbb R\dsep L\in\mathbb R\vdash\\[-2pt]
  &\quad a-L\colon\mathbb N\to\mathbb R,\\[-2pt]
  &\quad
    \forall n\in\mathbb N\,((a-L)_n\coloneqq a_n-L).
  \end{aligned}%
}{PointwiseRealSequenceOperationsDef}{%
  \DeltaRow{Mengen}{c\dsep L\dsep n}
  \DeltaRow{Folgen}{a\dsep b\dsep a+b\dsep a-b}%
  \DeltaRow{Folgen}{-a\dsep ca\dsep a\mathbin{\cdot}b}%
  \DeltaRow{Folgen}{|a|\dsep a-L}
}
\begin{tabproofsplitwide}
  \proofpartwide{Typisierung der rechten Terme}
    \proofstepwidestar[1]{%
      a,b\colon\mathbb N\to\mathbb R}{%
      \rA}
    \proofstepwidestar[2]{%
      c,L\in\mathbb R}{%
      \rA}
    \proofstepwidestar[3]{%
      n\in\mathbb N}{%
      \rA}
    \proofstepwidestar[1,3]{%
      a_n,b_n\in\mathbb R}{%
      \FormulaRefAuto{IndexedFamilyCoordinateTyping}[thm]{1,3}}
    \proofstepwidestar[1,3]{%
      -a_n,-b_n\in\mathbb R}{%
      \FormulaRefAuto{RealCutAdditiveClosure}[thm]{4},\ \FormulaRefAuto{RealSequenceConventionalNotation}[def]}
    \proofstepwidestar[1,3]{%
      a_n+b_n,\ a_n-b_n\in\mathbb R}{%
      \FormulaRefAuto{RealCutAdditiveClosure}[thm]{4,5},\ \FormulaRefAuto{RealSequenceConventionalNotation}[def]}
    \proofstepwidestar[1,3]{%
      a_nb_n\in\mathbb R}{%
      \FormulaRefAuto{RealCutMultiplicativeClosure}[thm]{4},\ \FormulaRefAuto{RealSequenceConventionalNotation}[def]}
    \proofstepwidestar[1,3]{%
      |a_n|\in\mathbb R}{%
      \FormulaRefAuto{RealAbsoluteValue}[def]{4},\ \FormulaRefAuto{RealCutAdditiveClosure}[thm]{4}}
    \proofstepwidestar[1,2,3]{%
      c\,a_n\in\mathbb R}{%
      \FormulaRefAuto{RealCutMultiplicativeClosure}[thm]{2,4}}
    \proofstepwidestar[1,2,3]{%
      a_n-L\in\mathbb R}{%
      \FormulaRefAuto{RealCutAdditiveClosure}[thm]{2,4}}
    \proofstepwidestar[1,2]{%
      \begin{aligned}[t]&a+b,\ a-b,\ a\mathbin{\cdot}b,\ -a,\ |a|,\ ca,\ a-L\\[-2pt]&\qquad\colon\mathbb N\to\mathbb R\end{aligned}}{%
      \FormulaRefAuto{TypedFunctionDefinitionPrinciple}[thm]{5,6,7,8,9,10}}
  \closeproofpartwide
\end{tabproofsplitwide}

\FormulaDefDeltaK[Punktweiser Kehrwert und Quotient reeller Folgen]{%
  \begin{aligned}[t]
  &b\colon\mathbb N\to\mathbb R\dsep
    \forall n\in\mathbb N\,(b_n\neq0)\vdash\\[-2pt]
  &\quad b^{-1}\colon\mathbb N\to\mathbb R,\\[-2pt]
  &\quad
    \forall n\in\mathbb N\,
      ((b^{-1})_n\coloneqq\tfrac1{b_n});\\[2pt]
  &a,b\colon\mathbb N\to\mathbb R\dsep
    \forall n\in\mathbb N\,(b_n\neq0)\vdash\\[-2pt]
  &\quad\frac ab\colon\mathbb N\to\mathbb R,\\[-2pt]
  &\quad
    \forall n\in\mathbb N\,
      \left(\left(\frac ab\right)_n\coloneqq\frac{a_n}{b_n}\right).
  \end{aligned}%
}{PointwiseRealSequenceReciprocalQuotientDef}{%
  \DeltaRow{Mengen}{n}
  \DeltaRow{Funktionen}{a\dsep b\dsep b^{-1}\dsep a/b}
}
\begin{tabproofsplitwide}
  \proofpartwide{Nichtnullbedingung und Typisierung}
    \proofstepwidestar[1]{%
      a,b\colon\mathbb N\to\mathbb R}{%
      \rA}
    \proofstepwidestar[2]{%
      \forall n\in\mathbb N\,(b_n\neq0)}{%
      \rA}
    \proofstepwidestar[3]{%
      n\in\mathbb N}{%
      \rA}
    \proofstepwidestar[1,3]{%
      a_n,b_n\in\mathbb R}{%
      \FormulaRefAuto{IndexedFamilyCoordinateTyping}[thm]{1,3}}
    \proofstepwidestar[2,3]{%
      b_n\neq0}{%
      \rRE{\rUE{2},3}}
    \proofstepwidestar[1,2,3]{%
      \frac1{b_n}\in\mathbb R}{%
      \FormulaRefAuto{RealCutMultiplicativeClosure}[thm]{4,5},\ \FormulaRefAuto{RealSequenceConventionalNotation}[def]}
    \proofstepwidestar[1,2,3]{%
      \frac{a_n}{b_n}\in\mathbb R}{%
      \FormulaRefAuto{RealCutMultiplicativeClosure}[thm]{4,6}}
    \proofstepwidestar[1,2]{%
      b^{-1},\ a/b\colon\mathbb N\to\mathbb R}{%
      \FormulaRefAuto{TypedFunctionDefinitionPrinciple}[thm]{6,7}}
  \closeproofpartwide
\end{tabproofsplitwide}

\section{Konvergenz}

Von nun an seien Folgen ohne anderslautenden Zusatz reelle Folgen, also
Funktionen \(a\colon\mathbb N\to\mathbb R\). Der Ausdruck
\(\varepsilon>0\) bedeutet stets
\(\varepsilon\in\mathbb R\land\RealZero\RealLt\varepsilon\).

\FormulaDefDeltaK[Konvergenz einer reellen Folge]{%
  \begin{aligned}[t]
  &a\colon\mathbb N\to\mathbb R\dsep L\in\mathbb R\vdash
    \bigl(a_n\longrightarrow L\bigr)\coloneqq{}\\[-2pt]
  &\quad\forall\varepsilon\in\mathbb R\,
    \Bigl(\varepsilon>0\rightarrow
      \exists N\in\mathbb N\,
      \forall n\in\mathbb N\,
        (N\leq n\rightarrow|a_n-L|<\varepsilon)\Bigr)
  \end{aligned}%
}{RealSequenceConvergenceDef}{%
  \DeltaRow{Mengen}{L\dsep\varepsilon\dsep N\dsep n}
  \DeltaRow{Funktionen}{a}
  \DeltaRow{Neue Schreibweise}{a_n\to L}
}

Gilt \(a_n\to L\), so heißt \(L\) ein \textbf{Grenzwert} der Folge. Man
schreibt auch
\[
  \lim_{n\to\infty}a_n=L.
\]
Die Zahl \(N\) darf von \(\varepsilon\) abhängen, aber nicht vom später
gewählten \(n\). Die Forderung betrifft nur Indizes \(n\geq N\); endlich
viele Anfangsglieder haben deshalb keinen Einfluss auf Konvergenz und
Grenzwert.

\FormulaDefDeltaK[Konvergente und divergente Folge]{%
  a\colon\mathbb N\to\mathbb R\vdash
  \operatorname{Konv}(a)\coloneqq
    \exists L\in\mathbb R\,(a_n\to L)\dsep
  \operatorname{Div}(a)\coloneqq\neg\operatorname{Konv}(a)%
}{RealSequenceConvergentDivergentDef}{%
  \DeltaRow{Mengen}{L}
  \DeltaRow{Funktionen}{a}
  \DeltaRow{Neue Prädikate}{\operatorname{Konv}\dsep\operatorname{Div}}
}

\FormulaThmDeltaK[Eindeutigkeit des Grenzwerts]{%
  a\colon\mathbb N\to\mathbb R\dsep
  L,K\in\mathbb R\dsep
  a_n\to L\dsep a_n\to K\vdash L=K%
}{RealSequenceLimitUnique}{%
  \DeltaRow{Mengen}{L\dsep K\dsep\varepsilon\dsep N\dsep n}
  \DeltaRow{Funktionen}{a}
}
\begin{tabproofsplitwide}
  \proofpartwideindR[Gemeinsame natürliche Indexschranke]{%
    N,M\in\mathbb N\vdash N+M\in\mathbb N\land(N\leq N+M\land M\leq N+M)
  }{N,M\in\mathbb N\vdash N+M\in\mathbb N\land(N\leq N+M\land M\leq N+M)}[B21CommonNaturalBound]
    \proofstepwidestar[1]{%
      N,M\in\mathbb N}{%
      \rA}
    \proofstepwidestar[1]{%
      N+M\in\mathbb N}{%
      \FormulaRefAuto{PeanoAddClosure}[thm]{1}}
    \proofstepwidestar[1]{%
      N\subseteq N+M}{%
      \FormulaRefAuto{PeanoAddLeftSubset}[thm]{1}}
    \proofstepwidestar[1]{%
      N\leq N+M}{%
      \FormulaRefAuto{NaturalOrderOperator}[def]{1,2,3},\ \FormulaRefAuto{NaturalOrderNotation}[def]}
    \proofstepwidestar[1]{%
      M\subseteq M+N}{%
      \FormulaRefAuto{PeanoAddLeftSubset}[thm]{1}}
    \proofstepwidestar[1]{%
      M+N=N+M}{%
      \FormulaRefAuto{PeanoAddCommutative}[thm]{1}}
    \proofstepwidestar[1]{%
      M\leq N+M}{%
      \FormulaRefAuto{NaturalOrderOperator}[def]{1,2,5},\ \FormulaRefAuto{NaturalOrderNotation}[def],\ \rIE{6}}
    \proofstepwidestar[1]{%
      N+M\in\mathbb N\land(N\leq N+M\land M\leq N+M)}{%
      \rAI{2,\rAI{4,7}}}
  \closeproofpartwideindR

  \proofpartwide{Widerspruch bei verschiedenen Grenzwerten}
    \proofstepwidestar[1]{%
      a\colon\mathbb N\to\mathbb R}{%
      \rA}
    \proofstepwidestar[2]{%
      L,K\in\mathbb R}{%
      \rA}
    \proofstepwidestar[3]{%
      a_n\to L\land a_n\to K}{%
      \rA}
    \proofstepwidestar[4]{%
      L\neq K}{%
      \rA}
    \proofstepwidestar[2,4]{%
      0<|L-K|}{%
      \FormulaRefAuto{RealAbsoluteValueNonnegativeZero}[thm]{2,4},\ \FormulaRefAuto{RealCutTotalOrder}[thm]}
    \proofstepwidestar[2,4]{%
      0<\frac{|L-K|}{3}}{%
      \FormulaRefAuto{RealCutsCompleteOrderedField}[thm]{2,5}}
    \proofstepwidestar[1,2,3,4]{%
      \exists N\in\mathbb N\,\forall k\in\mathbb N\,(N\leq k\rightarrow|a_k-L|<|L-K|/3)}{%
      \FormulaRefAuto{RealSequenceConvergenceDef}[def]{1,2,\rAEa{3},6}}
    \proofstepwidestar[1,2,3,4]{%
      \exists M\in\mathbb N\,\forall k\in\mathbb N\,(M\leq k\rightarrow|a_k-K|<|L-K|/3)}{%
      \FormulaRefAuto{RealSequenceConvergenceDef}[def]{1,2,\rAEb{3},6}}
    \proofstepwidestar[9]{%
      \begin{aligned}[t]&N\in\mathbb N\land{}\\[-2pt]&\forall k\in\mathbb N\,(N\leq k\rightarrow|a_k-L|<|L-K|/3)\end{aligned}}{%
      \rA}
    \proofstepwidestar[10]{%
      \begin{aligned}[t]&M\in\mathbb N\land{}\\[-2pt]&\forall k\in\mathbb N\,(M\leq k\rightarrow|a_k-K|<|L-K|/3)\end{aligned}}{%
      \rA}
    \proofstepwidestar[9,10]{%
      N+M\in\mathbb N\land(N\leq N+M\land M\leq N+M)}{%
      \FormulaRefAuto{B21CommonNaturalBound}[thm]{\rAEa{9},\rAEa{10}}}
    \proofstepwidestar[9,10]{%
      |a_{N+M}-L|<|L-K|/3}{%
      \rRE{\rRE{\rUE{\rAEb{9}},\rAEa{11}},\rAEa{\rAEb{11}}}}
    \proofstepwidestar[9,10]{%
      |a_{N+M}-K|<|L-K|/3}{%
      \rRE{\rRE{\rUE{\rAEb{10}},\rAEa{11}},\rAEb{\rAEb{11}}}}
    \proofstepwidestar[1,2,9,10]{%
      |L-K|\leq|a_{N+M}-L|+|a_{N+M}-K|}{%
      \FormulaRefAuto{RealDistanceTriangle}[thm]{2,\FormulaRefAuto{IndexedFamilyCoordinateTyping}[thm]{1,\rAEa{11}}},\ \FormulaRefAuto{RealDistanceBasicLaws}[thm]}
    \proofstepwidestar[1,2,9,10]{%
      |L-K|<2|L-K|/3}{%
      \FormulaRefAuto{RealCutOrderArithmeticCompatibility}[thm]{14,12,13},\ \FormulaRefAuto{RealCutTotalOrder}[thm]}
    \proofstepwidestar[1,2,4,9,10]{%
      \bot}{%
      \FormulaRefAuto{RealCutsCompleteOrderedField}[thm]{5,15}}
    \proofstepwidestar[1,2,3,4,9]{%
      \bot}{%
      \rEE{8,10,16}}
    \proofstepwidestar[1,2,3,4]{%
      \bot}{%
      \rEE{7,9,17}}
    \proofstepwidestar[1,2,3]{%
      L=K}{%
      \rCE{4,18}}
  \closeproofpartwide
\end{tabproofsplitwide}

Wegen der Eindeutigkeit ist \(\lim_{n\to\infty}a_n\) ein eindeutig
bestimmter Ausdruck, sobald die Folge konvergiert.

\FormulaThmDeltaK[Konvergenz konstanter Folgen]{%
  c\in\mathbb R\vdash
  (\operatorname{const}_{\mathbb N,c})_n\to c%
}{ConstantRealSequenceConverges}{%
  \DeltaRow{Mengen}{c\dsep n}
  \DeltaRow{Funktionen}{\operatorname{const}_{\mathbb N,c}}
}
\begin{tabproofsplitwide}
  \proofpartwide{Konstante Folge}
    \proofstepwidestar[1]{%
      c\in\mathbb R}{%
      \rA}
    \proofstepwidestar[2]{%
      \varepsilon\in\mathbb R\land\varepsilon>0}{%
      \rA}
    \proofstepwidestar[1]{%
      \operatorname{const}_{\mathbb N,c}\colon\mathbb N\to\mathbb R}{%
      \FormulaRefAuto{ConstantIndexedFamilyDef}[def]{1}}
    \proofstepwidestar[4]{%
      n\in\mathbb N\land0\leq n}{%
      \rA}
    \proofstepwidestar[1,4]{%
      (\operatorname{const}_{\mathbb N,c})_n=c}{%
      \FormulaRefAuto{ConstantIndexedFamilyDef}[def]{1,\rAEa{4}}}
    \proofstepwidestar[1,2,4]{%
      |(\operatorname{const}_{\mathbb N,c})_n-c|=0<\varepsilon}{%
      \FormulaRefAuto{RealCutAdditiveLaws}[thm]{1},\ \FormulaRefAuto{RealAbsoluteValueNonnegativeZero}[thm],\ \rIE{5},\ 2}
    \proofstepwidestar[1,2]{%
      \exists N\in\mathbb N\,\forall n\in\mathbb N\,(N\leq n\rightarrow|(\operatorname{const}_{\mathbb N,c})_n-c|<\varepsilon)}{%
      \rEI{\rAI{\FormulaRefAuto{PeanoZero}[ax],\rUI{\rRI{4,6}}}}}
    \proofstepwidestar[1]{%
      (\operatorname{const}_{\mathbb N,c})_n\to c}{%
      \FormulaRefAuto{RealSequenceConvergenceDef}[def]{3,1,\rUI{\rRI{2,7}}}}
  \closeproofpartwide
\end{tabproofsplitwide}

\FormulaThmDeltaK[Endliche Änderungen verändern den Grenzwert nicht]{%
  \begin{aligned}[t]
  &a,b\colon\mathbb N\to\mathbb R\dsep
    \exists N_0\in\mathbb N\,
      \forall n\in\mathbb N\,(N_0\leq n\rightarrow a_n=b_n)\\[-2pt]
  &\quad\vdash
    \forall L\in\mathbb R\,(a_n\to L\leftrightarrow b_n\to L)
  \end{aligned}%
}{FiniteModificationPreservesLimit}{%
  \DeltaRow{Mengen}{L\dsep N_0\dsep n}
  \DeltaRow{Funktionen}{a\dsep b}
}
\begin{tabproofsplitwide}
  \proofpartwideindR[Übertragung bei schließlich gleichen Gliedern]{%
    \begin{aligned}[t]&a,b\colon\mathbb N\to\mathbb R\dsep L\in\mathbb R\dsep a_n\to L\dsep{}\\[-2pt]&\exists N_0\in\mathbb N\,\forall n\in\mathbb N\,(N_0\leq n\rightarrow a_n=b_n)\vdash b_n\to L\end{aligned}
  }{\begin{aligned}[t]&a,b\colon\mathbb N\to\mathbb R\dsep L\in\mathbb R\dsep a_n\to L\dsep{}\\[-2pt]&\exists N_0\in\mathbb N\,\forall n\in\mathbb N\,(N_0\leq n\rightarrow a_n=b_n)\vdash b_n\to L\end{aligned}}[B21EventuallyEqualLimitTransfer]
    \proofstepwidestar[1]{%
      a,b\colon\mathbb N\to\mathbb R}{%
      \rA}
    \proofstepwidestar[2]{%
      L\in\mathbb R}{%
      \rA}
    \proofstepwidestar[3]{%
      \exists N_0\in\mathbb N\,\forall k\in\mathbb N\,(N_0\leq k\rightarrow a_k=b_k)}{%
      \rA}
    \proofstepwidestar[4]{%
      a_n\to L}{%
      \rA}
    \proofstepwidestar[5]{%
      \varepsilon\in\mathbb R\land\varepsilon>0}{%
      \rA}
    \proofstepwidestar[1,2,4,5]{%
      \exists N\in\mathbb N\,\forall k\in\mathbb N\,(N\leq k\rightarrow|a_k-L|<\varepsilon)}{%
      \FormulaRefAuto{RealSequenceConvergenceDef}[def]{1,2,4,5}}
    \proofstepwidestar[7]{%
      \begin{aligned}[t]&N_0\in\mathbb N\land{}\\[-2pt]&\forall k\in\mathbb N\,(N_0\leq k\rightarrow a_k=b_k)\end{aligned}}{%
      \rA}
    \proofstepwidestar[8]{%
      \begin{aligned}[t]&N\in\mathbb N\land{}\\[-2pt]&\forall k\in\mathbb N\,(N\leq k\rightarrow|a_k-L|<\varepsilon)\end{aligned}}{%
      \rA}
    \proofstepwidestar[7,8]{%
      N_0+N\in\mathbb N\land(N_0\leq N_0+N\land N\leq N_0+N)}{%
      \FormulaRefAuto{B21CommonNaturalBound}[thm]{\rAEa{7},\rAEa{8}}}
    \proofstepwidestar[10]{%
      n\in\mathbb N\land N_0+N\leq n}{%
      \rA}
    \proofstepwidestar[7,8,10]{%
      N_0\leq n\land N\leq n}{%
      \FormulaRefAuto{PeanoLeqTransitive}[thm]{9,10}}
    \proofstepwidestar[7,8,10]{%
      a_n=b_n}{%
      \rRE{\rRE{\rUE{\rAEb{7}},\rAEa{10}},\rAEa{11}}}
    \proofstepwidestar[7,8,10]{%
      |a_n-L|<\varepsilon}{%
      \rRE{\rRE{\rUE{\rAEb{8}},\rAEa{10}},\rAEb{11}}}
    \proofstepwidestar[7,8,10]{%
      |b_n-L|<\varepsilon}{%
      \rIE{12,13}}
    \proofstepwidestar[7,8]{%
      \exists M\in\mathbb N\,\forall n\in\mathbb N\,(M\leq n\rightarrow|b_n-L|<\varepsilon)}{%
      \rEI{\rAI{\rAEa{9},\rUI{\rRI{10,14}}}}}
    \proofstepwidestar[1,2,4,5,7]{%
      \exists M\in\mathbb N\,\forall n\in\mathbb N\,(M\leq n\rightarrow|b_n-L|<\varepsilon)}{%
      \rEE{6,8,15}}
    \proofstepwidestar[1,2,3,4,5]{%
      \exists M\in\mathbb N\,\forall n\in\mathbb N\,(M\leq n\rightarrow|b_n-L|<\varepsilon)}{%
      \rEE{3,7,16}}
    \proofstepwidestar[1,2,3,4]{%
      b_n\to L}{%
      \FormulaRefAuto{RealSequenceConvergenceDef}[def]{1,2,\rUI{\rRI{5,17}}}}
  \closeproofpartwideindR

  \proofpartwide{Äquivalenz}
    \proofstepwidestar[1]{%
      a,b\colon\mathbb N\to\mathbb R}{%
      \rA}
    \proofstepwidestar[2]{%
      \exists N_0\in\mathbb N\,\forall n\in\mathbb N\,(N_0\leq n\rightarrow a_n=b_n)}{%
      \rA}
    \proofstepwidestar[3]{%
      L\in\mathbb R}{%
      \rA}
    \proofstepwidestar[2]{%
      \exists N_0\in\mathbb N\,\forall n\in\mathbb N\,(N_0\leq n\rightarrow b_n=a_n)}{%
      \FormulaRefAuto{a=b\vdash b=a}[thm]{2}}
    \proofstepwidestar[1,2,3]{%
      a_n\to L\rightarrow b_n\to L}{%
      \FormulaRefAuto{B21EventuallyEqualLimitTransfer}[thm]{1,3,2}}
    \proofstepwidestar[1,2,3]{%
      b_n\to L\rightarrow a_n\to L}{%
      \FormulaRefAuto{B21EventuallyEqualLimitTransfer}[thm]{1,3,4}}
    \proofstepwidestar[1,2,3]{%
      a_n\to L\leftrightarrow b_n\to L}{%
      \rLRI{5,6}}
    \proofstepwidestar[1,2]{%
      \forall L\in\mathbb R\,(a_n\to L\leftrightarrow b_n\to L)}{%
      \rUI{\rRI{3,7}}}
  \closeproofpartwide
\end{tabproofsplitwide}

\FormulaThmDeltaK[Restfolgen haben denselben Grenzwert]{%
  a\colon\mathbb N\to\mathbb R\dsep r\in\mathbb N\dsep L\in\mathbb R
  \vdash
  a_n\to L\leftrightarrow a^{\langle r\rangle}_n\to L%
}{SequenceTailSameLimit}{%
  \DeltaRow{Mengen}{L\dsep r\dsep n}
  \DeltaRow{Funktionen}{a\dsep a^{\langle r\rangle}}
}
\begin{tabproofsplitwide}
  \proofpartwideindR[Hinrichtung]{%
    a\colon\mathbb N\to\mathbb R\dsep r\in\mathbb N\dsep L\in\mathbb R\dsep a_n\to L\vdash a^{\langle r\rangle}_n\to L
  }{a\colon\mathbb N\to\mathbb R\dsep r\in\mathbb N\dsep L\in\mathbb R\dsep a_n\to L\vdash a^{\langle r\rangle}_n\to L}[B21TailLimitForward]
    \proofstepwidestar[1]{%
      a\colon\mathbb N\to\mathbb R}{%
      \rA}
    \proofstepwidestar[2]{%
      r\in\mathbb N\land L\in\mathbb R}{%
      \rA}
    \proofstepwidestar[3]{%
      a_n\to L}{%
      \rA}
    \proofstepwidestar[4]{%
      \varepsilon\in\mathbb R\land\varepsilon>0}{%
      \rA}
    \proofstepwidestar[1,2,3,4]{%
      \exists N\in\mathbb N\,\forall k\in\mathbb N\,(N\leq k\rightarrow|a_k-L|<\varepsilon)}{%
      \FormulaRefAuto{RealSequenceConvergenceDef}[def]{1,2,3,4}}
    \proofstepwidestar[6]{%
      \begin{aligned}[t]&N\in\mathbb N\land{}\\[-2pt]&\forall k\in\mathbb N\,(N\leq k\rightarrow|a_k-L|<\varepsilon)\end{aligned}}{%
      \rA}
    \proofstepwidestar[7]{%
      n\in\mathbb N\land N\leq n}{%
      \rA}
    \proofstepwidestar[2,7]{%
      n+r\in\mathbb N\land n\leq n+r}{%
      \FormulaRefAuto{B21CommonNaturalBound}[thm]{\rAEa{7},\rAEa{2}}}
    \proofstepwidestar[2,6,7]{%
      N\leq n+r}{%
      \FormulaRefAuto{PeanoLeqTransitive}[thm]{\rAEa{6},7,8}}
    \proofstepwidestar[1,2,6,7]{%
      |a^{\langle r\rangle}_n-L|<\varepsilon}{%
      \rRE{\rRE{\rUE{\rAEb{6}},\rAEa{8}},9},\ \FormulaRefAuto{SequenceTailDef}[def]{1,2,7}}
    \proofstepwidestar[1,2,6]{%
      \exists N\in\mathbb N\,\forall n\in\mathbb N\,(N\leq n\rightarrow|a^{\langle r\rangle}_n-L|<\varepsilon)}{%
      \rEI{\rAI{\rAEa{6},\rUI{\rRI{7,10}}}}}
    \proofstepwidestar[1,2,3,4]{%
      \exists N\in\mathbb N\,\forall n\in\mathbb N\,(N\leq n\rightarrow|a^{\langle r\rangle}_n-L|<\varepsilon)}{%
      \rEE{5,6,11}}
    \proofstepwidestar[1,2,3]{%
      a^{\langle r\rangle}_n\to L}{%
      \FormulaRefAuto{RealSequenceConvergenceDef}[def]{\FormulaRefAuto{SequenceTailDef}[def]{1,2},2,\rUI{\rRI{4,12}}}}
  \closeproofpartwideindR

  \proofpartwideindR[Rückrichtung]{%
    a\colon\mathbb N\to\mathbb R\dsep r\in\mathbb N\dsep L\in\mathbb R\dsep a^{\langle r\rangle}_n\to L\vdash a_n\to L
  }{a\colon\mathbb N\to\mathbb R\dsep r\in\mathbb N\dsep L\in\mathbb R\dsep a^{\langle r\rangle}_n\to L\vdash a_n\to L}[B21TailLimitBackward]
    \proofstepwidestar[1]{%
      a\colon\mathbb N\to\mathbb R}{%
      \rA}
    \proofstepwidestar[2]{%
      r\in\mathbb N\land L\in\mathbb R}{%
      \rA}
    \proofstepwidestar[3]{%
      a^{\langle r\rangle}_n\to L}{%
      \rA}
    \proofstepwidestar[4]{%
      \varepsilon\in\mathbb R\land\varepsilon>0}{%
      \rA}
    \proofstepwidestar[1,2,3,4]{%
      \exists N\in\mathbb N\,\forall k\in\mathbb N\,(N\leq k\rightarrow|a^{\langle r\rangle}_k-L|<\varepsilon)}{%
      \FormulaRefAuto{RealSequenceConvergenceDef}[def]{1,2,3,4}}
    \proofstepwidestar[6]{%
      \begin{aligned}[t]&N\in\mathbb N\land{}\\[-2pt]&\forall k\in\mathbb N\,(N\leq k\rightarrow|a^{\langle r\rangle}_k-L|<\varepsilon)\end{aligned}}{%
      \rA}
    \proofstepwidestar[7]{%
      n\in\mathbb N\land r+N\leq n}{%
      \rA}
    \proofstepwidestar[2,6,7]{%
      \exists t\in\mathbb N\,(r+N+t=n)}{%
      \FormulaRefAuto{m\in\mathbb N\dsep n\in\mathbb N\vdash m\leq n\leftrightarrow\exists k\in\mathbb N\,(m+k=n)}[thm]{\FormulaRefAuto{PeanoAddClosure}[thm]{2,6},7}}
    \proofstepwidestar[9]{%
      t\in\mathbb N\land r+N+t=n}{%
      \rA}
    \proofstepwidestar[2,6,9]{%
      N+t\in\mathbb N\land N\leq N+t}{%
      \FormulaRefAuto{B21CommonNaturalBound}[thm]{\rAEa{6},\rAEa{9}}}
    \proofstepwidestar[2,6,9]{%
      (N+t)+r=n}{%
      \FormulaRefAuto{PeanoAddAssociative}[thm]{2,6,9},\ \FormulaRefAuto{PeanoAddCommutative}[thm],\ \rAEb{9}}
    \proofstepwidestar[1,2,6,9]{%
      |a_n-L|<\varepsilon}{%
      \rRE{\rRE{\rUE{\rAEb{6}},\rAEa{10}},\rAEb{10}},\ \FormulaRefAuto{SequenceTailDef}[def],\ \rIE{11}}
    \proofstepwidestar[1,2,6,7]{%
      |a_n-L|<\varepsilon}{%
      \rEE{8,9,12}}
    \proofstepwidestar[1,2,6]{%
      \exists M\in\mathbb N\,\forall n\in\mathbb N\,(M\leq n\rightarrow|a_n-L|<\varepsilon)}{%
      \rEI{\rAI{\FormulaRefAuto{PeanoAddClosure}[thm]{2,6},\rUI{\rRI{7,13}}}}}
    \proofstepwidestar[1,2,3,4]{%
      \exists M\in\mathbb N\,\forall n\in\mathbb N\,(M\leq n\rightarrow|a_n-L|<\varepsilon)}{%
      \rEE{5,6,14}}
    \proofstepwidestar[1,2,3]{%
      a_n\to L}{%
      \FormulaRefAuto{RealSequenceConvergenceDef}[def]{1,2,\rUI{\rRI{4,15}}}}
  \closeproofpartwideindR

  \proofpartwide{Zusammenfassung}
    \proofstepwidestar[1]{%
      a\colon\mathbb N\to\mathbb R}{%
      \rA}
    \proofstepwidestar[2]{%
      r\in\mathbb N\land L\in\mathbb R}{%
      \rA}
    \proofstepwidestar[1,2]{%
      a_n\to L\leftrightarrow a^{\langle r\rangle}_n\to L}{%
      \rLRI{\FormulaRefAuto{B21TailLimitForward}[thm]{1,2},\FormulaRefAuto{B21TailLimitBackward}[thm]{1,2}}}
  \closeproofpartwide
\end{tabproofsplitwide}

\FormulaThmDeltaK[Teilfolgen einer konvergenten Folge]{%
  \begin{aligned}[t]
  &a\colon\mathbb N\to\mathbb R\dsep L\in\mathbb R\dsep a_n\to L\dsep
    \varphi\colon\mathbb N\to\mathbb N\dsep{}\\[-2pt]
  &\quad\forall n\in\mathbb N\,
    \bigl(\varphi(n)<\varphi(\Succ(n))\bigr)
    \vdash a_{\varphi(n)}\to L
  \end{aligned}%
}{SubsequenceOfConvergentSequence}{%
  \DeltaRow{Mengen}{L\dsep n}
  \DeltaRow{Funktionen}{a\dsep\varphi}
}
\begin{tabproofsplitwide}
  \proofpartwideindR[Eine streng wachsende Indexfolge dominiert den Index]{%
    \varphi\colon\mathbb N\to\mathbb N\dsep\forall n\in\mathbb N\,(\varphi(n)<\varphi(\Succ(n)))\vdash\forall n\in\mathbb N\,(n\leq\varphi(n))
  }{\varphi\colon\mathbb N\to\mathbb N\dsep\forall n\in\mathbb N\,(\varphi(n)<\varphi(\Succ(n)))\vdash\forall n\in\mathbb N\,(n\leq\varphi(n))}[B21IncreasingIndexDominates]
    \proofstepwidestar[1]{%
      \varphi\colon\mathbb N\to\mathbb N}{%
      \rA}
    \proofstepwidestar[2]{%
      \forall n\in\mathbb N\,(\varphi(n)<\varphi(\Succ(n)))}{%
      \rA}
    \proofstepwidestar[1]{%
      0\leq\varphi(0)}{%
      \FormulaRefAuto{PeanoZeroLeq}[thm]{\FormulaRefAuto{F\colon A\to B\dsep x\in A\vdash F(x)\in B}[thm]{1,\FormulaRefAuto{PeanoZero}[ax]}}}
    \proofstepwidestar[4]{%
      n\in\mathbb N\land n\leq\varphi(n)}{%
      \rA}
    \proofstepwidestar[1,4]{%
      \varphi(n),\varphi(\Succ(n))\in\mathbb N}{%
      \FormulaRefAuto{F\colon A\to B\dsep x\in A\vdash F(x)\in B}[thm]{1,\rAEa{4},\FormulaRefAuto{PeanoSuccClosure}[ax]}}
    \proofstepwidestar[2,4]{%
      \varphi(n)<\varphi(\Succ(n))}{%
      \rRE{\rUE{2},\rAEa{4}}}
    \proofstepwidestar[1,4]{%
      n<\Succ(\varphi(n))}{%
      \FormulaRefAuto{PeanoLeqSuccStrict}[thm]{4,5}}
    \proofstepwidestar[1,2,4]{%
      \Succ(\varphi(n))\leq\varphi(\Succ(n))}{%
      \FormulaRefAuto{PeanoLtAddOneLeq}[thm]{5,6},\ \FormulaRefAuto{PeanoAddOneSucc}[thm]}
    \proofstepwidestar[1,2,4]{%
      n<\varphi(\Succ(n))}{%
      \FormulaRefAuto{PeanoLtLeqTransitive}[thm]{4,5,7,8}}
    \proofstepwidestar[1,2,4]{%
      \Succ(n)\leq\varphi(\Succ(n))}{%
      \FormulaRefAuto{PeanoLtAddOneLeq}[thm]{4,5,9},\ \FormulaRefAuto{PeanoAddOneSucc}[thm]}
    \proofstepwidestar[1,2]{%
      \forall n\in\mathbb N\,(n\leq\varphi(n))}{%
      \FormulaRefAuto{PeanoInductionRule}[ax]{3,\rUI{\rRI{4,10}}}}
  \closeproofpartwideindR

  \proofpartwide{Übertragung der Epsilon-Schranke}
    \proofstepwidestar[1]{%
      a\colon\mathbb N\to\mathbb R\land L\in\mathbb R}{%
      \rA}
    \proofstepwidestar[2]{%
      a_n\to L}{%
      \rA}
    \proofstepwidestar[3]{%
      \varphi\colon\mathbb N\to\mathbb N}{%
      \rA}
    \proofstepwidestar[4]{%
      \forall n\in\mathbb N\,(\varphi(n)<\varphi(\Succ(n)))}{%
      \rA}
    \proofstepwidestar[5]{%
      \varepsilon\in\mathbb R\land\varepsilon>0}{%
      \rA}
    \proofstepwidestar[1,2,5]{%
      \exists N\in\mathbb N\,\forall k\in\mathbb N\,(N\leq k\rightarrow|a_k-L|<\varepsilon)}{%
      \FormulaRefAuto{RealSequenceConvergenceDef}[def]{1,2,5}}
    \proofstepwidestar[7]{%
      \begin{aligned}[t]&N\in\mathbb N\land{}\\[-2pt]&\forall k\in\mathbb N\,(N\leq k\rightarrow|a_k-L|<\varepsilon)\end{aligned}}{%
      \rA}
    \proofstepwidestar[8]{%
      n\in\mathbb N\land N\leq n}{%
      \rA}
    \proofstepwidestar[3,4,8]{%
      n\leq\varphi(n)}{%
      \rRE{\rUE{\FormulaRefAuto{B21IncreasingIndexDominates}[thm]{3,4}},\rAEa{8}}}
    \proofstepwidestar[3,8]{%
      \varphi(n)\in\mathbb N}{%
      \FormulaRefAuto{F\colon A\to B\dsep x\in A\vdash F(x)\in B}[thm]{3,\rAEa{8}}}
    \proofstepwidestar[3,4,7,8]{%
      N\leq\varphi(n)}{%
      \FormulaRefAuto{PeanoLeqTransitive}[thm]{7,8,10,9}}
    \proofstepwidestar[3,4,7,8]{%
      |a_{\varphi(n)}-L|<\varepsilon}{%
      \rRE{\rRE{\rUE{\rAEb{7}},10},11}}
    \proofstepwidestar[3,4,7]{%
      \exists N\in\mathbb N\,\forall n\in\mathbb N\,(N\leq n\rightarrow|a_{\varphi(n)}-L|<\varepsilon)}{%
      \rEI{\rAI{\rAEa{7},\rUI{\rRI{8,12}}}}}
    \proofstepwidestar[1,2,3,4,5]{%
      \exists N\in\mathbb N\,\forall n\in\mathbb N\,(N\leq n\rightarrow|a_{\varphi(n)}-L|<\varepsilon)}{%
      \rEE{6,7,13}}
    \proofstepwidestar[1,2,3,4]{%
      a_{\varphi(n)}\to L}{%
      \FormulaRefAuto{RealSequenceConvergenceDef}[def]{\FormulaRefAuto{CompositionDef}[def]{3,1},\rUI{\rRI{5,14}}}}
  \closeproofpartwide
\end{tabproofsplitwide}

\section{Nullfolgen und Abschätzungen}

\FormulaDefDeltaK[Nullfolge]{%
  a\colon\mathbb N\to\mathbb R\vdash
  \bigl(a\text{ ist eine Nullfolge}\bigr)
  \coloneqq(a_n\to0)%
}{NullSequenceDef}{%
  \DeltaRow{Funktionen}{a}
}

\FormulaThmDeltaK[Konvergenz als Nullfolgeneigenschaft]{%
  a\colon\mathbb N\to\mathbb R\dsep L\in\mathbb R\vdash
  a_n\to L\leftrightarrow(a-L)_n\to0%
}{ConvergenceIffDifferenceNullSequence}{%
  \DeltaRow{Mengen}{L\dsep n}
  \DeltaRow{Funktionen}{a\dsep a-L}
}
\begin{tabproofsplitwide}
  \proofpartwide{Identische Epsilon-Bedingungen}
    \proofstepwidestar[1]{%
      a\colon\mathbb N\to\mathbb R}{%
      \rA}
    \proofstepwidestar[2]{%
      L\in\mathbb R}{%
      \rA}
    \proofstepwidestar[3]{%
      n\in\mathbb N}{%
      \rA}
    \proofstepwidestar[1,2,3]{%
      (a-L)_n=a_n-L}{%
      \FormulaRefAuto{PointwiseRealSequenceOperationsDef}[def]{1,2,3}}
    \proofstepwidestar[1,2,3]{%
      |(a-L)_n-0|=|a_n-L|}{%
      \FormulaRefAuto{RealCutAdditiveLaws}[thm]{\FormulaRefAuto{IndexedFamilyCoordinateTyping}[thm]{1,3},2},\ \rIE{4}}
    \proofstepwidestar[1,2]{%
      \forall n\in\mathbb N\,(|(a-L)_n-0|=|a_n-L|)}{%
      \rUI{\rRI{3,5}}}
    \proofstepwidestar[1,2]{%
      a_n\to L\leftrightarrow(a-L)_n\to0}{%
      \FormulaRefAuto{RealSequenceConvergenceDef}[def]{1,2,\FormulaRefAuto{PointwiseRealSequenceOperationsDef}[def],6}}
  \closeproofpartwide
\end{tabproofsplitwide}

\FormulaThmDeltaK[Majorantenkriterium für Nullfolgen]{%
  a,b\colon\mathbb N\to\mathbb R\dsep
  b_n\to0\dsep
  \exists N_0\in\mathbb N\,
    \forall n\in\mathbb N\,
      (N_0\leq n\rightarrow|a_n|\leq b_n)
  \vdash a_n\to0%
}{NullSequenceMajorantCriterion}{%
  \DeltaRow{Mengen}{N_0\dsep n}
  \DeltaRow{Funktionen}{a\dsep b}
}
\begin{tabproofsplitwide}
  \proofpartwide{Gemeinsamer Anfangsindex und Majorante}
    \proofstepwidestar[1]{%
      a,b\colon\mathbb N\to\mathbb R}{%
      \rA}
    \proofstepwidestar[2]{%
      b_n\to0}{%
      \rA}
    \proofstepwidestar[3]{%
      \exists N_0\in\mathbb N\,\forall k\in\mathbb N\,(N_0\leq k\rightarrow|a_k|\leq b_k)}{%
      \rA}
    \proofstepwidestar[4]{%
      \varepsilon\in\mathbb R\land\varepsilon>0}{%
      \rA}
    \proofstepwidestar[1,2,4]{%
      \exists N\in\mathbb N\,\forall k\in\mathbb N\,(N\leq k\rightarrow|b_k|<\varepsilon)}{%
      \FormulaRefAuto{RealSequenceConvergenceDef}[def]{1,2,4},\ \FormulaRefAuto{RealCutAdditiveLaws}[thm]}
    \proofstepwidestar[6]{%
      \begin{aligned}[t]&N_0\in\mathbb N\land{}\\[-2pt]&\forall k\in\mathbb N\,(N_0\leq k\rightarrow|a_k|\leq b_k)\end{aligned}}{%
      \rA}
    \proofstepwidestar[7]{%
      \begin{aligned}[t]&N\in\mathbb N\land{}\\[-2pt]&\forall k\in\mathbb N\,(N\leq k\rightarrow|b_k|<\varepsilon)\end{aligned}}{%
      \rA}
    \proofstepwidestar[6,7]{%
      N_0+N\in\mathbb N\land(N_0\leq N_0+N\land N\leq N_0+N)}{%
      \FormulaRefAuto{B21CommonNaturalBound}[thm]{6,7}}
    \proofstepwidestar[9]{%
      n\in\mathbb N\land N_0+N\leq n}{%
      \rA}
    \proofstepwidestar[6,7,9]{%
      N_0\leq n\land N\leq n}{%
      \FormulaRefAuto{PeanoLeqTransitive}[thm]{8,9}}
    \proofstepwidestar[1,6,7,9]{%
      |a_n|\leq b_n\leq|b_n|<\varepsilon}{%
      \rUE{\rAEb{6}},\ \rUE{\rAEb{7}},\ 10,\ \FormulaRefAuto{RealAbsoluteValue}[def]{\FormulaRefAuto{IndexedFamilyCoordinateTyping}[thm]{1,9}}}
    \proofstepwidestar[1,6,7]{%
      \exists M\in\mathbb N\,\forall n\in\mathbb N\,(M\leq n\rightarrow|a_n|<\varepsilon)}{%
      \rEI{\rAI{\rAEa{8},\rUI{\rRI{9,11}}}}}
    \proofstepwidestar[1,2,4,6]{%
      \exists M\in\mathbb N\,\forall n\in\mathbb N\,(M\leq n\rightarrow|a_n|<\varepsilon)}{%
      \rEE{5,7,12}}
    \proofstepwidestar[1,2,3,4]{%
      \exists M\in\mathbb N\,\forall n\in\mathbb N\,(M\leq n\rightarrow|a_n|<\varepsilon)}{%
      \rEE{3,6,13}}
    \proofstepwidestar[1,2,3]{%
      a_n\to0}{%
      \FormulaRefAuto{RealSequenceConvergenceDef}[def]{1,\rUI{\rRI{4,14}}},\ \FormulaRefAuto{RealCutAdditiveLaws}[thm]}
  \closeproofpartwide
\end{tabproofsplitwide}

\FormulaThmDeltaK[Einschließungsprinzip]{%
  \begin{aligned}[t]
  &a,b,c\colon\mathbb N\to\mathbb R\dsep L\in\mathbb R\dsep
    a_n\to L\dsep c_n\to L\dsep{}\\[-2pt]
  &\quad\exists N_0\in\mathbb N\,
    \forall n\in\mathbb N\,
      \bigl(N_0\leq n\rightarrow
        (a_n\leq b_n\land b_n\leq c_n)\bigr)
    \vdash b_n\to L
  \end{aligned}%
}{RealSequenceSqueezeTheorem}{%
  \DeltaRow{Mengen}{L\dsep N_0\dsep n}
  \DeltaRow{Funktionen}{a\dsep b\dsep c}
}
\begin{tabproofsplitwide}
  \proofpartwide{Einschließung im Epsilon-Intervall}
    \proofstepwidestar[1]{%
      a,b,c\colon\mathbb N\to\mathbb R\land L\in\mathbb R}{%
      \rA}
    \proofstepwidestar[2]{%
      a_n\to L\land c_n\to L}{%
      \rA}
    \proofstepwidestar[3]{%
      \exists N_0\in\mathbb N\,\forall k\in\mathbb N\,(N_0\leq k\rightarrow(a_k\leq b_k\land b_k\leq c_k))}{%
      \rA}
    \proofstepwidestar[4]{%
      \varepsilon\in\mathbb R\land\varepsilon>0}{%
      \rA}
    \proofstepwidestar[1,2,3,4]{%
      \begin{aligned}[t]&\exists N\in\mathbb N\,\forall k\in\mathbb N\,\bigl(N\leq k\rightarrow{}\\[-2pt]&\quad(|a_k-L|<\varepsilon\land|c_k-L|<\varepsilon\\[-2pt]&\qquad\land a_k\leq b_k\land b_k\leq c_k)\bigr)\end{aligned}}{%
      \FormulaRefAuto{RealSequenceConvergenceDef}[def]{1,2,4},\ \FormulaRefAuto{B21CommonNaturalBound}[thm],\ \rEE{3}}
    \proofstepwidestar[6]{%
      \begin{aligned}[t]&N\in\mathbb N\land\forall k\in\mathbb N\,\bigl(N\leq k\rightarrow{}\\[-2pt]&\quad(|a_k-L|<\varepsilon\land|c_k-L|<\varepsilon\\[-2pt]&\qquad\land a_k\leq b_k\land b_k\leq c_k)\bigr)\end{aligned}}{%
      \rA}
    \proofstepwidestar[7]{%
      n\in\mathbb N\land N\leq n}{%
      \rA}
    \proofstepwidestar[1,4,6,7]{%
      L-\varepsilon<a_n\leq b_n\leq c_n<L+\varepsilon}{%
      \rRE{\rRE{\rUE{\rAEb{6}},\rAEa{7}},\rAEb{7}},\ \FormulaRefAuto{RealEpsilonNeighborhoodInterval}[thm]{1,4}}
    \proofstepwidestar[1,4,6,7]{%
      |b_n-L|<\varepsilon}{%
      \FormulaRefAuto{RealEpsilonNeighborhoodInterval}[thm]{1,4,8},\ \FormulaRefAuto{RealCutTotalOrder}[thm]}
    \proofstepwidestar[1,4,6]{%
      \exists N\in\mathbb N\,\forall n\in\mathbb N\,(N\leq n\rightarrow|b_n-L|<\varepsilon)}{%
      \rEI{\rAI{\rAEa{6},\rUI{\rRI{7,9}}}}}
    \proofstepwidestar[1,2,3,4]{%
      \exists N\in\mathbb N\,\forall n\in\mathbb N\,(N\leq n\rightarrow|b_n-L|<\varepsilon)}{%
      \rEE{5,6,10}}
    \proofstepwidestar[1,2,3]{%
      b_n\to L}{%
      \FormulaRefAuto{RealSequenceConvergenceDef}[def]{1,\rUI{\rRI{4,11}}}}
  \closeproofpartwide
\end{tabproofsplitwide}

\section{Beschränktheit und Ordnung}

\FormulaDefDeltaK[Beschränkte reelle Folge]{%
  a\colon\mathbb N\to\mathbb R\vdash
  \operatorname{Beschr}(a)\coloneqq
  \exists C\in\mathbb R\,
    \bigl(C\geq0\land
      \forall n\in\mathbb N\,(|a_n|\leq C)\bigr)%
}{BoundedRealSequenceDef}{%
  \DeltaRow{Mengen}{C\dsep n}
  \DeltaRow{Funktionen}{a}
  \DeltaRow{Neues Prädikat}{\operatorname{Beschr}}
}

\FormulaThmDeltaK[Konvergente Folgen sind beschränkt]{%
  a\colon\mathbb N\to\mathbb R\dsep L\in\mathbb R\dsep a_n\to L
  \vdash\operatorname{Beschr}(a)%
}{ConvergentRealSequenceBounded}{%
  \DeltaRow{Mengen}{L\dsep C\dsep N\dsep n}
  \DeltaRow{Funktionen}{a}
}
\begin{tabproofsplitwide}
  \proofpartwideindR[Anfangsschranke: Induktionsanfang]{%
    a\colon\mathbb N\to\mathbb R\vdash\exists C\in\mathbb R\,(0\leq C\land\forall k\in\NatSeg 0\,(|a_k|\leq C))
  }{a\colon\mathbb N\to\mathbb R\vdash\exists C\in\mathbb R\,(0\leq C\land\forall k\in\NatSeg 0\,(|a_k|\leq C))}[B21FinitePrefixBase]
    \proofstepwidestar[1]{%
      a\colon\mathbb N\to\mathbb R}{%
      \rA}
    \proofstepwidestar[2]{%
      k\in\NatSeg 0}{%
      \rA}
    \proofstepwidestar[2]{%
      \bot}{%
      \FormulaRefAuto{PeanoNatSegZero}[thm],\ \text{Def. }\ensuremath{\varnothing}(2)}
    \proofstepwidestar[1,2]{%
      |a_k|\leq0}{%
      \FormulaRefAuto{\bot\vdash P}[thm]{3}}
    \proofstepwidestar[1]{%
      \forall k\in\NatSeg 0\,(|a_k|\leq0)}{%
      \rUI{\rRI{2,4}}}
    \proofstepwidestar[1]{%
      \exists C\in\mathbb R\,(0\leq C\land\forall k\in\NatSeg 0\,(|a_k|\leq C))}{%
      \rEI{\rAI{\FormulaRefAuto{RealCutAdditiveClosure}[thm],\rAI{\FormulaRefAuto{RealCutTotalOrder}[thm],5}}}}
  \closeproofpartwideindR

  \proofpartwideindR[Anfangsschranke: Induktionsschritt]{%
    \begin{aligned}[t]&a\colon\mathbb N\to\mathbb R\dsep N\in\mathbb N\dsep{}\\[-2pt]&\exists C\in\mathbb R\,(0\leq C\land\forall k\in\NatSeg N\,(|a_k|\leq C))\vdash{}\\[-2pt]&\exists C\in\mathbb R\,(0\leq C\land\forall k\in\NatSeg{\Succ(N)}\,(|a_k|\leq C))\end{aligned}
  }{\begin{aligned}[t]&a\colon\mathbb N\to\mathbb R\dsep N\in\mathbb N\dsep{}\\[-2pt]&\exists C\in\mathbb R\,(0\leq C\land\forall k\in\NatSeg N\,(|a_k|\leq C))\vdash{}\\[-2pt]&\exists C\in\mathbb R\,(0\leq C\land\forall k\in\NatSeg{\Succ(N)}\,(|a_k|\leq C))\end{aligned}}[B21FinitePrefixStep]
    \proofstepwidestar[1]{%
      a\colon\mathbb N\to\mathbb R}{%
      \rA}
    \proofstepwidestar[2]{%
      N\in\mathbb N}{%
      \rA}
    \proofstepwidestar[3]{%
      \exists C\in\mathbb R\,(0\leq C\land\forall k\in\NatSeg N\,(|a_k|\leq C))}{%
      \rA}
    \proofstepwidestar[4]{%
      C\in\mathbb R\land(0\leq C\land\forall k\in\NatSeg N\,(|a_k|\leq C))}{%
      \rA}
    \proofstepwidestar[1,2]{%
      a_N\in\mathbb R}{%
      \FormulaRefAuto{IndexedFamilyCoordinateTyping}[thm]{1,2}}
    \proofstepwidestar[1,2,4]{%
      C+|a_N|\in\mathbb R\land0\leq C+|a_N|}{%
      \FormulaRefAuto{RealAbsoluteValueNonnegativeZero}[thm]{5},\ \FormulaRefAuto{RealCutAdditiveClosure}[thm]{4,5},\ \FormulaRefAuto{RealCutOrderArithmeticCompatibility}[thm]}
    \proofstepwidestar[7]{%
      k\in\NatSeg{\Succ(N)}}{%
      \rA}
    \proofstepwidestar[2,7]{%
      k\in\NatSeg N\lor k=N}{%
      \FormulaRefAuto{PeanoNatSegSuccElim}[thm]{2,7}}
    \proofstepwidestar[9]{%
      k\in\NatSeg N}{%
      \rA}
    \proofstepwidestar[1,2,4,9]{%
      |a_k|\leq C+|a_N|}{%
      \rRE{\rUE{\rAEb{\rAEb{4}}},9},\ \FormulaRefAuto{RealAbsoluteValueNonnegativeZero}[thm]{5},\ \FormulaRefAuto{RealCutOrderArithmeticCompatibility}[thm]}
    \proofstepwidestar[11]{%
      k=N}{%
      \rA}
    \proofstepwidestar[1,2,4,11]{%
      |a_k|\leq C+|a_N|}{%
      \FormulaRefAuto{RealCutOrderArithmeticCompatibility}[thm]{4,5},\ \rIE{11}}
    \proofstepwidestar[1,2,4,7]{%
      |a_k|\leq C+|a_N|}{%
      \rOE{8,9,10,11,12}}
    \proofstepwidestar[1,2,4]{%
      \forall k\in\NatSeg{\Succ(N)}\,(|a_k|\leq C+|a_N|)}{%
      \rUI{\rRI{7,13}}}
    \proofstepwidestar[1,2,4]{%
      \exists C\in\mathbb R\,(0\leq C\land\forall k\in\NatSeg{\Succ(N)}\,(|a_k|\leq C))}{%
      \rEI{\rAI{\rAEa{6},\rAI{\rAEb{6},14}}}}
    \proofstepwidestar[1,2,3]{%
      \exists C\in\mathbb R\,(0\leq C\land\forall k\in\NatSeg{\Succ(N)}\,(|a_k|\leq C))}{%
      \rEE{3,4,15}}
  \closeproofpartwideindR

  \proofpartwideindR[Anfangsschranke: Induktionsschluss]{%
    a\colon\mathbb N\to\mathbb R\dsep N\in\mathbb N\vdash\exists C\in\mathbb R\,(0\leq C\land\forall k\in\NatSeg N\,(|a_k|\leq C))
  }{a\colon\mathbb N\to\mathbb R\dsep N\in\mathbb N\vdash\exists C\in\mathbb R\,(0\leq C\land\forall k\in\NatSeg N\,(|a_k|\leq C))}[B21FinitePrefixBound]
    \proofstepwidestar[1]{%
      a\colon\mathbb N\to\mathbb R}{%
      \rA}
    \proofstepwidestar[2]{%
      N\in\mathbb N}{%
      \rA}
    \proofstepwidestar[1,2]{%
      \exists C\in\mathbb R\,(0\leq C\land\forall k\in\NatSeg N\,(|a_k|\leq C))}{%
      \FormulaRefAuto{PeanoInductionRule}[ax]{\FormulaRefAuto{B21FinitePrefixBase}[thm]{1},\FormulaRefAuto{B21FinitePrefixStep}[thm]{1},2}}
  \closeproofpartwideindR

  \proofpartwideindR[Eine Restschranke liefert eine globale Schranke]{%
    \begin{aligned}[t]&a\colon\mathbb N\to\mathbb R\dsep N\in\mathbb N\dsep C\in\mathbb R\dsep0\leq C\dsep{}\\[-2pt]&\forall k\in\mathbb N\,(N\leq k\rightarrow|a_k|\leq C)\vdash\operatorname{Beschr}(a)\end{aligned}
  }{\begin{aligned}[t]&a\colon\mathbb N\to\mathbb R\dsep N\in\mathbb N\dsep C\in\mathbb R\dsep0\leq C\dsep{}\\[-2pt]&\forall k\in\mathbb N\,(N\leq k\rightarrow|a_k|\leq C)\vdash\operatorname{Beschr}(a)\end{aligned}}[B21EventuallyBoundedIsBounded]
    \proofstepwidestar[1]{%
      a\colon\mathbb N\to\mathbb R}{%
      \rA}
    \proofstepwidestar[2]{%
      N\in\mathbb N\land(C\in\mathbb R\land0\leq C)}{%
      \rA}
    \proofstepwidestar[3]{%
      \forall k\in\mathbb N\,(N\leq k\rightarrow|a_k|\leq C)}{%
      \rA}
    \proofstepwidestar[1,2]{%
      \exists D\in\mathbb R\,(0\leq D\land\forall k\in\NatSeg N\,(|a_k|\leq D))}{%
      \FormulaRefAuto{B21FinitePrefixBound}[thm]{1,\rAEa{2}}}
    \proofstepwidestar[5]{%
      D\in\mathbb R\land(0\leq D\land\forall k\in\NatSeg N\,(|a_k|\leq D))}{%
      \rA}
    \proofstepwidestar[6]{%
      k\in\mathbb N}{%
      \rA}
    \proofstepwidestar[2,6]{%
      k\in\NatSeg N\lor N\leq k}{%
      \FormulaRefAuto{PeanoOrderTrichotomy}[thm]{2,6},\ \FormulaRefAuto{PeanoNatSegMembership}[thm]}
    \proofstepwidestar[8]{%
      k\in\NatSeg N}{%
      \rA}
    \proofstepwidestar[2,5,8]{%
      |a_k|\leq C+D}{%
      \rRE{\rUE{\rAEb{\rAEb{5}}},8},\ \FormulaRefAuto{RealCutOrderArithmeticCompatibility}[thm]{2,5}}
    \proofstepwidestar[10]{%
      N\leq k}{%
      \rA}
    \proofstepwidestar[2,3,5,6,10]{%
      |a_k|\leq C+D}{%
      \rRE{\rRE{\rUE{3},6},10},\ \FormulaRefAuto{RealCutOrderArithmeticCompatibility}[thm]{2,5}}
    \proofstepwidestar[2,3,5,6]{%
      |a_k|\leq C+D}{%
      \rOE{7,8,9,10,11}}
    \proofstepwidestar[2,3,5]{%
      \forall k\in\mathbb N\,(|a_k|\leq C+D)}{%
      \rUI{\rRI{6,12}}}
    \proofstepwidestar[1,2,3,5]{%
      \operatorname{Beschr}(a)}{%
      \FormulaRefAuto{BoundedRealSequenceDef}[def]{1,\rEI{\rAI{\FormulaRefAuto{RealCutAdditiveClosure}[thm]{2,5},\rAI{\FormulaRefAuto{RealCutOrderArithmeticCompatibility}[thm]{2,5},13}}}}}
    \proofstepwidestar[1,2,3]{%
      \operatorname{Beschr}(a)}{%
      \rEE{4,5,14}}
  \closeproofpartwideindR

  \proofpartwide{Konvergenz liefert eine Restschranke}
    \proofstepwidestar[1]{%
      a\colon\mathbb N\to\mathbb R\land L\in\mathbb R}{%
      \rA}
    \proofstepwidestar[2]{%
      a_n\to L}{%
      \rA}
    \proofstepwidestar[1,2]{%
      \exists N\in\mathbb N\,\forall k\in\mathbb N\,(N\leq k\rightarrow|a_k-L|<1)}{%
      \FormulaRefAuto{RealSequenceConvergenceDef}[def]{1,2,\FormulaRefAuto{RealCutsCompleteOrderedField}[thm]}}
    \proofstepwidestar[4]{%
      N\in\mathbb N\land\forall k\in\mathbb N\,(N\leq k\rightarrow|a_k-L|<1)}{%
      \rA}
    \proofstepwidestar[5]{%
      k\in\mathbb N\land N\leq k}{%
      \rA}
    \proofstepwidestar[1,4,5]{%
      |a_k|\leq|a_k-L|+|L|<1+|L|}{%
      \FormulaRefAuto{RealAbsoluteValueTriangle}[thm]{1,5},\ \rRE{\rRE{\rUE{\rAEb{4}},\rAEa{5}},\rAEb{5}},\ \FormulaRefAuto{RealCutOrderArithmeticCompatibility}[thm]}
    \proofstepwidestar[1,4]{%
      \forall k\in\mathbb N\,(N\leq k\rightarrow|a_k|\leq1+|L|)}{%
      \rUI{\rRI{5,6}},\ \FormulaRefAuto{RealCutTotalOrder}[thm]}
    \proofstepwidestar[1,4]{%
      \operatorname{Beschr}(a)}{%
      \FormulaRefAuto{B21EventuallyBoundedIsBounded}[thm]{1,4,\FormulaRefAuto{RealAbsoluteValueNonnegativeZero}[thm]{1},\FormulaRefAuto{RealCutAdditiveClosure}[thm],7}}
    \proofstepwidestar[1,2]{%
      \operatorname{Beschr}(a)}{%
      \rEE{3,4,8}}
  \closeproofpartwide
\end{tabproofsplitwide}

\FormulaThmDeltaK[Ordnung bleibt im Grenzübergang erhalten]{%
  \begin{aligned}[t]
  &a,b\colon\mathbb N\to\mathbb R\dsep A,B\in\mathbb R\dsep
    a_n\to A\dsep b_n\to B\dsep{}\\[-2pt]
  &\quad\exists N_0\in\mathbb N\,
    \forall n\in\mathbb N\,(N_0\leq n\rightarrow a_n\leq b_n)
    \vdash A\leq B
  \end{aligned}%
}{RealSequenceLimitPreservesOrder}{%
  \DeltaRow{Mengen}{A\dsep B\dsep N_0\dsep n}
  \DeltaRow{Funktionen}{a\dsep b}
}
\begin{tabproofsplitwide}
  \proofpartwide{Widerspruch bei umgekehrter Grenzwertordnung}
    \proofstepwidestar[1]{%
      a,b\colon\mathbb N\to\mathbb R\land A,B\in\mathbb R}{%
      \rA}
    \proofstepwidestar[2]{%
      a_n\to A\land b_n\to B}{%
      \rA}
    \proofstepwidestar[3]{%
      \exists N_0\in\mathbb N\,\forall k\in\mathbb N\,(N_0\leq k\rightarrow a_k\leq b_k)}{%
      \rA}
    \proofstepwidestar[4]{%
      \neg(A\leq B)}{%
      \rA}
    \proofstepwidestar[1,4]{%
      B<A\land0<(A-B)/3}{%
      \FormulaRefAuto{RealCutTotalOrder}[thm]{1,4},\ \FormulaRefAuto{RealCutsCompleteOrderedField}[thm]}
    \proofstepwidestar[1,2,3,4]{%
      \begin{aligned}[t]&\exists N\in\mathbb N\,\bigl(|a_N-A|<(A-B)/3\\[-2pt]&\qquad\land|b_N-B|<(A-B)/3\land a_N\leq b_N\bigr)\end{aligned}}{%
      \FormulaRefAuto{RealSequenceConvergenceDef}[def]{1,2,5},\ \FormulaRefAuto{B21CommonNaturalBound}[thm],\ \rEE{3}}
    \proofstepwidestar[7]{%
      \begin{aligned}[t]&N\in\mathbb N\land\bigl(|a_N-A|<(A-B)/3\\[-2pt]&\qquad\land|b_N-B|<(A-B)/3\land a_N\leq b_N\bigr)\end{aligned}}{%
      \rA}
    \proofstepwidestar[1,4,7]{%
      a_N>A-(A-B)/3}{%
      \FormulaRefAuto{RealEpsilonNeighborhoodInterval}[thm]{1,5,7}}
    \proofstepwidestar[1,4,7]{%
      b_N<B+(A-B)/3}{%
      \FormulaRefAuto{RealEpsilonNeighborhoodInterval}[thm]{1,5,7}}
    \proofstepwidestar[1,4]{%
      B+(A-B)/3<A-(A-B)/3}{%
      \FormulaRefAuto{RealCutsCompleteOrderedField}[thm]{1,5}}
    \proofstepwidestar[1,4,7]{%
      \bot}{%
      \FormulaRefAuto{RealCutTotalOrder}[thm]{7,8,9,10}}
    \proofstepwidestar[1,2,3,4]{%
      \bot}{%
      \rEE{6,7,11}}
    \proofstepwidestar[1,2,3]{%
      A\leq B}{%
      \rCE{4,12}}
  \closeproofpartwide
\end{tabproofsplitwide}

Aus einer schließlich strikten Ungleichung \(a_n<b_n\) folgt im Allgemeinen
nur \(A\leq B\), nicht \(A<B\). Beispielsweise gilt
\(0<1/(n+1)\), obwohl beide Seiten gegen \(0\) konvergieren.

\FormulaDefDeltaK[Monotone reelle Folgen]{%
  a\colon\mathbb N\to\mathbb R\vdash
  \begin{aligned}[t]
  \operatorname{Mon}_{\uparrow}(a)&\coloneqq
    \forall n\in\mathbb N\,(a_n\leq a_{\Succ(n)}),\\[-2pt]
  \operatorname{Mon}_{\downarrow}(a)&\coloneqq
    \forall n\in\mathbb N\,(a_{\Succ(n)}\leq a_n)
  \end{aligned}%
}{MonotoneRealSequenceDef}{%
  \DeltaRow{Mengen}{n}
  \DeltaRow{Funktionen}{a}
  \DeltaRow{Neue Prädikate}{\operatorname{Mon}_{\uparrow}\dsep
    \operatorname{Mon}_{\downarrow}}
}

\FormulaThmDeltaK[Monotone Konvergenz]{%
  \begin{aligned}[t]
  &a\colon\mathbb N\to\mathbb R\dsep
    \operatorname{Mon}_{\uparrow}(a)\dsep
    \exists C\in\mathbb R\,\forall n\in\mathbb N\,(a_n\leq C)
    \vdash
    a_n\to\sup_{\mathbb R,\RealLe}\bigl(a[\mathbb N]\bigr),\\[-2pt]
  &a\colon\mathbb N\to\mathbb R\dsep
    \operatorname{Mon}_{\downarrow}(a)\dsep
    \exists C\in\mathbb R\,\forall n\in\mathbb N\,(C\leq a_n)
    \vdash
    a_n\to\inf_{\mathbb R,\RealLe}\bigl(a[\mathbb N]\bigr).
  \end{aligned}%
}{MonotoneRealSequenceConvergence}{%
  \DeltaRow{Mengen}{C\dsep n\dsep a[\mathbb N]}
  \DeltaRow{Funktionen}{a}
  \DeltaRow{Relationsoperatoren}{\RealLe}
}
\begin{tabproofsplitwide}
  \proofpartwideindR[Wachsende Folgen: Vergleich beliebiger Indizes]{%
    a\colon\mathbb N\to\mathbb R\dsep\operatorname{Mon}_{\uparrow}(a)\dsep N,n\in\mathbb N\dsep N\leq n\vdash a_N\leq a_n
  }{a\colon\mathbb N\to\mathbb R\dsep\operatorname{Mon}_{\uparrow}(a)\dsep N,n\in\mathbb N\dsep N\leq n\vdash a_N\leq a_n}[B21MonotoneIncreasingNesting]
    \proofstepwidestar[1]{%
      a\colon\mathbb N\to\mathbb R}{%
      \rA}
    \proofstepwidestar[2]{%
      \operatorname{Mon}_{\uparrow}(a)}{%
      \rA}
    \proofstepwidestar[3]{%
      (N\in\mathbb N\land n\in\mathbb N)\land N\leq n}{%
      \rA}
    \proofstepwidestar[1,3]{%
      a_N\leq a_{N+0}}{%
      \FormulaRefAuto{PeanoAddRightZero}[thm]{3},\ \FormulaRefAuto{RealCutTotalOrder}[thm]{\FormulaRefAuto{IndexedFamilyCoordinateTyping}[thm]{1,3}}}
    \proofstepwidestar[5]{%
      k\in\mathbb N\land a_N\leq a_{N+k}}{%
      \rA}
    \proofstepwidestar[1,2,3,5]{%
      a_{N+k}\leq a_{N+(k+1)}}{%
      \FormulaRefAuto{MonotoneRealSequenceDef}[def]{1,2},\ \FormulaRefAuto{PeanoAddClosure}[thm]{3,5},\ \FormulaRefAuto{PeanoAddAddOneRight}[thm],\ \FormulaRefAuto{PeanoAddOneSucc}[thm]}
    \proofstepwidestar[1,2,3,5]{%
      a_N\leq a_{N+(k+1)}}{%
      \FormulaRefAuto{RealCutTotalOrder}[thm]{5,6}}
    \proofstepwidestar[1,2,3]{%
      \forall k\in\mathbb N\,(a_N\leq a_{N+k})}{%
      \FormulaRefAuto{PeanoInductionRuleAddOne}[thm]{4,\rUI{\rRI{5,7}}}}
    \proofstepwidestar[3]{%
      \exists k\in\mathbb N\,(N+k=n)}{%
      \FormulaRefAuto{m\in\mathbb N\dsep n\in\mathbb N\vdash m\leq n\leftrightarrow\exists k\in\mathbb N\,(m+k=n)}[thm]{3}}
    \proofstepwidestar[10]{%
      k\in\mathbb N\land N+k=n}{%
      \rA}
    \proofstepwidestar[1,2,3,10]{%
      a_N\leq a_n}{%
      \rIE{\rAEb{10},\rRE{\rUE{8},\rAEa{10}}}}
    \proofstepwidestar[1,2,3]{%
      a_N\leq a_n}{%
      \rEE{9,10,11}}
  \closeproofpartwideindR

  \proofpartwideindR[Fallende Folgen: Vergleich beliebiger Indizes]{%
    a\colon\mathbb N\to\mathbb R\dsep\operatorname{Mon}_{\downarrow}(a)\dsep N,n\in\mathbb N\dsep N\leq n\vdash a_n\leq a_N
  }{a\colon\mathbb N\to\mathbb R\dsep\operatorname{Mon}_{\downarrow}(a)\dsep N,n\in\mathbb N\dsep N\leq n\vdash a_n\leq a_N}[B21MonotoneDecreasingNesting]
    \proofstepwidestar[1]{%
      a\colon\mathbb N\to\mathbb R}{%
      \rA}
    \proofstepwidestar[2]{%
      \operatorname{Mon}_{\downarrow}(a)}{%
      \rA}
    \proofstepwidestar[3]{%
      (N\in\mathbb N\land n\in\mathbb N)\land N\leq n}{%
      \rA}
    \proofstepwidestar[1]{%
      -a\colon\mathbb N\to\mathbb R}{%
      \FormulaRefAuto{PointwiseRealSequenceOperationsDef}[def]{1}}
    \proofstepwidestar[1,2]{%
      \operatorname{Mon}_{\uparrow}(-a)}{%
      \FormulaRefAuto{MonotoneRealSequenceDef}[def]{1,2},\ \FormulaRefAuto{PointwiseRealSequenceOperationsDef}[def],\ \FormulaRefAuto{RealCutNegationOrder}[thm]}
    \proofstepwidestar[1,2,3]{%
      (-a)_N\leq(-a)_n}{%
      \FormulaRefAuto{B21MonotoneIncreasingNesting}[thm]{4,5,3}}
    \proofstepwidestar[1,2,3]{%
      a_n\leq a_N}{%
      \FormulaRefAuto{PointwiseRealSequenceOperationsDef}[def]{1,6},\ \FormulaRefAuto{RealCutNegationOrder}[thm]}
  \closeproofpartwideindR

  \proofpartwide{Konvergenz gegen das Supremum}
    \proofstepwidestar[1]{%
      a\colon\mathbb N\to\mathbb R}{%
      \rA}
    \proofstepwidestar[2]{%
      \operatorname{Mon}_{\uparrow}(a)}{%
      \rA}
    \proofstepwidestar[3]{%
      \exists C\in\mathbb R\,\forall k\in\mathbb N\,(a_k\leq C)}{%
      \rA}
    \proofstepwidestar[1]{%
      a_0\in a[\mathbb N]\subseteq\mathbb R}{%
      \FormulaRefAuto{IndexedFamilyCoordinateTyping}[thm]{1,\FormulaRefAuto{PeanoZero}[ax]},\ \FormulaRefAuto{IndexedFamilyImageMembership}[thm]{1},\ \FormulaRefAuto{IndexedFamilyImageSubsetCodomain}[thm]{1}}
    \proofstepwidestar[1,3]{%
      s\coloneqq\sup_{\mathbb R,\RealLe}(a[\mathbb N])\in\mathbb R}{%
      \FormulaRefAuto{RealLeastUpperBoundProperty}[thm]{4,3}}
    \proofstepwidestar[6]{%
      \varepsilon\in\mathbb R\land\varepsilon>0}{%
      \rA}
    \proofstepwidestar[1,3,6]{%
      \exists N\in\mathbb N\,(s-\varepsilon<a_N)}{%
      \FormulaRefAuto{RealLeastUpperBoundProperty}[thm]{5},\ \FormulaRefAuto{RealCutsCompleteOrderedField}[thm]{6},\ \FormulaRefAuto{IndexedFamilyImageMembership}[thm]{1}}
    \proofstepwidestar[8]{%
      N\in\mathbb N\land s-\varepsilon<a_N}{%
      \rA}
    \proofstepwidestar[9]{%
      n\in\mathbb N\land N\leq n}{%
      \rA}
    \proofstepwidestar[1,2,3,6,8,9]{%
      s-\varepsilon<a_N\leq a_n\leq s<s+\varepsilon}{%
      \FormulaRefAuto{B21MonotoneIncreasingNesting}[thm]{1,2,8,9},\ \FormulaRefAuto{RealLeastUpperBoundProperty}[thm]{5},\ \FormulaRefAuto{RealCutOrderArithmeticCompatibility}[thm]{6}}
    \proofstepwidestar[1,2,3,6,8,9]{%
      |a_n-s|<\varepsilon}{%
      \FormulaRefAuto{RealEpsilonNeighborhoodInterval}[thm]{1,5,6,10}}
    \proofstepwidestar[1,2,3,6,8]{%
      \exists N\in\mathbb N\,\forall n\in\mathbb N\,(N\leq n\rightarrow|a_n-s|<\varepsilon)}{%
      \rEI{\rAI{\rAEa{8},\rUI{\rRI{9,11}}}}}
    \proofstepwidestar[1,2,3,6]{%
      \exists N\in\mathbb N\,\forall n\in\mathbb N\,(N\leq n\rightarrow|a_n-s|<\varepsilon)}{%
      \rEE{7,8,12}}
    \proofstepwidestar[1,2,3]{%
      a_n\to\sup_{\mathbb R,\RealLe}(a[\mathbb N])}{%
      \FormulaRefAuto{RealSequenceConvergenceDef}[def]{1,5,\rUI{\rRI{6,13}}}}
  \closeproofpartwide

  \proofpartwide{Konvergenz gegen das Infimum}
    \proofstepwidestar[1]{%
      a\colon\mathbb N\to\mathbb R}{%
      \rA}
    \proofstepwidestar[2]{%
      \operatorname{Mon}_{\downarrow}(a)}{%
      \rA}
    \proofstepwidestar[3]{%
      \exists C\in\mathbb R\,\forall k\in\mathbb N\,(C\leq a_k)}{%
      \rA}
    \proofstepwidestar[1]{%
      a_0\in a[\mathbb N]\subseteq\mathbb R}{%
      \FormulaRefAuto{IndexedFamilyCoordinateTyping}[thm]{1,\FormulaRefAuto{PeanoZero}[ax]},\ \FormulaRefAuto{IndexedFamilyImageMembership}[thm]{1},\ \FormulaRefAuto{IndexedFamilyImageSubsetCodomain}[thm]{1}}
    \proofstepwidestar[1,3]{%
      r\coloneqq\inf_{\mathbb R,\RealLe}(a[\mathbb N])\in\mathbb R}{%
      \FormulaRefAuto{RealInfimumCharacterization}[thm]{4,3}}
    \proofstepwidestar[6]{%
      \varepsilon\in\mathbb R\land\varepsilon>0}{%
      \rA}
    \proofstepwidestar[1,3,6]{%
      \exists N\in\mathbb N\,(a_N<r+\varepsilon)}{%
      \FormulaRefAuto{RealInfimumCharacterization}[thm]{5},\ \FormulaRefAuto{RealCutsCompleteOrderedField}[thm]{6},\ \FormulaRefAuto{IndexedFamilyImageMembership}[thm]{1}}
    \proofstepwidestar[8]{%
      N\in\mathbb N\land a_N<r+\varepsilon}{%
      \rA}
    \proofstepwidestar[9]{%
      n\in\mathbb N\land N\leq n}{%
      \rA}
    \proofstepwidestar[1,2,3,6,8,9]{%
      r-\varepsilon<r\leq a_n\leq a_N<r+\varepsilon}{%
      \FormulaRefAuto{B21MonotoneDecreasingNesting}[thm]{1,2,8,9},\ \FormulaRefAuto{RealInfimumCharacterization}[thm]{5},\ \FormulaRefAuto{RealCutOrderArithmeticCompatibility}[thm]{6}}
    \proofstepwidestar[1,2,3,6,8,9]{%
      |a_n-r|<\varepsilon}{%
      \FormulaRefAuto{RealEpsilonNeighborhoodInterval}[thm]{1,5,6,10}}
    \proofstepwidestar[1,2,3,6,8]{%
      \exists N\in\mathbb N\,\forall n\in\mathbb N\,(N\leq n\rightarrow|a_n-r|<\varepsilon)}{%
      \rEI{\rAI{\rAEa{8},\rUI{\rRI{9,11}}}}}
    \proofstepwidestar[1,2,3,6]{%
      \exists N\in\mathbb N\,\forall n\in\mathbb N\,(N\leq n\rightarrow|a_n-r|<\varepsilon)}{%
      \rEE{7,8,12}}
    \proofstepwidestar[1,2,3]{%
      a_n\to\inf_{\mathbb R,\RealLe}(a[\mathbb N])}{%
      \FormulaRefAuto{RealSequenceConvergenceDef}[def]{1,5,\rUI{\rRI{6,13}}}}
  \closeproofpartwide
\end{tabproofsplitwide}

\section{Grenzwertregeln}

\FormulaThmDeltaK[Lineare Grenzwertregeln]{%
  a,b\colon\mathbb N\to\mathbb R\dsep
  A,B,c\in\mathbb R\dsep
  a_n\to A\dsep b_n\to B
  \vdash
  \begin{aligned}[t]
  (a+b)_n&\to A+B,\\[-2pt]
  (a-b)_n&\to A-B,\\[-2pt]
  (ca)_n&\to cA
  \end{aligned}%
}{RealSequenceLinearLimitRules}{%
  \DeltaRow{Mengen}{A\dsep B\dsep c\dsep n}
  \DeltaRow{Funktionen}{a\dsep b\dsep a+b\dsep a-b\dsep ca}
}
\begin{tabproofsplitwide}
  \proofpartwideindR[Summe]{%
    a,b\colon\mathbb N\to\mathbb R\dsep A,B\in\mathbb R\dsep a_n\to A\dsep b_n\to B\vdash(a+b)_n\to A+B
  }{a,b\colon\mathbb N\to\mathbb R\dsep A,B\in\mathbb R\dsep a_n\to A\dsep b_n\to B\vdash(a+b)_n\to A+B}[B21SequenceSumLimit]
    \proofstepwidestar[1]{%
      a,b\colon\mathbb N\to\mathbb R\land A,B\in\mathbb R}{%
      \rA}
    \proofstepwidestar[2]{%
      a_n\to A\land b_n\to B}{%
      \rA}
    \proofstepwidestar[3]{%
      \varepsilon\in\mathbb R\land\varepsilon>0}{%
      \rA}
    \proofstepwidestar[3]{%
      0<\varepsilon/2}{%
      \FormulaRefAuto{RealCutsCompleteOrderedField}[thm]{3}}
    \proofstepwidestar[1,2,3]{%
      \begin{aligned}[t]&\exists N\in\mathbb N\,\forall k\in\mathbb N\,(N\leq k\rightarrow{}\\[-2pt]&\qquad(|a_k-A|<\varepsilon/2\land|b_k-B|<\varepsilon/2))\end{aligned}}{%
      \FormulaRefAuto{RealSequenceConvergenceDef}[def]{1,2,4},\ \FormulaRefAuto{B21CommonNaturalBound}[thm]}
    \proofstepwidestar[6]{%
      \begin{aligned}[t]&N\in\mathbb N\land\forall k\in\mathbb N\,(N\leq k\rightarrow{}\\[-2pt]&\qquad(|a_k-A|<\varepsilon/2\land|b_k-B|<\varepsilon/2))\end{aligned}}{%
      \rA}
    \proofstepwidestar[7]{%
      n\in\mathbb N\land N\leq n}{%
      \rA}
    \proofstepwidestar[1,7]{%
      |(a+b)_n-(A+B)|\leq|a_n-A|+|b_n-B|}{%
      \FormulaRefAuto{PointwiseRealSequenceOperationsDef}[def]{1,7},\ \FormulaRefAuto{RealAbsoluteValueTriangle}[thm],\ \FormulaRefAuto{RealCutAdditiveLaws}[thm]}
    \proofstepwidestar[1,3,6,7]{%
      |(a+b)_n-(A+B)|<\varepsilon}{%
      \rRE{\rRE{\rUE{\rAEb{6}},\rAEa{7}},\rAEb{7}},\ \FormulaRefAuto{RealCutOrderArithmeticCompatibility}[thm]{8,3}}
    \proofstepwidestar[1,3,6]{%
      \exists N\in\mathbb N\,\forall n\in\mathbb N\,(N\leq n\rightarrow|(a+b)_n-(A+B)|<\varepsilon)}{%
      \rEI{\rAI{\rAEa{6},\rUI{\rRI{7,9}}}}}
    \proofstepwidestar[1,2,3]{%
      \exists N\in\mathbb N\,\forall n\in\mathbb N\,(N\leq n\rightarrow|(a+b)_n-(A+B)|<\varepsilon)}{%
      \rEE{5,6,10}}
    \proofstepwidestar[1,2]{%
      (a+b)_n\to A+B}{%
      \FormulaRefAuto{RealSequenceConvergenceDef}[def]{1,\FormulaRefAuto{PointwiseRealSequenceOperationsDef}[def],\rUI{\rRI{3,11}}}}
  \closeproofpartwideindR

  \proofpartwideindR[Differenz]{%
    a,b\colon\mathbb N\to\mathbb R\dsep A,B\in\mathbb R\dsep a_n\to A\dsep b_n\to B\vdash(a-b)_n\to A-B
  }{a,b\colon\mathbb N\to\mathbb R\dsep A,B\in\mathbb R\dsep a_n\to A\dsep b_n\to B\vdash(a-b)_n\to A-B}[B21SequenceDifferenceLimit]
    \proofstepwidestar[1]{%
      a,b\colon\mathbb N\to\mathbb R\land A,B\in\mathbb R}{%
      \rA}
    \proofstepwidestar[2]{%
      a_n\to A\land b_n\to B}{%
      \rA}
    \proofstepwidestar[3]{%
      \varepsilon\in\mathbb R\land\varepsilon>0}{%
      \rA}
    \proofstepwidestar[3]{%
      0<\varepsilon/2}{%
      \FormulaRefAuto{RealCutsCompleteOrderedField}[thm]{3}}
    \proofstepwidestar[1,2,3]{%
      \begin{aligned}[t]&\exists N\in\mathbb N\,\forall k\in\mathbb N\,(N\leq k\rightarrow{}\\[-2pt]&\qquad(|a_k-A|<\varepsilon/2\land|b_k-B|<\varepsilon/2))\end{aligned}}{%
      \FormulaRefAuto{RealSequenceConvergenceDef}[def]{1,2,4},\ \FormulaRefAuto{B21CommonNaturalBound}[thm]}
    \proofstepwidestar[6]{%
      \begin{aligned}[t]&N\in\mathbb N\land\forall k\in\mathbb N\,(N\leq k\rightarrow{}\\[-2pt]&\qquad(|a_k-A|<\varepsilon/2\land|b_k-B|<\varepsilon/2))\end{aligned}}{%
      \rA}
    \proofstepwidestar[7]{%
      n\in\mathbb N\land N\leq n}{%
      \rA}
    \proofstepwidestar[1,7]{%
      |(a-b)_n-(A-B)|\leq|a_n-A|+|b_n-B|}{%
      \FormulaRefAuto{PointwiseRealSequenceOperationsDef}[def]{1,7},\ \FormulaRefAuto{RealAbsoluteValueTriangle}[thm],\ \FormulaRefAuto{RealCutAdditiveLaws}[thm]}
    \proofstepwidestar[1,3,6,7]{%
      |(a-b)_n-(A-B)|<\varepsilon}{%
      \rRE{\rRE{\rUE{\rAEb{6}},\rAEa{7}},\rAEb{7}},\ \FormulaRefAuto{RealCutOrderArithmeticCompatibility}[thm]{8,3}}
    \proofstepwidestar[1,3,6]{%
      \exists N\in\mathbb N\,\forall n\in\mathbb N\,(N\leq n\rightarrow|(a-b)_n-(A-B)|<\varepsilon)}{%
      \rEI{\rAI{\rAEa{6},\rUI{\rRI{7,9}}}}}
    \proofstepwidestar[1,2,3]{%
      \exists N\in\mathbb N\,\forall n\in\mathbb N\,(N\leq n\rightarrow|(a-b)_n-(A-B)|<\varepsilon)}{%
      \rEE{5,6,10}}
    \proofstepwidestar[1,2]{%
      (a-b)_n\to A-B}{%
      \FormulaRefAuto{RealSequenceConvergenceDef}[def]{1,\FormulaRefAuto{PointwiseRealSequenceOperationsDef}[def],\rUI{\rRI{3,11}}}}
  \closeproofpartwideindR

  \proofpartwideindR[Skalarmultiplikation ohne Sonderfall bei Null]{%
    a\colon\mathbb N\to\mathbb R\dsep A,c\in\mathbb R\dsep a_n\to A\vdash(ca)_n\to cA
  }{a\colon\mathbb N\to\mathbb R\dsep A,c\in\mathbb R\dsep a_n\to A\vdash(ca)_n\to cA}[B21SequenceScalarLimit]
    \proofstepwidestar[1]{%
      a\colon\mathbb N\to\mathbb R\land A,c\in\mathbb R}{%
      \rA}
    \proofstepwidestar[2]{%
      a_n\to A}{%
      \rA}
    \proofstepwidestar[3]{%
      \varepsilon\in\mathbb R\land\varepsilon>0}{%
      \rA}
    \proofstepwidestar[1,3]{%
      0<\varepsilon/(|c|+1)}{%
      \FormulaRefAuto{RealAbsoluteValueNonnegativeZero}[thm]{1},\ \FormulaRefAuto{RealCutsCompleteOrderedField}[thm]{3}}
    \proofstepwidestar[1,2,3]{%
      \exists N\in\mathbb N\,\forall k\in\mathbb N\,(N\leq k\rightarrow|a_k-A|<\varepsilon/(|c|+1))}{%
      \FormulaRefAuto{RealSequenceConvergenceDef}[def]{1,2,4}}
    \proofstepwidestar[6]{%
      \begin{aligned}[t]&N\in\mathbb N\land{}\\[-2pt]&\forall k\in\mathbb N\,(N\leq k\rightarrow|a_k-A|<\varepsilon/(|c|+1))\end{aligned}}{%
      \rA}
    \proofstepwidestar[7]{%
      n\in\mathbb N\land N\leq n}{%
      \rA}
    \proofstepwidestar[1,7]{%
      |(ca)_n-cA|=|c|\,|a_n-A|}{%
      \FormulaRefAuto{PointwiseRealSequenceOperationsDef}[def]{1,7},\ \FormulaRefAuto{RealAbsoluteValueNegationProduct}[thm],\ \FormulaRefAuto{RealCutMultiplicativeLaws}[thm]}
    \proofstepwidestar[1,3,6,7]{%
      |(ca)_n-cA|\leq(|c|+1)|a_n-A|<\varepsilon}{%
      \rRE{\rRE{\rUE{\rAEb{6}},\rAEa{7}},\rAEb{7}},\ \FormulaRefAuto{RealCutsCompleteOrderedField}[thm]{8,4}}
    \proofstepwidestar[1,3,6]{%
      \exists N\in\mathbb N\,\forall n\in\mathbb N\,(N\leq n\rightarrow|(ca)_n-cA|<\varepsilon)}{%
      \rEI{\rAI{\rAEa{6},\rUI{\rRI{7,9}}}}}
    \proofstepwidestar[1,2,3]{%
      \exists N\in\mathbb N\,\forall n\in\mathbb N\,(N\leq n\rightarrow|(ca)_n-cA|<\varepsilon)}{%
      \rEE{5,6,10}}
    \proofstepwidestar[1,2]{%
      (ca)_n\to cA}{%
      \FormulaRefAuto{RealSequenceConvergenceDef}[def]{1,\FormulaRefAuto{PointwiseRealSequenceOperationsDef}[def],\rUI{\rRI{3,11}}}}
  \closeproofpartwideindR
\end{tabproofsplitwide}

\FormulaThmDeltaK[Produktregel für Grenzwerte]{%
  a,b\colon\mathbb N\to\mathbb R\dsep
  A,B\in\mathbb R\dsep
  a_n\to A\dsep b_n\to B
  \vdash (a\mathbin{\cdot}b)_n\to AB%
}{RealSequenceProductLimitRule}{%
  \DeltaRow{Mengen}{A\dsep B\dsep n}
  \DeltaRow{Funktionen}{a\dsep b\dsep a\mathbin{\cdot}b}
}
\begin{tabproofsplitwide}
  \proofpartwide{Produktabschätzung}
    \proofstepwidestar[1]{%
      a,b\colon\mathbb N\to\mathbb R\land A,B\in\mathbb R}{%
      \rA}
    \proofstepwidestar[2]{%
      a_n\to A\land b_n\to B}{%
      \rA}
    \proofstepwidestar[1,2]{%
      \operatorname{Beschr}(a)}{%
      \FormulaRefAuto{ConvergentRealSequenceBounded}[thm]{1,\rAEa{2}}}
    \proofstepwidestar[1,2]{%
      \exists C\in\mathbb R\,(0\leq C\land\forall k\in\mathbb N\,(|a_k|\leq C))}{%
      \FormulaRefAuto{BoundedRealSequenceDef}[def]{1,3}}
    \proofstepwidestar[5]{%
      C\in\mathbb R\land(0\leq C\land\forall k\in\mathbb N\,(|a_k|\leq C))}{%
      \rA}
    \proofstepwidestar[1,5]{%
      D\coloneqq C+|B|+1}{%
      \text{Abkürzung}}
    \proofstepwidestar[1,5]{%
      D\in\mathbb R\land0<D\land C\leq D\land|B|\leq D}{%
      \FormulaRefAuto{RealAbsoluteValueNonnegativeZero}[thm]{1},\ \FormulaRefAuto{RealCutsCompleteOrderedField}[thm]{5,6}}
    \proofstepwidestar[8]{%
      \varepsilon\in\mathbb R\land\varepsilon>0}{%
      \rA}
    \proofstepwidestar[1,5,8]{%
      0<\varepsilon/(2D)}{%
      \FormulaRefAuto{RealCutsCompleteOrderedField}[thm]{7,8}}
    \proofstepwidestar[1,2,5,8]{%
      \begin{aligned}[t]&\exists N\in\mathbb N\,\forall k\in\mathbb N\,(N\leq k\rightarrow{}\\[-2pt]&\quad(|a_k-A|<\varepsilon/(2D)\land|b_k-B|<\varepsilon/(2D)))\end{aligned}}{%
      \FormulaRefAuto{RealSequenceConvergenceDef}[def]{1,2,9},\ \FormulaRefAuto{B21CommonNaturalBound}[thm]}
    \proofstepwidestar[11]{%
      \begin{aligned}[t]&N\in\mathbb N\land\forall k\in\mathbb N\,(N\leq k\rightarrow{}\\[-2pt]&\quad(|a_k-A|<\varepsilon/(2D)\land|b_k-B|<\varepsilon/(2D)))\end{aligned}}{%
      \rA}
    \proofstepwidestar[12]{%
      n\in\mathbb N\land N\leq n}{%
      \rA}
    \proofstepwidestar[1,12]{%
      a_nb_n-AB=a_n(b_n-B)+B(a_n-A)}{%
      \FormulaRefAuto{RealCutMultiplicativeLaws}[thm]{1,\FormulaRefAuto{IndexedFamilyCoordinateTyping}[thm]{1,12}},\ \FormulaRefAuto{RealCutAdditiveLaws}[thm]}
    \proofstepwidestar[1,12]{%
      |a_nb_n-AB|\leq|a_n|\,|b_n-B|+|B|\,|a_n-A|}{%
      \rIE{13},\ \FormulaRefAuto{RealAbsoluteValueTriangle}[thm],\ \FormulaRefAuto{RealAbsoluteValueNegationProduct}[thm]}
    \proofstepwidestar[1,5,11,12]{%
      |a_nb_n-AB|\leq D(|b_n-B|+|a_n-A|)}{%
      \rRE{\rUE{\rAEb{\rAEb{5}}},\rAEa{12}},\ \FormulaRefAuto{RealCutOrderArithmeticCompatibility}[thm]{7,14}}
    \proofstepwidestar[1,5,8,11,12]{%
      |(a\mathbin{\cdot}b)_n-AB|<\varepsilon}{%
      \rRE{\rRE{\rUE{\rAEb{11}},\rAEa{12}},\rAEb{12}},\ \FormulaRefAuto{RealCutsCompleteOrderedField}[thm]{15,9},\ \FormulaRefAuto{PointwiseRealSequenceOperationsDef}[def]}
    \proofstepwidestar[1,5,8,11]{%
      \exists N\in\mathbb N\,\forall n\in\mathbb N\,(N\leq n\rightarrow|(a\mathbin{\cdot}b)_n-AB|<\varepsilon)}{%
      \rEI{\rAI{\rAEa{11},\rUI{\rRI{12,16}}}}}
    \proofstepwidestar[1,2,5,8]{%
      \exists N\in\mathbb N\,\forall n\in\mathbb N\,(N\leq n\rightarrow|(a\mathbin{\cdot}b)_n-AB|<\varepsilon)}{%
      \rEE{10,11,17}}
    \proofstepwidestar[1,2,5]{%
      (a\mathbin{\cdot}b)_n\to AB}{%
      \FormulaRefAuto{RealSequenceConvergenceDef}[def]{1,\FormulaRefAuto{PointwiseRealSequenceOperationsDef}[def],\rUI{\rRI{8,18}}}}
    \proofstepwidestar[1,2]{%
      (a\mathbin{\cdot}b)_n\to AB}{%
      \rEE{4,5,19}}
  \closeproofpartwide
\end{tabproofsplitwide}

\FormulaThmDeltaK[Kehrwertregel für Grenzwerte]{%
  b\colon\mathbb N\to\mathbb R\dsep
  B\in\mathbb R\dsep b_n\to B\dsep B\neq0\dsep
  \forall n\in\mathbb N\,(b_n\neq0)
  \vdash (b^{-1})_n\to\frac1B%
}{RealSequenceReciprocalLimitRule}{%
  \DeltaRow{Mengen}{B\dsep n}
  \DeltaRow{Funktionen}{b\dsep b^{-1}}
}
\begin{tabproofsplitwide}
  \proofpartwide{Nennertrennung und Fehlerabschätzung}
    \proofstepwidestar[1]{%
      b\colon\mathbb N\to\mathbb R\land B\in\mathbb R}{%
      \rA}
    \proofstepwidestar[2]{%
      b_n\to B}{%
      \rA}
    \proofstepwidestar[3]{%
      B\neq0\land\forall k\in\mathbb N\,(b_k\neq0)}{%
      \rA}
    \proofstepwidestar[4]{%
      \varepsilon\in\mathbb R\land\varepsilon>0}{%
      \rA}
    \proofstepwidestar[1,3,4]{%
      0<|B|/2\land0<\varepsilon|B|^2/2}{%
      \FormulaRefAuto{RealAbsoluteValueNonnegativeZero}[thm]{1,3},\ \FormulaRefAuto{RealCutsCompleteOrderedField}[thm]{4}}
    \proofstepwidestar[1,2,3,4]{%
      \begin{aligned}[t]&\exists N\in\mathbb N\,\forall k\in\mathbb N\,(N\leq k\rightarrow{}\\[-2pt]&\qquad(|b_k-B|<|B|/2\land|b_k-B|<\varepsilon|B|^2/2))\end{aligned}}{%
      \FormulaRefAuto{RealSequenceConvergenceDef}[def]{1,2,5},\ \FormulaRefAuto{B21CommonNaturalBound}[thm]}
    \proofstepwidestar[7]{%
      \begin{aligned}[t]&N\in\mathbb N\land\forall k\in\mathbb N\,(N\leq k\rightarrow{}\\[-2pt]&\qquad(|b_k-B|<|B|/2\land|b_k-B|<\varepsilon|B|^2/2))\end{aligned}}{%
      \rA}
    \proofstepwidestar[8]{%
      n\in\mathbb N\land N\leq n}{%
      \rA}
    \proofstepwidestar[7,8]{%
      |b_n-B|<|B|/2\land|b_n-B|<\varepsilon|B|^2/2}{%
      \rRE{\rRE{\rUE{\rAEb{7}},\rAEa{8}},\rAEb{8}}}
    \proofstepwidestar[1,8]{%
      |B|\leq|B-b_n|+|b_n|}{%
      \FormulaRefAuto{RealAbsoluteValueTriangle}[thm]{1,\FormulaRefAuto{IndexedFamilyCoordinateTyping}[thm]{1,8}},\ \FormulaRefAuto{RealCutAdditiveLaws}[thm]}
    \proofstepwidestar[1,3,7,8]{%
      0<|B|/2<|b_n|}{%
      \FormulaRefAuto{RealCutsCompleteOrderedField}[thm]{5,9,10},\ \FormulaRefAuto{RealAbsoluteValueNegationProduct}[thm]}
    \proofstepwidestar[1,3,8]{%
      \left|\frac1{b_n}-\frac1B\right|=\frac{|b_n-B|}{|b_n|\,|B|}}{%
      \FormulaRefAuto{RealCutMultiplicativeLaws}[thm]{1,3,\rRE{\rUE{\rAEb{3}},\rAEa{8}}},\ \FormulaRefAuto{RealAbsoluteValueNegationProduct}[thm]}
    \proofstepwidestar[1,3,4,7,8]{%
      \left|\frac1{b_n}-\frac1B\right|\leq\frac{2|b_n-B|}{|B|^2}<\varepsilon}{%
      \FormulaRefAuto{RealCutsCompleteOrderedField}[thm]{12,11,9}}
    \proofstepwidestar[1,3,4,7]{%
      \exists N\in\mathbb N\,\forall n\in\mathbb N\,(N\leq n\rightarrow|(b^{-1})_n-1/B|<\varepsilon)}{%
      \rEI{\rAI{\rAEa{7},\rUI{\rRI{8,13}}}},\ \FormulaRefAuto{PointwiseRealSequenceReciprocalQuotientDef}[def]}
    \proofstepwidestar[1,2,3,4]{%
      \exists N\in\mathbb N\,\forall n\in\mathbb N\,(N\leq n\rightarrow|(b^{-1})_n-1/B|<\varepsilon)}{%
      \rEE{6,7,14}}
    \proofstepwidestar[1,2,3]{%
      (b^{-1})_n\to1/B}{%
      \FormulaRefAuto{RealSequenceConvergenceDef}[def]{1,\FormulaRefAuto{PointwiseRealSequenceReciprocalQuotientDef}[def],\rUI{\rRI{4,15}}}}
  \closeproofpartwide
\end{tabproofsplitwide}

\FormulaThmDeltaK[Quotientenregel für Grenzwerte]{%
  \begin{aligned}[t]
  &a,b\colon\mathbb N\to\mathbb R\dsep A,B\in\mathbb R\dsep
    a_n\to A\dsep b_n\to B\dsep B\neq0\dsep{}\\[-2pt]
  &\quad\forall n\in\mathbb N\,(b_n\neq0)
    \vdash \left(\frac ab\right)_n\to\frac AB
  \end{aligned}%
}{RealSequenceQuotientLimitRule}{%
  \DeltaRow{Mengen}{A\dsep B\dsep n}
  \DeltaRow{Funktionen}{a\dsep b\dsep a/b}
}
\begin{tabproofsplitwide}
  \proofpartwide{Quotient als Produkt mit dem Kehrwert}
    \proofstepwidestar[1]{%
      a,b\colon\mathbb N\to\mathbb R\land A,B\in\mathbb R}{%
      \rA}
    \proofstepwidestar[2]{%
      a_n\to A\land b_n\to B}{%
      \rA}
    \proofstepwidestar[3]{%
      B\neq0\land\forall n\in\mathbb N\,(b_n\neq0)}{%
      \rA}
    \proofstepwidestar[1,3]{%
      b^{-1}\colon\mathbb N\to\mathbb R}{%
      \FormulaRefAuto{PointwiseRealSequenceReciprocalQuotientDef}[def]{1,3}}
    \proofstepwidestar[1,2,3]{%
      (b^{-1})_n\to1/B}{%
      \FormulaRefAuto{RealSequenceReciprocalLimitRule}[thm]{1,2,3}}
    \proofstepwidestar[1,2,3]{%
      (a\mathbin{\cdot}b^{-1})_n\to A(1/B)}{%
      \FormulaRefAuto{RealSequenceProductLimitRule}[thm]{1,4,2,5}}
    \proofstepwidestar[7]{%
      n\in\mathbb N}{%
      \rA}
    \proofstepwidestar[1,3,7]{%
      (a/b)_n=(a\mathbin{\cdot}b^{-1})_n}{%
      \FormulaRefAuto{PointwiseRealSequenceReciprocalQuotientDef}[def]{1,3,7},\ \FormulaRefAuto{PointwiseRealSequenceOperationsDef}[def]{1,4,7}}
    \proofstepwidestar[1,3]{%
      a/b=a\mathbin{\cdot}b^{-1}}{%
      \FormulaRefAuto{IndexedFamilyExtensionality}[thm]{1,4,\rUI{\rRI{7,8}}}}
    \proofstepwidestar[1,2,3]{%
      (a/b)_n\to A/B}{%
      \rIE{9,6},\ \FormulaRefAuto{RealSequenceConventionalNotation}[def]}
  \closeproofpartwide
\end{tabproofsplitwide}

\FormulaThmDeltaK[Stetigkeit des Betrags für Folgen]{%
  a\colon\mathbb N\to\mathbb R\dsep A\in\mathbb R\dsep a_n\to A
  \vdash (|a|)_n\to|A|%
}{AbsoluteValuePreservesSequenceLimits}{%
  \DeltaRow{Mengen}{A\dsep n}
  \DeltaRow{Funktionen}{a\dsep |a|}
}
\begin{tabproofsplitwide}
  \proofpartwideindR[Umgekehrte Dreiecksungleichung]{%
    x,y\in\mathbb R\vdash\bigl||x|-|y|\bigr|\leq|x-y|
  }{x,y\in\mathbb R\vdash\bigl||x|-|y|\bigr|\leq|x-y|}[B21ReverseTriangle]
    \proofstepwidestar[1]{%
      x,y\in\mathbb R}{%
      \rA}
    \proofstepwidestar[1]{%
      |x|\leq|x-y|+|y|}{%
      \FormulaRefAuto{RealAbsoluteValueTriangle}[thm]{1},\ \FormulaRefAuto{RealCutAdditiveLaws}[thm]}
    \proofstepwidestar[1]{%
      |y|\leq|x-y|+|x|}{%
      \FormulaRefAuto{RealAbsoluteValueTriangle}[thm]{1},\ \FormulaRefAuto{RealAbsoluteValueNegationProduct}[thm],\ \FormulaRefAuto{RealCutAdditiveLaws}[thm]}
    \proofstepwidestar[1]{%
      -|x-y|\leq|x|-|y|\leq|x-y|}{%
      \FormulaRefAuto{RealCutOrderArithmeticCompatibility}[thm]{2,3}}
    \proofstepwidestar[1]{%
      \bigl||x|-|y|\bigr|\leq|x-y|}{%
      \FormulaRefAuto{RealAbsoluteValue}[def]{4},\ \FormulaRefAuto{RealCutNegationOrder}[thm]}
  \closeproofpartwideindR

  \proofpartwide{Übertragung der Konvergenz}
    \proofstepwidestar[1]{%
      a\colon\mathbb N\to\mathbb R\land A\in\mathbb R}{%
      \rA}
    \proofstepwidestar[2]{%
      a_n\to A}{%
      \rA}
    \proofstepwidestar[3]{%
      \varepsilon\in\mathbb R\land\varepsilon>0}{%
      \rA}
    \proofstepwidestar[1,2,3]{%
      \exists N\in\mathbb N\,\forall k\in\mathbb N\,(N\leq k\rightarrow|a_k-A|<\varepsilon)}{%
      \FormulaRefAuto{RealSequenceConvergenceDef}[def]{1,2,3}}
    \proofstepwidestar[5]{%
      N\in\mathbb N\land\forall k\in\mathbb N\,(N\leq k\rightarrow|a_k-A|<\varepsilon)}{%
      \rA}
    \proofstepwidestar[6]{%
      n\in\mathbb N\land N\leq n}{%
      \rA}
    \proofstepwidestar[1,5,6]{%
      \bigl|(|a|)_n-|A|\bigr|\leq|a_n-A|<\varepsilon}{%
      \FormulaRefAuto{B21ReverseTriangle}[thm]{1,\FormulaRefAuto{IndexedFamilyCoordinateTyping}[thm]{1,6}},\ \FormulaRefAuto{PointwiseRealSequenceOperationsDef}[def],\ \rRE{\rRE{\rUE{\rAEb{5}},\rAEa{6}},\rAEb{6}}}
    \proofstepwidestar[1,5]{%
      \exists N\in\mathbb N\,\forall n\in\mathbb N\,(N\leq n\rightarrow\bigl|(|a|)_n-|A|\bigr|<\varepsilon)}{%
      \rEI{\rAI{\rAEa{5},\rUI{\rRI{6,7}}}}}
    \proofstepwidestar[1,2,3]{%
      \exists N\in\mathbb N\,\forall n\in\mathbb N\,(N\leq n\rightarrow\bigl|(|a|)_n-|A|\bigr|<\varepsilon)}{%
      \rEE{4,5,8}}
    \proofstepwidestar[1,2]{%
      (|a|)_n\to|A|}{%
      \FormulaRefAuto{RealSequenceConvergenceDef}[def]{1,\FormulaRefAuto{PointwiseRealSequenceOperationsDef}[def],\rUI{\rRI{3,9}}}}
  \closeproofpartwide
\end{tabproofsplitwide}

\section{Cauchy-Folgen}

\FormulaDefDeltaK[Cauchy-Folge in den reellen Zahlen]{%
  \begin{aligned}[t]
  &a\colon\mathbb N\to\mathbb R\vdash
    \operatorname{Cauchy}(a)\coloneqq{}\\[-2pt]
  &\quad\forall\varepsilon\in\mathbb R\,
    \Bigl(\varepsilon>0\rightarrow
      \exists N\in\mathbb N\,
      \forall m,n\in\mathbb N\,\bigl((N\leq m\land N\leq n)\\[-2pt]
  &\hspace{54mm}\rightarrow|a_m-a_n|<\varepsilon\bigr)\Bigr)
  \end{aligned}%
}{RealCauchySequenceDef}{%
  \DeltaRow{Mengen}{\varepsilon\dsep N\dsep m\dsep n}
  \DeltaRow{Funktionen}{a}
  \DeltaRow{Neues Prädikat}{\operatorname{Cauchy}}
}

Die Cauchy-Bedingung erwähnt keinen vermuteten Grenzwert. Sie verlangt, dass
späte Folgenglieder untereinander beliebig nahe liegen. Ob daraus ein
Grenzwert folgt, hängt von der Vollständigkeit der Zielmenge ab.

\FormulaThmDeltaK[Konvergenz impliziert die Cauchy-Eigenschaft]{%
  a\colon\mathbb N\to\mathbb R\dsep L\in\mathbb R\dsep a_n\to L
  \vdash\operatorname{Cauchy}(a)%
}{ConvergentSequenceIsCauchy}{%
  \DeltaRow{Mengen}{L\dsep\varepsilon\dsep N\dsep m\dsep n}
  \DeltaRow{Funktionen}{a}
}
\begin{tabproofsplitwide}
  \proofpartwide{Zwei späte Glieder}
    \proofstepwidestar[1]{%
      a\colon\mathbb N\to\mathbb R\land L\in\mathbb R}{%
      \rA}
    \proofstepwidestar[2]{%
      a_n\to L}{%
      \rA}
    \proofstepwidestar[3]{%
      \varepsilon\in\mathbb R\land\varepsilon>0}{%
      \rA}
    \proofstepwidestar[3]{%
      0<\varepsilon/2}{%
      \FormulaRefAuto{RealCutsCompleteOrderedField}[thm]{3}}
    \proofstepwidestar[1,2,3]{%
      \exists N\in\mathbb N\,\forall k\in\mathbb N\,(N\leq k\rightarrow|a_k-L|<\varepsilon/2)}{%
      \FormulaRefAuto{RealSequenceConvergenceDef}[def]{1,2,4}}
    \proofstepwidestar[6]{%
      N\in\mathbb N\land\forall k\in\mathbb N\,(N\leq k\rightarrow|a_k-L|<\varepsilon/2)}{%
      \rA}
    \proofstepwidestar[7]{%
      (m\in\mathbb N\land n\in\mathbb N)\land(N\leq m\land N\leq n)}{%
      \rA}
    \proofstepwidestar[1,3,6,7]{%
      |a_m-a_n|\leq|a_m-L|+|a_n-L|<\varepsilon}{%
      \FormulaRefAuto{RealDistanceTriangle}[thm]{1,\FormulaRefAuto{IndexedFamilyCoordinateTyping}[thm]{1,7}},\ \FormulaRefAuto{RealDistanceBasicLaws}[thm],\ \rUE{\rAEb{6}},\ 7,\ \FormulaRefAuto{RealCutOrderArithmeticCompatibility}[thm]}
    \proofstepwidestar[1,3,6]{%
      \begin{aligned}[t]&\exists N\in\mathbb N\,\forall m,n\in\mathbb N\,\\[-2pt]&\qquad((N\leq m\land N\leq n)\rightarrow|a_m-a_n|<\varepsilon)\end{aligned}}{%
      \rEI{\rAI{\rAEa{6},\rUI{\rUI{\rRI{7,8}}}}}}
    \proofstepwidestar[1,2,3]{%
      \begin{aligned}[t]&\exists N\in\mathbb N\,\forall m,n\in\mathbb N\,\\[-2pt]&\qquad((N\leq m\land N\leq n)\rightarrow|a_m-a_n|<\varepsilon)\end{aligned}}{%
      \rEE{5,6,9}}
    \proofstepwidestar[1,2]{%
      \operatorname{Cauchy}(a)}{%
      \FormulaRefAuto{RealCauchySequenceDef}[def]{1,\rUI{\rRI{3,10}}}}
  \closeproofpartwide
\end{tabproofsplitwide}

\FormulaThmDeltaK[Cauchy-Folgen sind beschränkt]{%
  a\colon\mathbb N\to\mathbb R\dsep\operatorname{Cauchy}(a)
  \vdash\operatorname{Beschr}(a)%
}{RealCauchySequenceBounded}{%
  \DeltaRow{Mengen}{N\dsep n}
  \DeltaRow{Funktionen}{a}
}
\begin{tabproofsplitwide}
  \proofpartwide{Restschranke um ein festes Folgenglied}
    \proofstepwidestar[1]{%
      a\colon\mathbb N\to\mathbb R}{%
      \rA}
    \proofstepwidestar[2]{%
      \operatorname{Cauchy}(a)}{%
      \rA}
    \proofstepwidestar[1,2]{%
      \begin{aligned}[t]&\exists N\in\mathbb N\,\forall m,n\in\mathbb N\,\\[-2pt]&\qquad((N\leq m\land N\leq n)\rightarrow|a_m-a_n|<1)\end{aligned}}{%
      \FormulaRefAuto{RealCauchySequenceDef}[def]{1,2,\FormulaRefAuto{RealCutsCompleteOrderedField}[thm]}}
    \proofstepwidestar[4]{%
      \begin{aligned}[t]&N\in\mathbb N\land\forall m,n\in\mathbb N\,\\[-2pt]&\qquad((N\leq m\land N\leq n)\rightarrow|a_m-a_n|<1)\end{aligned}}{%
      \rA}
    \proofstepwidestar[5]{%
      n\in\mathbb N\land N\leq n}{%
      \rA}
    \proofstepwidestar[4,5]{%
      |a_n-a_N|<1}{%
      \rUE{\rUE{\rAEb{4}}},\ 5,\ \FormulaRefAuto{PeanoLeqReflexive}[thm]{\rAEa{4}}}
    \proofstepwidestar[1,4,5]{%
      |a_n|\leq|a_n-a_N|+|a_N|<1+|a_N|}{%
      \FormulaRefAuto{RealAbsoluteValueTriangle}[thm]{\FormulaRefAuto{IndexedFamilyCoordinateTyping}[thm]{1,4},\FormulaRefAuto{IndexedFamilyCoordinateTyping}[thm]{1,5}},\ \FormulaRefAuto{RealCutOrderArithmeticCompatibility}[thm]{6}}
    \proofstepwidestar[1,4]{%
      \forall n\in\mathbb N\,(N\leq n\rightarrow|a_n|\leq1+|a_N|)}{%
      \rUI{\rRI{5,7}},\ \FormulaRefAuto{RealCutTotalOrder}[thm]}
    \proofstepwidestar[1,4]{%
      \operatorname{Beschr}(a)}{%
      \FormulaRefAuto{B21EventuallyBoundedIsBounded}[thm]{1,4,\FormulaRefAuto{RealAbsoluteValueNonnegativeZero}[thm]{\FormulaRefAuto{IndexedFamilyCoordinateTyping}[thm]{1,4}},\FormulaRefAuto{RealCutAdditiveClosure}[thm],8}}
    \proofstepwidestar[1,2]{%
      \operatorname{Beschr}(a)}{%
      \rEE{3,4,9}}
  \closeproofpartwide
\end{tabproofsplitwide}

\FormulaThmDeltaK[Cauchy-Kriterium in den reellen Zahlen]{%
  a\colon\mathbb N\to\mathbb R\vdash
  \operatorname{Konv}(a)\leftrightarrow\operatorname{Cauchy}(a)%
}{RealSequenceCauchyCriterion}{%
  \DeltaRow{Funktionen}{a}
}
\begin{tabproofsplitwide}
  \proofpartwideindR[Hinrichtung]{%
    a\colon\mathbb N\to\mathbb R\dsep\operatorname{Konv}(a)\vdash\operatorname{Cauchy}(a)
  }{a\colon\mathbb N\to\mathbb R\dsep\operatorname{Konv}(a)\vdash\operatorname{Cauchy}(a)}[B21CauchyCriterionForward]
    \proofstepwidestar[1]{%
      a\colon\mathbb N\to\mathbb R}{%
      \rA}
    \proofstepwidestar[2]{%
      \operatorname{Konv}(a)}{%
      \rA}
    \proofstepwidestar[1,2]{%
      \exists L\in\mathbb R\,(a_n\to L)}{%
      \FormulaRefAuto{RealSequenceConvergentDivergentDef}[def]{1,2}}
    \proofstepwidestar[4]{%
      L\in\mathbb R\land a_n\to L}{%
      \rA}
    \proofstepwidestar[1,4]{%
      \operatorname{Cauchy}(a)}{%
      \FormulaRefAuto{ConvergentSequenceIsCauchy}[thm]{1,4}}
    \proofstepwidestar[1,2]{%
      \operatorname{Cauchy}(a)}{%
      \rEE{3,4,5}}
  \closeproofpartwideindR

  \proofpartwideindR[Rückrichtung: Infima der Restwertemengen]{%
    a\colon\mathbb N\to\mathbb R\dsep\operatorname{Cauchy}(a)\vdash\operatorname{Konv}(a)
  }{a\colon\mathbb N\to\mathbb R\dsep\operatorname{Cauchy}(a)\vdash\operatorname{Konv}(a)}[B21CauchyCriterionBackward]
    \proofstepwidestar[1]{%
      a\colon\mathbb N\to\mathbb R}{%
      \rA}
    \proofstepwidestar[2]{%
      \operatorname{Cauchy}(a)}{%
      \rA}
    \proofstepwidestar[1,2]{%
      \operatorname{Beschr}(a)}{%
      \FormulaRefAuto{RealCauchySequenceBounded}[thm]{1,2}}
    \proofstepwidestar[1,2]{%
      \exists C\in\mathbb R\,(0\leq C\land\forall k\in\mathbb N\,(|a_k|\leq C))}{%
      \FormulaRefAuto{BoundedRealSequenceDef}[def]{1,3}}
    \proofstepwidestar[5]{%
      C\in\mathbb R\land(0\leq C\land\forall k\in\mathbb N\,(|a_k|\leq C))}{%
      \rA}
    \proofstepwidestar[1]{%
      T_N\coloneqq a[\{k\in\mathbb N\mid N\leq k\}]\quad(N\in\mathbb N)}{%
      \text{Abkürzung}}
    \proofstepwidestar[1,5]{%
      \begin{aligned}[t]&\forall N\in\mathbb N\,\bigl(a_N\in T_N\subseteq\mathbb R\land{}\\[-2pt]&\qquad\forall x\in T_N\,(-C\leq x\land x\leq C)\bigr)\end{aligned}}{%
      6,\ \FormulaRefAuto{PeanoLeqReflexive}[thm],\ \FormulaRefAuto{IndexedFamilyImageMembership}[thm]{1},\ \FormulaRefAuto{RealAbsoluteValue}[def]{5}}
    \proofstepwidestar[1,5]{%
      \begin{aligned}[t]&\exists\alpha\colon\mathbb N\to\mathbb R\,\forall N\in\mathbb N\,\\[-2pt]&\qquad\alpha_N=\inf_{\mathbb R,\RealLe}(T_N)\end{aligned}}{%
      \FormulaRefAuto{RealInfimumCharacterization}[thm]{7},\ \FormulaRefAuto{TypedFunctionDefinitionPrinciple}[thm],\ \FormulaRefAuto{IndexedFamilyDef}[def]}
    \proofstepwidestar[9]{%
      \begin{aligned}[t]&\alpha\colon\mathbb N\to\mathbb R\land\forall N\in\mathbb N\,\\[-2pt]&\qquad\alpha_N=\inf_{\mathbb R,\RealLe}(T_N)\end{aligned}}{%
      \rA}
    \proofstepwidestar[1,5,9]{%
      \forall N\in\mathbb N\,(-C\leq\alpha_N\land\alpha_N\leq a_N\land a_N\leq C)}{%
      \FormulaRefAuto{RealInfimumCharacterization}[thm]{7,9}}
    \proofstepwidestar[1,5,9]{%
      \alpha_0\in\alpha[\mathbb N]\subseteq\mathbb R}{%
      \FormulaRefAuto{IndexedFamilyImageMembership}[thm]{9},\ \FormulaRefAuto{IndexedFamilyCoordinateTyping}[thm]{9,\FormulaRefAuto{PeanoZero}[ax]},\ \FormulaRefAuto{IndexedFamilyImageSubsetCodomain}[thm]{9}}
    \proofstepwidestar[1,5,9]{%
      \begin{aligned}[t]&\exists L\in\mathbb R\,\bigl(\forall k\in\mathbb N\,(\alpha_k\leq L)\land{}\\[-2pt]&\forall u\in\mathbb R\,((\forall k\in\mathbb N\,(\alpha_k\leq u))\rightarrow L\leq u)\bigr)\end{aligned}}{%
      \FormulaRefAuto{RealLeastUpperBoundProperty}[thm]{11,10},\ \FormulaRefAuto{IndexedFamilyImageMembership}[thm]{9}}
    \proofstepwidestar[13]{%
      \begin{aligned}[t]&L\in\mathbb R\land\bigl(\forall k\in\mathbb N\,(\alpha_k\leq L)\land{}\\[-2pt]&\forall u\in\mathbb R\,((\forall k\in\mathbb N\,(\alpha_k\leq u))\rightarrow L\leq u)\bigr)\end{aligned}}{%
      \rA}
    \proofstepwidestar[14]{%
      \varepsilon\in\mathbb R\land\varepsilon>0}{%
      \rA}
    \proofstepwidestar[14]{%
      0<\varepsilon/2<\varepsilon}{%
      \FormulaRefAuto{RealCutsCompleteOrderedField}[thm]{14}}
    \proofstepwidestar[1,2,14]{%
      \begin{aligned}[t]&\exists N\in\mathbb N\,\forall m,n\in\mathbb N\,\\[-2pt]&\qquad((N\leq m\land N\leq n)\rightarrow|a_m-a_n|<\varepsilon/2)\end{aligned}}{%
      \FormulaRefAuto{RealCauchySequenceDef}[def]{1,2,15}}
    \proofstepwidestar[17]{%
      \begin{aligned}[t]&N\in\mathbb N\land\forall m,n\in\mathbb N\,\\[-2pt]&\qquad((N\leq m\land N\leq n)\rightarrow|a_m-a_n|<\varepsilon/2)\end{aligned}}{%
      \rA}
    \proofstepwidestar[18]{%
      n\in\mathbb N\land N\leq n}{%
      \rA}
    \proofstepwidestar[19]{%
      x\in T_N}{%
      \rA}
    \proofstepwidestar[1,17,19]{%
      \exists m\in\mathbb N\,(N\leq m\land x=a_m)}{%
      6,\ \FormulaRefAuto{IndexedFamilyImageMembership}[thm]{1,19}}
    \proofstepwidestar[21]{%
      m\in\mathbb N\land(N\leq m\land x=a_m)}{%
      \rA}
    \proofstepwidestar[17,18,21]{%
      |a_m-a_n|<\varepsilon/2}{%
      \rUE{\rUE{\rAEb{17}}},\ 18,\ 21}
    \proofstepwidestar[1,14,17,18,21]{%
      a_n-\varepsilon/2\leq x}{%
      \FormulaRefAuto{RealEpsilonNeighborhoodInterval}[thm]{\FormulaRefAuto{IndexedFamilyCoordinateTyping}[thm]{1,18},15,22},\ \rIE{\rAEb{\rAEb{21}}},\ \FormulaRefAuto{RealCutTotalOrder}[thm]}
    \proofstepwidestar[1,14,17,18,19]{%
      a_n-\varepsilon/2\leq x}{%
      \rEE{20,21,23}}
    \proofstepwidestar[1,14,17,18]{%
      \forall x\in T_N\,(a_n-\varepsilon/2\leq x)}{%
      \rUI{\rRI{19,24}}}
    \proofstepwidestar[1,5,9,14,17,18]{%
      a_n-\varepsilon/2\leq\alpha_N}{%
      \FormulaRefAuto{RealInfimumCharacterization}[thm]{7,9,25}}
    \proofstepwidestar[1,5,9,13,14,17,18]{%
      a_n-\varepsilon/2\leq L}{%
      \rRE{\rUE{\rAEa{\rAEb{13}}},\rAEa{17}},\ \FormulaRefAuto{RealCutTotalOrder}[thm]{26}}
    \proofstepwidestar[28]{%
      k\in\mathbb N}{%
      \rA}
    \proofstepwidestar[17,28]{%
      N+k\in\mathbb N\land(N\leq N+k\land k\leq N+k)}{%
      \FormulaRefAuto{B21CommonNaturalBound}[thm]{\rAEa{17},28}}
    \proofstepwidestar[1,17,28]{%
      a_{N+k}\in T_k}{%
      6,\ \FormulaRefAuto{IndexedFamilyImageMembership}[thm]{1},\ 29}
    \proofstepwidestar[1,5,9,17,28]{%
      \alpha_k\leq a_{N+k}}{%
      \FormulaRefAuto{RealInfimumCharacterization}[thm]{7,9,30}}
    \proofstepwidestar[17,18,28]{%
      |a_{N+k}-a_n|<\varepsilon/2}{%
      \rUE{\rUE{\rAEb{17}}},\ 18,\ 29}
    \proofstepwidestar[1,5,9,14,17,18,28]{%
      \alpha_k\leq a_n+\varepsilon/2}{%
      \FormulaRefAuto{RealEpsilonNeighborhoodInterval}[thm]{1,15,32},\ \FormulaRefAuto{RealCutTotalOrder}[thm]{31}}
    \proofstepwidestar[1,5,9,14,17,18]{%
      \forall k\in\mathbb N\,(\alpha_k\leq a_n+\varepsilon/2)}{%
      \rUI{\rRI{28,33}}}
    \proofstepwidestar[1,5,9,13,14,17,18]{%
      L\leq a_n+\varepsilon/2}{%
      \rUE{\rAEb{\rAEb{13}}},\ \FormulaRefAuto{RealCutAdditiveClosure}[thm]{1,14},\ 34}
    \proofstepwidestar[1,5,9,13,14,17,18]{%
      |a_n-L|\leq\varepsilon/2<\varepsilon}{%
      \FormulaRefAuto{RealAbsoluteValue}[def],\ \FormulaRefAuto{RealCutNegationOrder}[thm],\ \FormulaRefAuto{RealCutOrderArithmeticCompatibility}[thm]{27,35,15}}
    \proofstepwidestar[1,5,9,13,14,17]{%
      \exists N\in\mathbb N\,\forall n\in\mathbb N\,(N\leq n\rightarrow|a_n-L|<\varepsilon)}{%
      \rEI{\rAI{\rAEa{17},\rUI{\rRI{18,36}}}}}
    \proofstepwidestar[1,2,5,9,13,14]{%
      \exists N\in\mathbb N\,\forall n\in\mathbb N\,(N\leq n\rightarrow|a_n-L|<\varepsilon)}{%
      \rEE{16,17,37}}
    \proofstepwidestar[1,2,5,9,13]{%
      a_n\to L}{%
      \FormulaRefAuto{RealSequenceConvergenceDef}[def]{1,\rAEa{13},\rUI{\rRI{14,38}}}}
    \proofstepwidestar[1,2,5,9,13]{%
      \operatorname{Konv}(a)}{%
      \FormulaRefAuto{RealSequenceConvergentDivergentDef}[def]{1,\rEI{\rAI{\rAEa{13},39}}}}
    \proofstepwidestar[1,2,5,9]{%
      \operatorname{Konv}(a)}{%
      \rEE{12,13,40}}
    \proofstepwidestar[1,2,5]{%
      \operatorname{Konv}(a)}{%
      \rEE{8,9,41}}
    \proofstepwidestar[1,2]{%
      \operatorname{Konv}(a)}{%
      \rEE{4,5,42}}
  \closeproofpartwideindR

  \proofpartwide{Äquivalenz}
    \proofstepwidestar[1]{%
      a\colon\mathbb N\to\mathbb R}{%
      \rA}
    \proofstepwidestar[1]{%
      \operatorname{Konv}(a)\leftrightarrow\operatorname{Cauchy}(a)}{%
      \rLRI{\FormulaRefAuto{B21CauchyCriterionForward}[thm]{1},\FormulaRefAuto{B21CauchyCriterionBackward}[thm]{1}}}
  \closeproofpartwide
\end{tabproofsplitwide}

Die Rückrichtung des Cauchy-Kriteriums ist genau die Stelle, an der die
Vollständigkeit von \(\mathbb R\) benötigt wird. In \(\mathbb Q\) gibt es
Cauchy-Folgen, deren reeller Grenzwert irrational ist und die deshalb in
\(\mathbb Q\) keinen Grenzwert besitzen.

Die Definitionen dieses Kapitels trennen bereits zwei Arten von Argumenten.
Monotonie, Suprema und die Grenzwertregeln verwenden Ordnung und Arithmetik
von \(\mathbb R\). Dagegen benötigen die Eindeutigkeit des Grenzwerts und die
Hinrichtung des Cauchy-Kriteriums nur die Grundgesetze des reellen Abstands.
Band~43 erhebt genau diese Gesetze zu den Axiomen einer Metrik und formuliert
die abstandstheoretischen Aussagen in dieser Allgemeinheit neu.

\chapter{Die Mogiljanskaja-Familien}

\section{Die beiden Grundfamilien}

Für die spätere Konstruktion zweier unendlicher Halbgruppen werden zwei
disjunkte getaggte Schichten benötigt. In diesem Band geht es ausschließlich
um ihre familien- und mengentheoretischen Eigenschaften; Operationen auf den
späteren Trägermengen werden erst im Halbgruppenband definiert.
Die Zeichen \(0\) und \(1\) in den folgenden geordneten Paaren sind dabei
ausdrücklich die natürlichen Peano-Zahlen und dienen lediglich als Tags.

\FormulaDefDeltaK[Die einfach indizierte Grundschicht]{%
  L\coloneqq\{0\}\times\mathbb N%
}{MogiljanskajaLayerLDef}{%
  \DeltaRow{Mengen}{L}
}

\FormulaDefDeltaK[Die zweifach indizierte Grundschicht]{%
  D\coloneqq\{1\}\times(\mathbb N\times\mathbb N)%
}{MogiljanskajaLayerDDef}{%
  \DeltaRow{Mengen}{D}
}

\FormulaDefDeltaK[Die einfach indizierte Grundfamilie]{%
  \begin{aligned}[t]
    &a\colon\mathbb N\to L,\\[-2pt]
    &\forall i\in\mathbb N\,\bigl(a_i\coloneqq(0,i)\bigr)
  \end{aligned}%
}{MogiljanskajaAElementsDef}{%
  \DeltaRow{Mengen}{L\dsep a_i\dsep i}
  \DeltaRow{Funktionen}{a}
}

\begin{tabproofsplitwide}
  \proofpartwide{Typisierung der Grundfamilie}
    \proofstepwidestar[1]{%
      i\in\mathbb N}{%
      \rA}
    \proofstepwidestar[1]{%
      (0,i)\in\{0\}\times\mathbb N}{%
      \text{Def. }\ensuremath{\times},\ \text{Def. }\ensuremath{\{0\}},\ 1}
    \proofstepwidestar[]{%
      L=\{0\}\times\mathbb N}{%
      \FormulaRefAuto{MogiljanskajaLayerLDef}[def]}
    \proofstepwidestar[1]{%
      (0,i)\in L}{%
      \rIE{3,2}}
    \proofstepwidestar[]{%
      \forall i\in\mathbb N\,((0,i)\in L)}{%
      \rUI{\rRI{1,4}}}
  \closeproofpartwide
\end{tabproofsplitwide}

\FormulaDefDeltaK[Die zweifach indizierte Grundfamilie]{%
  \begin{aligned}[t]
    &b\colon\mathbb N\times\mathbb N\to D,\\[-2pt]
    &\forall i,j\in\mathbb N\,\bigl(b_{ij}\coloneqq(1,(i,j))\bigr)
  \end{aligned}%
}{MogiljanskajaBElementsDef}{%
  \DeltaRow{Mengen}{D\dsep b_{ij}\dsep i\dsep j}
  \DeltaRow{Funktionen}{b}
}

Hier steht \(b_{ij}\) für \(b((i,j))\). Die Familie ist also eine einzige
Funktion auf \(\mathbb N\times\mathbb N\), keine zweistufige Sammlung
voneinander unabhängiger Funktionen.

\FormulaThmDeltaK[Eindeutigkeit der (a)-Indizes]{%
  i,j\in\mathbb N\dsep a_i=a_j\vdash i=j%
}{MogiljanskajaAIndexUnique}{%
  \DeltaRow{Mengen}{a_i\dsep a_j\dsep i\dsep j}
  \DeltaRow{Funktionen}{a}
}
\begin{tabproofsplitwide}
  \proofpartwide{Gleichheit der zweiten Koordinate}
    \proofstepwidestar[1]{%
      i,j\in\mathbb N}{%
      \rA}
    \proofstepwidestar[2]{%
      a_i=a_j}{%
      \rA}
    \proofstepwidestar[1,2]{%
      (0,i)=(0,j)}{%
      \FormulaRefAuto{MogiljanskajaAElementsDef}[def]{1,2}}
    \proofstepwidestar[1,2]{%
      i=j}{%
      \FormulaRefAuto{TaggedPairRightInjective}[thm]{3}}
  \closeproofpartwide
\end{tabproofsplitwide}

\FormulaThmDeltaK[Eindeutigkeit der (b)-Indizes]{%
  i,j,m,n\in\mathbb N\dsep b_{ij}=b_{mn}
  \vdash i=m\land j=n%
}{MogiljanskajaBIndexUnique}{%
  \DeltaRow{Mengen}{b_{ij}\dsep b_{mn}\dsep i\dsep j\dsep m\dsep n}
  \DeltaRow{Funktionen}{b}
}
\begin{tabproofsplitwide}
  \proofpartwide{Gleichheit der beiden inneren Koordinaten}
    \proofstepwidestar[1]{%
      i,j,m,n\in\mathbb N}{%
      \rA}
    \proofstepwidestar[2]{%
      b_{ij}=b_{mn}}{%
      \rA}
    \proofstepwidestar[1,2]{%
      (1,(i,j))=(1,(m,n))}{%
      \FormulaRefAuto{MogiljanskajaBElementsDef}[def]{1,2}}
    \proofstepwidestar[1,2]{%
      i=m\land j=n}{%
      \FormulaRefAuto{TaggedNestedPairInjective}[thm]{3}}
  \closeproofpartwide
\end{tabproofsplitwide}

\FormulaThmDeltaK[Elementkriterium der (a)-Schicht]{%
  s\in L\leftrightarrow\exists i\in\mathbb N\,(s=a_i)%
}{MogiljanskajaLayerLMembership}{%
  \DeltaRow{Mengen}{L\dsep s\dsep a_i\dsep i}
  \DeltaRow{Funktionen}{a}
}
\begin{tabproofsplitwide}
  \proofpartwide{Darstellung der getaggten Schicht}
    \proofstepwidestar[]{%
      L=\{0\}\times\mathbb N}{%
      \FormulaRefAuto{MogiljanskajaLayerLDef}[def]}
    \proofstepwidestar[]{%
      s\in L\leftrightarrow\exists i\in\mathbb N\,(s=(0,i))}{%
      \FormulaRefAuto{TaggedProductLayerRepresentation}[thm],\ \rIE{1}}
    \proofstepwidestar[]{%
      \forall i\in\mathbb N\,(a_i=(0,i))}{%
      \FormulaRefAuto{MogiljanskajaAElementsDef}[def]}
    \proofstepwidestar[]{%
      s\in L\leftrightarrow\exists i\in\mathbb N\,(s=a_i)}{%
      2,\ \rUE{3}}
  \closeproofpartwide
\end{tabproofsplitwide}

\FormulaThmDeltaK[Elementkriterium der (b)-Schicht]{%
  s\in D\leftrightarrow
    \exists i,j\in\mathbb N\,(s=b_{ij})%
}{MogiljanskajaLayerDMembership}{%
  \DeltaRow{Mengen}{D\dsep s\dsep b_{ij}\dsep i\dsep j}
  \DeltaRow{Funktionen}{b}
}
\begin{tabproofsplitwide}
  \proofpartwide{Darstellung der doppelt indizierten Schicht}
    \proofstepwidestar[]{%
      D=\{1\}\times(\mathbb N\times\mathbb N)}{%
      \FormulaRefAuto{MogiljanskajaLayerDDef}[def]}
    \proofstepwidestar[]{%
      s\in D\leftrightarrow\exists i,j\in\mathbb N\,(s=(1,(i,j)))}{%
      \FormulaRefAuto{TaggedNestedProductLayerRepresentation}[thm],\ \rIE{1}}
    \proofstepwidestar[]{%
      \forall i,j\in\mathbb N\,(b_{ij}=(1,(i,j)))}{%
      \FormulaRefAuto{MogiljanskajaBElementsDef}[def]}
    \proofstepwidestar[]{%
      s\in D\leftrightarrow\exists i,j\in\mathbb N\,(s=b_{ij})}{%
      2,\ \rUE{\rUE{3}}}
  \closeproofpartwide
\end{tabproofsplitwide}

\FormulaThmDeltaK[Bildmenge der (a)-Familie]{%
  a[\mathbb N]=L%
}{MogiljanskajaLayerLImage}{%
  \DeltaRow{Mengen}{L\dsep a[\mathbb N]}
  \DeltaRow{Funktionen}{a}
}
\begin{tabproofsplitwide}
  \proofpartwide{Mengenextensionalität}
    \proofstepwidestar[]{%
      a\colon\mathbb N\to L}{%
      \FormulaRefAuto{MogiljanskajaAElementsDef}[def]}
    \proofstepwidestar[]{%
      s\in a[\mathbb N]\leftrightarrow\exists i\in\mathbb N\,(s=a_i)}{%
      \FormulaRefAuto{IndexedFamilyImageMembership}[thm]{1},\ \FormulaRefAuto{IndexedFamilyImageDef}[def]}
    \proofstepwidestar[]{%
      s\in L\leftrightarrow\exists i\in\mathbb N\,(s=a_i)}{%
      \FormulaRefAuto{MogiljanskajaLayerLMembership}[thm]}
    \proofstepwidestar[]{%
      a[\mathbb N]=L}{%
      \text{Extensionalität}(2,3)}
  \closeproofpartwide
\end{tabproofsplitwide}

\FormulaThmDeltaK[Bildmenge der (b)-Familie]{%
  b[\mathbb N\times\mathbb N]=D%
}{MogiljanskajaLayerDImage}{%
  \DeltaRow{Mengen}{D\dsep b[\mathbb N\times\mathbb N]}
  \DeltaRow{Funktionen}{b}
}
\begin{tabproofsplitwide}
  \proofpartwide{Mengenextensionalität}
    \proofstepwidestar[]{%
      b\colon\mathbb N\times\mathbb N\to D}{%
      \FormulaRefAuto{MogiljanskajaBElementsDef}[def]}
    \proofstepwidestar[]{%
      s\in b[\mathbb N\times\mathbb N]\leftrightarrow\exists i,j\in\mathbb N\,(s=b_{ij})}{%
      \FormulaRefAuto{IndexedFamilyImageMembership}[thm]{1},\ \FormulaRefAuto{IndexedFamilyImageDef}[def],\ \text{Def. }\ensuremath{\times}}
    \proofstepwidestar[]{%
      s\in D\leftrightarrow\exists i,j\in\mathbb N\,(s=b_{ij})}{%
      \FormulaRefAuto{MogiljanskajaLayerDMembership}[thm]}
    \proofstepwidestar[]{%
      b[\mathbb N\times\mathbb N]=D}{%
      \text{Extensionalität}(2,3)}
  \closeproofpartwide
\end{tabproofsplitwide}

\FormulaThmDeltaK[Die beiden Grundschichten sind disjunkt]{%
  L\cap D=\varnothing%
}{MogiljanskajaLayersDisjoint}{%
  \DeltaRow{Mengen}{L\dsep D}
}
\begin{tabproofsplitwide}
  \proofpartwide{Verschiedene Tags}
    \proofstepwidestar[]{%
      0\neq1}{%
      \FormulaRefAuto{PeanoZeroNeOne}[thm]}
    \proofstepwidestar[]{%
      (\{0\}\times\mathbb N)\cap(\{1\}\times(\mathbb N\times\mathbb N))=\varnothing}{%
      \FormulaRefAuto{TaggedProductLayersDisjoint}[thm]{1}}
    \proofstepwidestar[]{%
      L\cap D=\varnothing}{%
      \FormulaRefAuto{MogiljanskajaLayerLDef}[def],\ \FormulaRefAuto{MogiljanskajaLayerDDef}[def],\ \rIE{2}}
  \closeproofpartwide
\end{tabproofsplitwide}

\section{Marker und Reserveschichten}

\FormulaDefDeltaK[Marker und unendliche Reserveschicht]{%
  \begin{aligned}[t]
    c_0&\coloneqq b_{00},&
    c_1&\coloneqq b_{11},&
    c_2&\coloneqq b_{10},\\[-2pt]
    U&\coloneqq\bigl\{u\in D\bigm|
      \exists n,j\in\mathbb N\,
        u=b_{\Succ(\Succ(n)),j}\bigr\},\\[-2pt]
    V&\coloneqq U\cup\{c_2\},&
    P&\coloneqq\{c_0,c_1\},&
    D'&\coloneqq P\cup V
  \end{aligned}%
}{MogiljanskajaDPrimeDef}{%
  \DeltaRow{Mengen}{D\dsep D'\dsep U\dsep V\dsep P\dsep
    c_0\dsep c_1\dsep c_2\dsep b_{ij}\dsep n\dsep j\dsep u}
  \DeltaRow{Funktionen}{b}
}

Somit besteht \(U\) genau aus den Elementen der Zeilen mit erstem Index
mindestens \(2\). Die Menge \(V\) adjungiert den Marker \(c_2=b_{10}\),
und \(D'\) adjungiert zusätzlich \(c_0=b_{00}\) und \(c_1=b_{11}\).

\FormulaThmDeltaK[Elementkriterium der Reserveschicht]{%
  u\in U\eqvdash
  \exists n,j\in\mathbb N\,
    (u=b_{\Succ(\Succ(n)),j})%
}{MogiljanskajaReserveMembership}{%
  \DeltaRow{Mengen}{D\dsep U\dsep u\dsep b_{ij}\dsep n\dsep j}
  \DeltaRow{Funktionen}{b}
}
\begin{tabproofsplitwide}
  \proofpartwide{Die zusätzliche Trägerbedingung ist erfüllt}
    \proofstepwidestar[1]{%
      \exists n,j\in\mathbb N\,(u=b_{\Succ(\Succ(n)),j})}{%
      \rA}
    \proofstepwidestar[2]{%
      n\in\mathbb N\land(j\in\mathbb N\land u=b_{\Succ(\Succ(n)),j})}{%
      \rA}
    \proofstepwidestar[2]{%
      \Succ(\Succ(n))\in\mathbb N}{%
      \FormulaRefAuto{PeanoSuccClosure}[ax]{\FormulaRefAuto{PeanoSuccClosure}[ax]{\rAEa{2}}}}
    \proofstepwidestar[2]{%
      b_{\Succ(\Succ(n)),j}\in D}{%
      \FormulaRefAuto{MogiljanskajaBElementsDef}[def]{3,\rAEa{\rAEb{2}}},\ \FormulaRefAuto{IndexedFamilyCoordinateTyping}[thm]}
    \proofstepwidestar[2]{%
      u\in D}{%
      \rIE{\rAEb{\rAEb{2}},4}}
    \proofstepwidestar[2]{%
      u\in U}{%
      \FormulaRefAuto{MogiljanskajaDPrimeDef}[def]{5,\rEI{\rAI{\rAEa{2},\rEI{\rAEb{2}}}}}}
    \proofstepwidestar[1]{%
      u\in U}{%
      \rEE{1,2,6}}
    \proofstepwidestar[8]{%
      u\in U}{%
      \rA}
    \proofstepwidestar[8]{%
      \exists n,j\in\mathbb N\,(u=b_{\Succ(\Succ(n)),j})}{%
      \FormulaRefAuto{MogiljanskajaDPrimeDef}[def]{8}}
    \proofstepwidestar[]{%
      u\in U\leftrightarrow\exists n,j\in\mathbb N\,(u=b_{\Succ(\Succ(n)),j})}{%
      \rLRI{\rRI{8,9},\rRI{1,7}}}
  \closeproofpartwide
\end{tabproofsplitwide}

\FormulaThmDeltaK[Die Marker und das ausgesparte Rechteckelement]{%
  \begin{aligned}[t]
    &c_0,c_1,c_2\in D\dsep
      c_0,c_1,c_2\notin U\dsep
      c_0,c_1,c_2\text{ paarweise verschieden},\\[-2pt]
    &D'\subseteq D\dsep b_{01}\notin D'
  \end{aligned}%
}{MogiljanskajaReserveMarkerFacts}{%
  \DeltaRow{Mengen}{D\dsep D'\dsep U\dsep V\dsep P\dsep
    c_0\dsep c_1\dsep c_2\dsep b_{01}\dsep n\dsep j}
  \DeltaRow{Funktionen}{b}
}
\begin{tabproofsplitwide}
  \proofpartwideindR[Die ersten beiden Zeilen liegen außerhalb der Reserve]{%
    i\in\{0,1\}\dsep j\in\mathbb N\vdash b_{ij}\notin U
  }{i\in\{0,1\}\dsep j\in\mathbb N\vdash b_{ij}\notin U}[B21ReserveAvoidsFirstTwoRows]
    \proofstepwidestar[1]{%
      n\in\mathbb N}{%
      \rA}
    \proofstepwidestar[1]{%
      \Succ(n)\in\mathbb N}{%
      \FormulaRefAuto{PeanoSuccClosure}[ax]{1}}
    \proofstepwidestar[1]{%
      \Succ(\Succ(n))\neq0}{%
      \FormulaRefAuto{PeanoSuccNonzero}[ax]{2}}
    \proofstepwidestar[4]{%
      \Succ(\Succ(n))=1}{%
      \rA}
    \proofstepwidestar[1,4]{%
      \Succ(n)=0}{%
      \FormulaRefAuto{PeanoSuccInjective}[ax]{2,\FormulaRefAuto{PeanoZero}[ax],4},\ \FormulaRefAuto{PeanoOne}[def]}
    \proofstepwidestar[1,4]{%
      \bot}{%
      \rBI{\FormulaRefAuto{PeanoSuccNonzero}[ax]{1},5}}
    \proofstepwidestar[1]{%
      \Succ(\Succ(n))\neq1}{%
      \rCI{4,6}}
    \proofstepwidestar[]{%
      \forall n\in\mathbb N\,(\Succ(\Succ(n))\neq0\land\Succ(\Succ(n))\neq1)}{%
      \rUI{\rRI{1,\rAI{3,7}}}}
    \proofstepwidestar[9]{%
      i\in\{0,1\}\land j\in\mathbb N}{%
      \rA}
    \proofstepwidestar[9]{%
      i\in\mathbb N\land(i=0\lor i=1)}{%
      \text{Def. }\ensuremath{\{0,1\}},\ 9,\ \FormulaRefAuto{PeanoZero}[ax],\ \FormulaRefAuto{PeanoOneInN}[thm]}
    \proofstepwidestar[11]{%
      b_{ij}\in U}{%
      \rA}
    \proofstepwidestar[11]{%
      \exists n,k\in\mathbb N\,(b_{ij}=b_{\Succ(\Succ(n)),k})}{%
      \FormulaRefAuto{MogiljanskajaReserveMembership}[thm]{11}}
    \proofstepwidestar[13]{%
      n\in\mathbb N\land(k\in\mathbb N\land b_{ij}=b_{\Succ(\Succ(n)),k})}{%
      \rA}
    \proofstepwidestar[9,13]{%
      i=\Succ(\Succ(n))}{%
      \rAEa{\FormulaRefAuto{MogiljanskajaBIndexUnique}[thm]{10,9,13,\FormulaRefAuto{PeanoSuccClosure}[ax]}}}
    \proofstepwidestar[13]{%
      \Succ(\Succ(n))\neq0\land\Succ(\Succ(n))\neq1}{%
      \rRE{\rUE{8},\rAEa{13}}}
    \proofstepwidestar[16]{%
      i=0}{%
      \rA}
    \proofstepwidestar[9,13,16]{%
      \Succ(\Succ(n))=0}{%
      \rIE{14,16}}
    \proofstepwidestar[9,13,16]{%
      \bot}{%
      \rBI{\rAEa{15},17}}
    \proofstepwidestar[19]{%
      i=1}{%
      \rA}
    \proofstepwidestar[9,13,19]{%
      \Succ(\Succ(n))=1}{%
      \rIE{14,19}}
    \proofstepwidestar[9,13,19]{%
      \bot}{%
      \rBI{\rAEb{15},20}}
    \proofstepwidestar[9,13]{%
      \bot}{%
      \rOE{\rAEb{10},16,18,19,21}}
    \proofstepwidestar[9,11]{%
      \bot}{%
      \rEE{12,13,22}}
    \proofstepwidestar[9]{%
      b_{ij}\notin U}{%
      \rCI{11,23}}
  \closeproofpartwideindR

  \proofpartwideindR[Typisierung und Verschiedenheit der Marker]{%
    c_0,c_1,c_2\in D\dsep c_0\neq c_1\land(c_0\neq c_2\land c_1\neq c_2)
  }{c_0,c_1,c_2\in D\dsep c_0\neq c_1\land(c_0\neq c_2\land c_1\neq c_2)}[B21ReserveMarkerDistinct]
    \proofstepwidestar[]{%
      c_0,c_1,c_2\in D}{%
      \FormulaRefAuto{MogiljanskajaDPrimeDef}[def],\ \FormulaRefAuto{MogiljanskajaBElementsDef}[def]{\FormulaRefAuto{PeanoZero}[ax],\FormulaRefAuto{PeanoOneInN}[thm]}}
    \proofstepwidestar[2]{%
      c_0=c_1}{%
      \rA}
    \proofstepwidestar[2]{%
      0=1}{%
      \rAEa{\FormulaRefAuto{MogiljanskajaBIndexUnique}[thm]{\FormulaRefAuto{MogiljanskajaDPrimeDef}[def]{2}}}}
    \proofstepwidestar[2]{%
      \bot}{%
      \rBI{\FormulaRefAuto{PeanoZeroNeOne}[thm],3}}
    \proofstepwidestar[]{%
      c_0\neq c_1}{%
      \rCI{2,4}}
    \proofstepwidestar[6]{%
      c_0=c_2}{%
      \rA}
    \proofstepwidestar[6]{%
      0=1}{%
      \rAEa{\FormulaRefAuto{MogiljanskajaBIndexUnique}[thm]{\FormulaRefAuto{MogiljanskajaDPrimeDef}[def]{6}}}}
    \proofstepwidestar[6]{%
      \bot}{%
      \rBI{\FormulaRefAuto{PeanoZeroNeOne}[thm],7}}
    \proofstepwidestar[]{%
      c_0\neq c_2}{%
      \rCI{6,8}}
    \proofstepwidestar[10]{%
      c_1=c_2}{%
      \rA}
    \proofstepwidestar[10]{%
      1=0}{%
      \rAEb{\FormulaRefAuto{MogiljanskajaBIndexUnique}[thm]{\FormulaRefAuto{MogiljanskajaDPrimeDef}[def]{10}}}}
    \proofstepwidestar[10]{%
      \bot}{%
      \rBI{\FormulaRefAuto{PeanoOneNeZero}[thm],11}}
    \proofstepwidestar[]{%
      c_1\neq c_2}{%
      \rCI{10,12}}
    \proofstepwidestar[]{%
      c_0\neq c_1\land(c_0\neq c_2\land c_1\neq c_2)}{%
      \rAI{5,\rAI{9,13}}}
  \closeproofpartwideindR

  \proofpartwide{Reservenausschluss und Trägerinklusion}
    \proofstepwidestar[]{%
      c_0,c_1,c_2\notin U}{%
      \FormulaRefAuto{B21ReserveAvoidsFirstTwoRows}[thm]{\FormulaRefAuto{PeanoZero}[ax],\FormulaRefAuto{PeanoOneInN}[thm]},\ \FormulaRefAuto{MogiljanskajaDPrimeDef}[def]}
    \proofstepwidestar[]{%
      U\subseteq D}{%
      \FormulaRefAuto{MogiljanskajaDPrimeDef}[def]}
    \proofstepwidestar[]{%
      P\subseteq D\land\{c_2\}\subseteq D}{%
      \FormulaRefAuto{B21ReserveMarkerDistinct}[thm],\ \FormulaRefAuto{MogiljanskajaDPrimeDef}[def],\ \FormulaRefAuto{PairSetSubsetFromMembership}[thm]}
    \proofstepwidestar[]{%
      V\subseteq D}{%
      \FormulaRefAuto{UnionSubsetCommonSuperset}[thm]{2,\rAEb{3}},\ \FormulaRefAuto{MogiljanskajaDPrimeDef}[def]}
    \proofstepwidestar[]{%
      D'\subseteq D}{%
      \FormulaRefAuto{UnionSubsetCommonSuperset}[thm]{\rAEa{3},4},\ \FormulaRefAuto{MogiljanskajaDPrimeDef}[def]}
  \closeproofpartwide

  \proofpartwide{Das ausgesparte Rechteckelement}
    \proofstepwidestar[]{%
      b_{01}\notin U}{%
      \FormulaRefAuto{B21ReserveAvoidsFirstTwoRows}[thm]{\FormulaRefAuto{PeanoZero}[ax],\FormulaRefAuto{PeanoOneInN}[thm]}}
    \proofstepwidestar[2]{%
      b_{01}=c_0}{%
      \rA}
    \proofstepwidestar[2]{%
      1=0}{%
      \rAEb{\FormulaRefAuto{MogiljanskajaBIndexUnique}[thm]{\FormulaRefAuto{MogiljanskajaDPrimeDef}[def]{2}}}}
    \proofstepwidestar[2]{%
      \bot}{%
      \rBI{\FormulaRefAuto{PeanoOneNeZero}[thm],3}}
    \proofstepwidestar[]{%
      b_{01}\neq c_0}{%
      \rCI{2,4}}
    \proofstepwidestar[6]{%
      b_{01}=c_1}{%
      \rA}
    \proofstepwidestar[6]{%
      0=1}{%
      \rAEa{\FormulaRefAuto{MogiljanskajaBIndexUnique}[thm]{\FormulaRefAuto{MogiljanskajaDPrimeDef}[def]{6}}}}
    \proofstepwidestar[6]{%
      \bot}{%
      \rBI{\FormulaRefAuto{PeanoZeroNeOne}[thm],7}}
    \proofstepwidestar[]{%
      b_{01}\neq c_1}{%
      \rCI{6,8}}
    \proofstepwidestar[10]{%
      b_{01}=c_2}{%
      \rA}
    \proofstepwidestar[10]{%
      0=1}{%
      \rAEa{\FormulaRefAuto{MogiljanskajaBIndexUnique}[thm]{\FormulaRefAuto{MogiljanskajaDPrimeDef}[def]{10}}}}
    \proofstepwidestar[10]{%
      \bot}{%
      \rBI{\FormulaRefAuto{PeanoZeroNeOne}[thm],11}}
    \proofstepwidestar[]{%
      b_{01}\neq c_2}{%
      \rCI{10,12}}
    \proofstepwidestar[]{%
      b_{01}\notin P\land b_{01}\notin V}{%
      \FormulaRefAuto{NotInPairSet}[thm]{5,9},\ \FormulaRefAuto{NotInUnion}[thm]{1,13},\ \FormulaRefAuto{MogiljanskajaDPrimeDef}[def]}
    \proofstepwidestar[]{%
      b_{01}\notin D'}{%
      \FormulaRefAuto{NotInUnion}[thm]{\rAEa{14},\rAEb{14}},\ \FormulaRefAuto{MogiljanskajaDPrimeDef}[def]}
  \closeproofpartwide
\end{tabproofsplitwide}

\FormulaThmDeltaK[Der dritte Marker liegt außerhalb des festen Kerns]{%
  c_2\notin P\dsep P\cap U=\varnothing%
}{MogiljanskajaReserveCoreFacts}{%
  \DeltaRow{Mengen}{P\dsep U\dsep c_0\dsep c_1\dsep c_2}
}
\begin{tabproofsplitwide}
  \proofpartwide{Kern und Reserve sind disjunkt}
    \proofstepwidestar[]{%
      c_2\neq c_0\land c_2\neq c_1}{%
      \FormulaRefAuto{B21ReserveMarkerDistinct}[thm],\ \FormulaRefAuto{a=b\vdash b=a}[thm]}
    \proofstepwidestar[]{%
      c_2\notin P}{%
      \FormulaRefAuto{NotInPairSet}[thm]{1},\ \FormulaRefAuto{MogiljanskajaDPrimeDef}[def]}
    \proofstepwidestar[]{%
      c_0\notin U\land c_1\notin U}{%
      \FormulaRefAuto{MogiljanskajaReserveMarkerFacts}[thm]}
    \proofstepwidestar[]{%
      U\cap\{c_0,c_1\}=\varnothing}{%
      \FormulaRefAuto{PairSetDisjointFromTwoAvoided}[thm]{3}}
    \proofstepwidestar[]{%
      P\cap U=\varnothing}{%
      \FormulaRefAuto{MogiljanskajaDPrimeDef}[def],\ \text{Kommutativität von }\ensuremath{\cap}(4)}
  \closeproofpartwide
\end{tabproofsplitwide}

\section{Parametrisierungen und ihre Bildmengen}

\FormulaDefDeltaK[Die beiden Parametrisierungen der Produktschichten]{%
  \begin{aligned}[t]
    &\delta\colon\mathbb N\times\mathbb N\to D,\\[-2pt]
    &\forall n,j\in\mathbb N\,
      \bigl(\delta(n,j)\coloneqq b_{nj}\bigr);\\[2pt]
    &\upsilon\colon\mathbb N\times\mathbb N\to U,\\[-2pt]
    &\forall n,j\in\mathbb N\,
      \bigl(\upsilon(n,j)\coloneqq b_{\Succ(\Succ(n)),j}\bigr)
  \end{aligned}%
}{MogiljanskajaRowParametrizationsDef}{%
  \DeltaRow{Mengen}{D\dsep U\dsep b_{ij}\dsep n\dsep j}
  \DeltaRow{Funktionen}{b\dsep\delta\dsep\upsilon}
}

\FormulaThmDeltaK[Die Zeilenparametrisierungen sind bijektiv]{%
  \delta\colon\mathbb N\times\mathbb N\bij D\dsep
  \upsilon\colon\mathbb N\times\mathbb N\bij U%
}{MogiljanskajaRowParametrizationsBijective}{%
  \DeltaRow{Mengen}{D\dsep U\dsep n\dsep j\dsep m\dsep q\dsep u}
  \DeltaRow{Funktionen}{\delta\dsep\upsilon}
}
\begin{tabproofsplitwide}
  \proofpartwideindR[Unverschobene Zeilen]{%
    \delta\colon\mathbb N\times\mathbb N\bij D
  }{\delta\colon\mathbb N\times\mathbb N\bij D}[B21DeltaBijective]
    \proofstepwidestar[]{%
      \delta\colon\mathbb N\times\mathbb N\to D}{%
      \FormulaRefAuto{MogiljanskajaRowParametrizationsDef}[def]}
    \proofstepwidestar[2]{%
      n,j,m,q\in\mathbb N}{%
      \rA}
    \proofstepwidestar[3]{%
      \delta(n,j)=\delta(m,q)}{%
      \rA}
    \proofstepwidestar[2,3]{%
      b_{n,j}=b_{m,q}}{%
      \FormulaRefAuto{MogiljanskajaRowParametrizationsDef}[def]{2,3}}
    \proofstepwidestar[2,3]{%
      n=m\land j=q}{%
      \FormulaRefAuto{MogiljanskajaBIndexUnique}[thm]{2,4,\FormulaRefAuto{PeanoSuccClosure}[ax]}}
    \proofstepwidestar[2,3]{%
      n=m\land j=q}{%
      5}
    \proofstepwidestar[2,3]{%
      (n,j)=(m,q)}{%
      \FormulaRefAuto{(a,b)=(c,d)\eqvdash a=c\land b=d}[thm]{6}}
    \proofstepwidestar[]{%
      \delta\colon\mathbb N\times\mathbb N\inj D}{%
      \FormulaRefAuto{Injektive Funktion}[def]{1,\rUI{\rUI{\rUI{\rUI{\rRI{2,\rRI{3,7}}}}}}}}
    \proofstepwidestar[9]{%
      u\in D}{%
      \rA}
    \proofstepwidestar[9]{%
      \exists n,j\in\mathbb N\,(u=b_{n,j})}{%
      \FormulaRefAuto{MogiljanskajaLayerDMembership}[thm]{9}}
    \proofstepwidestar[11]{%
      n\in\mathbb N\land(j\in\mathbb N\land u=b_{n,j})}{%
      \rA}
    \proofstepwidestar[11]{%
      (n,j)\in\mathbb N\times\mathbb N\land \delta(n,j)=u}{%
      \FormulaRefAuto{MogiljanskajaRowParametrizationsDef}[def]{11},\ \text{Def. }\ensuremath{\times},\ \FormulaRefAuto{a=b\vdash b=a}[thm]}
    \proofstepwidestar[11]{%
      \exists p\in\mathbb N\times\mathbb N\,(\delta(p)=u)}{%
      \rEI{12}}
    \proofstepwidestar[9]{%
      \exists p\in\mathbb N\times\mathbb N\,(\delta(p)=u)}{%
      \rEE{10,11,13}}
    \proofstepwidestar[]{%
      \delta\colon\mathbb N\times\mathbb N\sur D}{%
      \FormulaRefAuto{Surjektive Funktion}[def]{1,\rUI{\rRI{9,14}}}}
    \proofstepwidestar[]{%
      \delta\colon\mathbb N\times\mathbb N\bij D}{%
      \FormulaRefAuto{Bijektive Funktion}[def]{8,15}}
  \closeproofpartwideindR

  \proofpartwideindR[Verschobene Zeilen]{%
    \upsilon\colon\mathbb N\times\mathbb N\bij U
  }{\upsilon\colon\mathbb N\times\mathbb N\bij U}[B21UpsilonBijective]
    \proofstepwidestar[]{%
      \upsilon\colon\mathbb N\times\mathbb N\to U}{%
      \FormulaRefAuto{MogiljanskajaRowParametrizationsDef}[def]}
    \proofstepwidestar[2]{%
      n,j,m,q\in\mathbb N}{%
      \rA}
    \proofstepwidestar[3]{%
      \upsilon(n,j)=\upsilon(m,q)}{%
      \rA}
    \proofstepwidestar[2,3]{%
      b_{\Succ(\Succ(n)),j}=b_{\Succ(\Succ(m)),q}}{%
      \FormulaRefAuto{MogiljanskajaRowParametrizationsDef}[def]{2,3}}
    \proofstepwidestar[2,3]{%
      \Succ(\Succ(n))=\Succ(\Succ(m))\land j=q}{%
      \FormulaRefAuto{MogiljanskajaBIndexUnique}[thm]{2,4,\FormulaRefAuto{PeanoSuccClosure}[ax]}}
    \proofstepwidestar[2,3]{%
      n=m\land j=q}{%
      \FormulaRefAuto{PeanoSuccInjective}[ax]{\FormulaRefAuto{PeanoSuccInjective}[ax]{\rAEa{5}}},\ \rAEb{5}}
    \proofstepwidestar[2,3]{%
      (n,j)=(m,q)}{%
      \FormulaRefAuto{(a,b)=(c,d)\eqvdash a=c\land b=d}[thm]{6}}
    \proofstepwidestar[]{%
      \upsilon\colon\mathbb N\times\mathbb N\inj U}{%
      \FormulaRefAuto{Injektive Funktion}[def]{1,\rUI{\rUI{\rUI{\rUI{\rRI{2,\rRI{3,7}}}}}}}}
    \proofstepwidestar[9]{%
      u\in U}{%
      \rA}
    \proofstepwidestar[9]{%
      \exists n,j\in\mathbb N\,(u=b_{\Succ(\Succ(n)),j})}{%
      \FormulaRefAuto{MogiljanskajaReserveMembership}[thm]{9}}
    \proofstepwidestar[11]{%
      n\in\mathbb N\land(j\in\mathbb N\land u=b_{\Succ(\Succ(n)),j})}{%
      \rA}
    \proofstepwidestar[11]{%
      (n,j)\in\mathbb N\times\mathbb N\land \upsilon(n,j)=u}{%
      \FormulaRefAuto{MogiljanskajaRowParametrizationsDef}[def]{11},\ \text{Def. }\ensuremath{\times},\ \FormulaRefAuto{a=b\vdash b=a}[thm]}
    \proofstepwidestar[11]{%
      \exists p\in\mathbb N\times\mathbb N\,(\upsilon(p)=u)}{%
      \rEI{12}}
    \proofstepwidestar[9]{%
      \exists p\in\mathbb N\times\mathbb N\,(\upsilon(p)=u)}{%
      \rEE{10,11,13}}
    \proofstepwidestar[]{%
      \upsilon\colon\mathbb N\times\mathbb N\sur U}{%
      \FormulaRefAuto{Surjektive Funktion}[def]{1,\rUI{\rRI{9,14}}}}
    \proofstepwidestar[]{%
      \upsilon\colon\mathbb N\times\mathbb N\bij U}{%
      \FormulaRefAuto{Bijektive Funktion}[def]{8,15}}
  \closeproofpartwideindR
\end{tabproofsplitwide}

Insbesondere gelten die präzisen Bildmengengleichheiten
\[
  \delta[\mathbb N\times\mathbb N]=D,
  \qquad
  \upsilon[\mathbb N\times\mathbb N]=U.
\]

\FormulaDefDeltaK[Die Zeilenrückverschiebung]{%
  \theta\coloneqq\delta\circ\upsilon^{-1}%
}{MogiljanskajaThetaDef}{%
  \DeltaRow{Mengen}{D\dsep U}
  \DeltaRow{Funktionen}{\delta\dsep\upsilon\dsep\theta}
}

\FormulaThmDeltaK[Die Zeilenrückverschiebung ist bijektiv]{%
  \theta\colon U\bij D%
}{MogiljanskajaThetaBijective}{%
  \DeltaRow{Mengen}{D\dsep U}
  \DeltaRow{Funktionen}{\delta\dsep\upsilon\dsep\theta}
}
\begin{tabproofsplitwide}
  \proofpartwide{Inverse und Komposition}
    \proofstepwidestar[]{%
      \upsilon\colon\mathbb N\times\mathbb N\bij U}{%
      \FormulaRefAuto{B21UpsilonBijective}[thm]}
    \proofstepwidestar[]{%
      \upsilon^{-1}\colon U\bij\mathbb N\times\mathbb N}{%
      \FormulaRefAuto{F\colon A\bij B\vdash F^{-1}\colon B\bij A}[thm]{1}}
    \proofstepwidestar[]{%
      \delta\colon\mathbb N\times\mathbb N\bij D}{%
      \FormulaRefAuto{B21DeltaBijective}[thm]}
    \proofstepwidestar[]{%
      \delta\circ\upsilon^{-1}\colon U\bij D}{%
      \FormulaRefAuto{BijectiveFunctionComposition}[thm]{2,3}}
    \proofstepwidestar[]{%
      \theta\colon U\bij D}{%
      \FormulaRefAuto{MogiljanskajaThetaDef}[def],\ \rIE{4}}
  \closeproofpartwide
\end{tabproofsplitwide}

\section{Die punktierte Verschiebung}

\FormulaThmDeltaK[Die Werte der Verschiebungsfolge liegen in der Reserve]{%
  n\in\mathbb N\vdash b_{\Succ(n),0}\in V%
}{MogiljanskajaShiftValueInV}{%
  \DeltaRow{Mengen}{U\dsep V\dsep c_2\dsep b_{ij}\dsep n}
  \DeltaRow{Funktionen}{b}
}
\begin{tabproofsplitwide}
  \proofpartwide{Nullindex und Nachfolgerindex}
    \proofstepwidestar[1]{%
      n\in\mathbb N}{%
      \rA}
    \proofstepwidestar[1]{%
      n=0\lor\exists m\in\mathbb N\,(n=m+1)}{%
      \FormulaRefAuto{PeanoIndexSplitAtOne}[thm]{1}}
    \proofstepwidestar[3]{%
      n=0}{%
      \rA}
    \proofstepwidestar[3]{%
      b_{\Succ(n),0}=c_2}{%
      \rIE{3},\ \FormulaRefAuto{PeanoOne}[def],\ \FormulaRefAuto{MogiljanskajaDPrimeDef}[def]}
    \proofstepwidestar[3]{%
      b_{\Succ(n),0}\in V}{%
      \FormulaRefAuto{MogiljanskajaDPrimeDef}[def]{4},\ \text{Def. }\ensuremath{\cup}}
    \proofstepwidestar[6]{%
      \exists m\in\mathbb N\,(n=m+1)}{%
      \rA}
    \proofstepwidestar[7]{%
      m\in\mathbb N\land n=m+1}{%
      \rA}
    \proofstepwidestar[7]{%
      n=\Succ(m)}{%
      \FormulaRefAuto{PeanoAddOneSucc}[thm]{\rAEa{7}},\ \rIE{\rAEb{7}}}
    \proofstepwidestar[7]{%
      b_{\Succ(n),0}\in U}{%
      \FormulaRefAuto{MogiljanskajaReserveMembership}[thm]{\rEI{\rAI{\rAEa{7},\rEI{\rAI{\FormulaRefAuto{PeanoZero}[ax],\rIE{8,\rII}}}}}}}
    \proofstepwidestar[7]{%
      b_{\Succ(n),0}\in V}{%
      \FormulaRefAuto{MogiljanskajaDPrimeDef}[def]{9},\ \text{Def. }\ensuremath{\cup}}
    \proofstepwidestar[6]{%
      b_{\Succ(n),0}\in V}{%
      \rEE{6,7,10}}
    \proofstepwidestar[1]{%
      b_{\Succ(n),0}\in V}{%
      \rOE{2,3,5,6,11}}
  \closeproofpartwide
\end{tabproofsplitwide}

\FormulaDefDeltaK[Die Folge für die punktierte Verschiebung]{%
  \begin{aligned}[t]
    &H\colon\mathbb N\to V,\\[-2pt]
    &\forall n\in\mathbb N\,
      \bigl(H_n\coloneqq b_{\Succ(n),0}\bigr)
  \end{aligned}%
}{MogiljanskajaShiftSequenceDef}{%
  \DeltaRow{Mengen}{V\dsep b_{ij}\dsep n}
  \DeltaRow{Funktionen}{b\dsep H}
}

\FormulaThmDeltaK[Die punktierte Verschiebung entfernt genau den dritten
Marker]{%
  H\colon\mathbb N\inj V\dsep H_0=c_2\dsep
  V\setminus\{c_2\}=U%
}{MogiljanskajaShiftSequenceFacts}{%
  \DeltaRow{Mengen}{U\dsep V\dsep c_2\dsep n\dsep m}
  \DeltaRow{Funktionen}{H}
}
\begin{tabproofsplitwide}
  \proofpartwideindR[Injektivität der Verschiebungsfolge]{%
    H\colon\mathbb N\inj V
  }{H\colon\mathbb N\inj V}[B21ShiftSequenceInjective]
    \proofstepwidestar[]{%
      H\colon\mathbb N\to V}{%
      \FormulaRefAuto{MogiljanskajaShiftSequenceDef}[def]}
    \proofstepwidestar[2]{%
      n,m\in\mathbb N}{%
      \rA}
    \proofstepwidestar[3]{%
      H_n=H_m}{%
      \rA}
    \proofstepwidestar[2,3]{%
      b_{\Succ(n),0}=b_{\Succ(m),0}}{%
      \FormulaRefAuto{MogiljanskajaShiftSequenceDef}[def]{2,3}}
    \proofstepwidestar[2,3]{%
      \Succ(n)=\Succ(m)}{%
      \rAEa{\FormulaRefAuto{MogiljanskajaBIndexUnique}[thm]{2,\FormulaRefAuto{PeanoSuccClosure}[ax],\FormulaRefAuto{PeanoZero}[ax],4}}}
    \proofstepwidestar[2,3]{%
      n=m}{%
      \FormulaRefAuto{PeanoSuccInjective}[ax]{2,5}}
    \proofstepwidestar[]{%
      H\colon\mathbb N\inj V}{%
      \FormulaRefAuto{InjectiveIndexedFamilyCriterion}[thm]{1,\rUI{\rUI{\rRI{2,\rRI{3,6}}}}}}
  \closeproofpartwideindR

  \proofpartwide{Anfangswert und punktierte Zielmenge}
    \proofstepwidestar[]{%
      H_0=b_{\Succ(0),0}=c_2}{%
      \FormulaRefAuto{MogiljanskajaShiftSequenceDef}[def]{\FormulaRefAuto{PeanoZero}[ax]},\ \FormulaRefAuto{PeanoOne}[def],\ \FormulaRefAuto{MogiljanskajaDPrimeDef}[def]}
    \proofstepwidestar[]{%
      c_2\notin U}{%
      \FormulaRefAuto{MogiljanskajaReserveMarkerFacts}[thm]}
    \proofstepwidestar[]{%
      \{c_2\}\cap U=\varnothing}{%
      \FormulaRefAuto{SingletonDisjointFromAvoidedSet}[thm]{2}}
    \proofstepwidestar[]{%
      (\{c_2\}\cup U)\setminus\{c_2\}=U}{%
      \FormulaRefAuto{RemoveDisjointUnionSummand}[thm]{3}}
    \proofstepwidestar[]{%
      V\setminus\{c_2\}=U}{%
      \FormulaRefAuto{MogiljanskajaDPrimeDef}[def],\ \text{Kommutativität von }\ensuremath{\cup}(4)}
  \closeproofpartwide
\end{tabproofsplitwide}

\FormulaThmDeltaK[Tatsächliche Bildmenge der Verschiebungsfolge]{%
  H[\mathbb N]
  =\{b_{\Succ(n),0}\mid n\in\mathbb N\}
  \dsep H[\mathbb N]\subsetneq V%
}{MogiljanskajaShiftSequenceImage}{%
  \DeltaRow{Mengen}{V\dsep b_{ij}\dsep n}
  \DeltaRow{Funktionen}{H}
}
\begin{tabproofsplitwide}
  \proofpartwide{Bildgleichheit und Inklusion}
    \proofstepwidestar[]{%
      H\colon\mathbb N\to V}{%
      \FormulaRefAuto{MogiljanskajaShiftSequenceDef}[def]}
    \proofstepwidestar[]{%
      x\in H[\mathbb N]\leftrightarrow\exists n\in\mathbb N\,(x=b_{\Succ(n),0})}{%
      \FormulaRefAuto{IndexedFamilyImageMembership}[thm]{1},\ \FormulaRefAuto{IndexedFamilyImageDef}[def],\ \FormulaRefAuto{MogiljanskajaShiftSequenceDef}[def]}
    \proofstepwidestar[]{%
      H[\mathbb N]=\{b_{\Succ(n),0}\mid n\in\mathbb N\}}{%
      \text{Extensionalität}(2),\ \FormulaRefAuto{TermImageSetUniqueExists}[thm]}
    \proofstepwidestar[]{%
      H[\mathbb N]\subseteq V}{%
      \FormulaRefAuto{F\colon A\to B\vdash F[A]\subseteq B}[thm]{1}}
  \closeproofpartwide

  \proofpartwide{Ein explizit ausgelassener Wert}
    \proofstepwidestar[]{%
      b_{\Succ(\Succ(0)),1}\in U}{%
      \FormulaRefAuto{MogiljanskajaReserveMembership}[thm]{\rEI{\rAI{\FormulaRefAuto{PeanoZero}[ax],\rEI{\rAI{\FormulaRefAuto{PeanoOneInN}[thm],\rII}}}}}}
    \proofstepwidestar[]{%
      b_{\Succ(\Succ(0)),1}\in V}{%
      \FormulaRefAuto{MogiljanskajaDPrimeDef}[def]{1},\ \text{Def. }\ensuremath{\cup}}
    \proofstepwidestar[3]{%
      b_{\Succ(\Succ(0)),1}\in H[\mathbb N]}{%
      \rA}
    \proofstepwidestar[3]{%
      \exists n\in\mathbb N\,(b_{\Succ(\Succ(0)),1}=H_n)}{%
      \FormulaRefAuto{IndexedFamilyImageMembership}[thm]{\FormulaRefAuto{MogiljanskajaShiftSequenceDef}[def],3},\ \FormulaRefAuto{IndexedFamilyImageDef}[def]}
    \proofstepwidestar[5]{%
      n\in\mathbb N\land b_{\Succ(\Succ(0)),1}=H_n}{%
      \rA}
    \proofstepwidestar[5]{%
      b_{\Succ(\Succ(0)),1}=b_{\Succ(n),0}}{%
      \FormulaRefAuto{MogiljanskajaShiftSequenceDef}[def]{\rAEa{5}},\ \rIE{\rAEb{5}}}
    \proofstepwidestar[5]{%
      1=0}{%
      \rAEb{\FormulaRefAuto{MogiljanskajaBIndexUnique}[thm]{\FormulaRefAuto{PeanoZero}[ax],\FormulaRefAuto{PeanoOneInN}[thm],\FormulaRefAuto{PeanoSuccClosure}[ax],5,6}}}
    \proofstepwidestar[5]{%
      \bot}{%
      \rBI{\FormulaRefAuto{PeanoOneNeZero}[thm],7}}
    \proofstepwidestar[3]{%
      \bot}{%
      \rEE{4,5,8}}
    \proofstepwidestar[]{%
      b_{\Succ(\Succ(0)),1}\notin H[\mathbb N]}{%
      \rCI{3,9}}
    \proofstepwidestar[]{%
      H[\mathbb N]\subseteq V}{%
      \FormulaRefAuto{IndexedFamilyImageSubsetCodomain}[thm]{\FormulaRefAuto{MogiljanskajaShiftSequenceDef}[def]},\ \FormulaRefAuto{IndexedFamilyImageDef}[def]}
    \proofstepwidestar[]{%
      H[\mathbb N]\subsetneq V}{%
      11,\ 2,\ 10,\ \text{Def. }\ensuremath{\subsetneq}}
  \closeproofpartwide
\end{tabproofsplitwide}

Damit ist ausdrücklich zwischen Ziel- und Bildmenge unterschieden:
\(H\colon\mathbb N\to V\), aber \(H[\mathbb N]\neq V\). Die Folge wird
nicht zur Aufzählung von \(V\), sondern zur Konstruktion einer Verschiebung
auf \(V\) verwendet.

\FormulaDefDeltaK[Die Bijektion in die um einen Marker erweiterte Reserve]{%
  \sigma\coloneqq(\SuccShift{H})^{-1}%
}{MogiljanskajaSigmaDef}{%
  \DeltaRow{Mengen}{U\dsep V}
  \DeltaRow{Funktionen}{H\dsep\SuccShift{H}\dsep\sigma}
}

\FormulaThmDeltaK[Die Markeradjunktion ist bijektiv]{%
  \sigma\colon U\bij V%
}{MogiljanskajaSigmaBijective}{%
  \DeltaRow{Mengen}{U\dsep V\dsep c_2}
  \DeltaRow{Funktionen}{H\dsep\SuccShift{H}\dsep\sigma}
}
\begin{tabproofsplitwide}
  \proofpartwide{Inverse der punktierten Verschiebung}
    \proofstepwidestar[]{%
      H\colon\mathbb N\inj V}{%
      \FormulaRefAuto{B21ShiftSequenceInjective}[thm]}
    \proofstepwidestar[]{%
      H_0=c_2\land V\setminus\{c_2\}=U}{%
      \FormulaRefAuto{MogiljanskajaShiftSequenceFacts}[thm]}
    \proofstepwidestar[]{%
      \SuccShift{H}\colon V\bij V\setminus\{H(0)\}}{%
      \FormulaRefAuto{SuccShiftPuncturedBijection}[thm]{1}}
    \proofstepwidestar[]{%
      \SuccShift{H}\colon V\bij U}{%
      \FormulaRefAuto{IndexedFamilyDef}[def]{\FormulaRefAuto{MogiljanskajaShiftSequenceDef}[def],\FormulaRefAuto{PeanoZero}[ax]},\ \rIE{2,3}}
    \proofstepwidestar[]{%
      (\SuccShift{H})^{-1}\colon U\bij V}{%
      \FormulaRefAuto{F\colon A\bij B\vdash F^{-1}\colon B\bij A}[thm]{4}}
    \proofstepwidestar[]{%
      \sigma\colon U\bij V}{%
      \FormulaRefAuto{MogiljanskajaSigmaDef}[def],\ \rIE{5}}
  \closeproofpartwide
\end{tabproofsplitwide}

\section{Familien mit den vorgegebenen Bildmengen}

Die Familien \(a\), \(b\), \(\delta\) und \(\upsilon\) parametrisieren
bereits \(L,D,D\) beziehungsweise \(U\). Aus \(\sigma\) entsteht ohne eine
weitere Fallunterscheidung eine bijektive Familie mit Bild \(V\).

\FormulaDefDeltaK[Eine bijektive Familie mit Bildmenge (V)]{%
  \nu\coloneqq\sigma\circ\upsilon%
}{MogiljanskajaVParametrizationDef}{%
  \DeltaRow{Mengen}{U\dsep V}
  \DeltaRow{Funktionen}{\upsilon\dsep\sigma\dsep\nu}
}

\FormulaThmDeltaK[Die Familie (nu) parametrisiert (V)]{%
  \nu\colon\mathbb N\times\mathbb N\bij V\dsep
  \nu[\mathbb N\times\mathbb N]=V%
}{MogiljanskajaVParametrizationBijective}{%
  \DeltaRow{Mengen}{V}
  \DeltaRow{Funktionen}{\nu}
}
\begin{tabproofsplitwide}
  \proofpartwide{Komposition und Bildmenge}
    \proofstepwidestar[]{%
      \upsilon\colon\mathbb N\times\mathbb N\bij U}{%
      \FormulaRefAuto{B21UpsilonBijective}[thm]}
    \proofstepwidestar[]{%
      \sigma\colon U\bij V}{%
      \FormulaRefAuto{MogiljanskajaSigmaBijective}[thm]}
    \proofstepwidestar[]{%
      \sigma\circ\upsilon\colon\mathbb N\times\mathbb N\bij V}{%
      \FormulaRefAuto{BijectiveFunctionComposition}[thm]{1,2}}
    \proofstepwidestar[]{%
      \nu\colon\mathbb N\times\mathbb N\bij V}{%
      \FormulaRefAuto{MogiljanskajaVParametrizationDef}[def],\ \rIE{3}}
    \proofstepwidestar[]{%
      \nu\colon\mathbb N\times\mathbb N\sur V}{%
      \FormulaRefAuto{Bijektive Funktion}[def]{4}}
    \proofstepwidestar[]{%
      \nu[\mathbb N\times\mathbb N]=V}{%
      \FormulaRefAuto{SurjectiveImageEqualsCodomain}[thm]{5}}
  \closeproofpartwide
\end{tabproofsplitwide}

Für \(D'=P\cup V\) müssen die beiden Marker aus \(P\) und die Familie
\(\nu\) disjunkt zusammengefügt werden. Die Tags \(0\) und \(1\) sorgen
dafür, dass sich die beiden Indexteile nicht überschneiden.

\FormulaDefDeltaK[Indexmenge für die Familie mit Bild (D')]{%
  I_{D'}\coloneqq
  (\{0\}\times\NatSeg{2})
  \cup
  (\{1\}\times(\mathbb N\times\mathbb N))%
}{MogiljanskajaDPrimeIndexSetDef}{%
  \DeltaRow{Mengen}{I_{D'}}
}

\FormulaDefDeltaK[Eine Familie mit Bildmenge (D')]{%
  \begin{aligned}[t]
    &\omega\colon I_{D'}\to D',\\[-2pt]
    &\forall i\in\NatSeg{2}\,
      \left(
        \omega((0,i))\coloneqq
        \begin{cases}
          c_0,& i=0,\\
          c_1,& i=1,
        \end{cases}
      \right),\\[-2pt]
    &\forall n,j\in\mathbb N\,
      \bigl(\omega((1,(n,j)))\coloneqq\nu(n,j)\bigr)
  \end{aligned}%
}{MogiljanskajaDPrimeParametrizationDef}{%
  \DeltaRow{Mengen}{I_{D'}\dsep D'\dsep c_0\dsep c_1\dsep n\dsep j}
  \DeltaRow{Funktionen}{\nu\dsep\omega}
}

\FormulaThmDeltaK[Die Familie (omega) parametrisiert (D')]{%
  \omega\colon I_{D'}\bij D'\dsep \omega[I_{D'}]=D'%
}{MogiljanskajaDPrimeParametrizationBijective}{%
  \DeltaRow{Mengen}{I_{D'}\dsep D'\dsep P\dsep V}
  \DeltaRow{Funktionen}{\omega}
}
\begin{tabproofsplitwide}
  \proofpartwideindR[Die beiden Indexteile]{%
    \omega\colon I_{D'}\bij D'
  }{\omega\colon I_{D'}\bij D'}[B21DPrimeFamilyBijective]
    \proofstepwidestar[]{%
      I_0\coloneqq\{0\}\times\NatSeg 2}{%
      \text{Abkürzung}}
    \proofstepwidestar[]{%
      I_1\coloneqq\{1\}\times(\mathbb N\times\mathbb N)}{%
      \text{Abkürzung}}
    \proofstepwidestar[]{%
      I_{D'}=I_0\cup I_1}{%
      \FormulaRefAuto{MogiljanskajaDPrimeIndexSetDef}[def],\ 1,\ 2}
    \proofstepwidestar[]{%
      I_0\cap I_1=\varnothing}{%
      \FormulaRefAuto{TaggedProductLayersDisjoint}[thm]{\FormulaRefAuto{PeanoZeroNeOne}[thm]},\ 1,\ 2}
    \proofstepwidestar[]{%
      \omega\colon I_{D'}\to D'}{%
      \FormulaRefAuto{MogiljanskajaDPrimeParametrizationDef}[def]}
    \proofstepwidestar[]{%
      I_0=\{(0,0),(0,1)\}}{%
      \FormulaRefAuto{TaggedProductLayerRepresentation}[thm],\ \FormulaRefAuto{PeanoNatSegTwoMembership}[thm]{\FormulaRefAuto{PeanoTwoInN}[thm]},\ 1,\ \text{Extensionalität}}
    \proofstepwidestar[]{%
      \omega(0,0)=c_0\land\omega(0,1)=c_1}{%
      \FormulaRefAuto{MogiljanskajaDPrimeParametrizationDef}[def]{\FormulaRefAuto{PeanoNatSegTwoMembership}[thm]}}
    \proofstepwidestar[]{%
      c_0\neq c_1}{%
      \FormulaRefAuto{B21ReserveMarkerDistinct}[thm]}
    \proofstepwidestar[]{%
      \omega\restriction_{I_0}\colon I_0\bij P}{%
      \FormulaRefAuto{RestrictionDef}[def]{5,3},\ \FormulaRefAuto{MogiljanskajaDPrimeDef}[def],\ 6,\ 7,\ 8,\ \FormulaRefAuto{Bijektive Funktion}[def]}
    \proofstepwidestar[]{%
      \nu\colon\mathbb N\times\mathbb N\bij V}{%
      \FormulaRefAuto{MogiljanskajaVParametrizationBijective}[thm]}
    \proofstepwidestar[]{%
      \forall x\in I_1\,\exists!p\in\mathbb N\times\mathbb N\,(x=(1,p))}{%
      \FormulaRefAuto{TaggedProductLayerUniqueRepresentation}[thm],\ 2}
    \proofstepwidestar[]{%
      \forall p\in\mathbb N\times\mathbb N\,(\omega(1,p)=\nu(p))}{%
      \FormulaRefAuto{MogiljanskajaDPrimeParametrizationDef}[def],\ \text{Def. }\ensuremath{\times}}
    \proofstepwidestar[]{%
      \omega\restriction_{I_1}\colon I_1\bij V}{%
      \FormulaRefAuto{RestrictionDef}[def]{5,3},\ 10,\ 11,\ 12,\ \FormulaRefAuto{Bijektive Funktion}[def]}
    \proofstepwidestar[]{%
      P\cap V=\varnothing}{%
      \FormulaRefAuto{MogiljanskajaReserveCoreFacts}[thm],\ \FormulaRefAuto{MogiljanskajaDPrimeDef}[def],\ \text{Def. }\ensuremath{\cap,\cup}}
    \proofstepwidestar[]{%
      \omega\colon I_{D'}\bij D'}{%
      \FormulaRefAuto{CaseFunBijectiveDisjointTargets}[thm]{9,13,4,14},\ \FormulaRefAuto{MogiljanskajaDPrimeDef}[def],\ 3,\ 7,\ 12}
  \closeproofpartwideindR

  \proofpartwide{Bildmenge}
    \proofstepwidestar[]{%
      \omega\colon I_{D'}\bij D'}{%
      \FormulaRefAuto{B21DPrimeFamilyBijective}[thm]}
    \proofstepwidestar[]{%
      \omega\colon I_{D'}\sur D'}{%
      \FormulaRefAuto{Bijektive Funktion}[def]{1}}
    \proofstepwidestar[]{%
      \omega[I_{D'}]=D'}{%
      \FormulaRefAuto{SurjectiveImageEqualsCodomain}[thm]{2}}
  \closeproofpartwide
\end{tabproofsplitwide}

\FormulaThmDeltaK[Mogiljanskaja-Bildmengen]{%
  \begin{aligned}[t]
    a[\mathbb N]&=L,&
    b[\mathbb N\times\mathbb N]&=D,\\[-2pt]
    \upsilon[\mathbb N\times\mathbb N]&=U,&
    \nu[\mathbb N\times\mathbb N]&=V,\\[-2pt]
    \omega[I_{D'}]&=D'
  \end{aligned}%
}{MogiljanskajaFamilyImages}{%
  \DeltaRow{Mengen}{L\dsep D\dsep U\dsep V\dsep D'\dsep I_{D'}}
  \DeltaRow{Funktionen}{a\dsep b\dsep\upsilon\dsep\nu\dsep\omega}
}
\begin{tabproofsplitwide}
  \proofpartwide{Zusammenstellung der fünf Bildmengen}
    \proofstepwidestar[]{%
      a[\mathbb N]=L}{%
      \FormulaRefAuto{MogiljanskajaLayerLImage}[thm]}
    \proofstepwidestar[]{%
      b[\mathbb N\times\mathbb N]=D}{%
      \FormulaRefAuto{MogiljanskajaLayerDImage}[thm]}
    \proofstepwidestar[]{%
      \upsilon\colon\mathbb N\times\mathbb N\bij U}{%
      \FormulaRefAuto{B21UpsilonBijective}[thm]}
    \proofstepwidestar[]{%
      \upsilon[\mathbb N\times\mathbb N]=U}{%
      \FormulaRefAuto{SurjectiveImageEqualsCodomain}[thm]{\FormulaRefAuto{Bijektive Funktion}[def]{3}}}
    \proofstepwidestar[]{%
      \nu[\mathbb N\times\mathbb N]=V}{%
      \FormulaRefAuto{MogiljanskajaVParametrizationBijective}[thm]}
    \proofstepwidestar[]{%
      \omega[I_{D'}]=D'}{%
      \FormulaRefAuto{MogiljanskajaDPrimeParametrizationBijective}[thm]}
  \closeproofpartwide
\end{tabproofsplitwide}

Diese Zusammenfassung ist die Schnittstelle zum späteren Halbgruppenband.
Dort müssen die Familien und ihre Typisierung nicht erneut aufgebaut werden;
es genügt, ihre Bildmengengleichheiten und die wenigen benötigten
Markerfakten zu zitieren.

% Die bandlokalen Tabellenanpassungen dürfen beim Gesamtsatz nicht in die
% nachfolgenden Bände hineinwirken.
\let\proofstepwidestar\BandTwentyOneOriginalProofStepWideStar
\let\RuleRefWithArg\BandTwentyOneOriginalRuleRefWithArg
\ExplSyntaxOn
\cs_set_eq:NN
  \mx_formula_ref_suffix:n
  \mx_band_twentyone_original_formula_ref_suffix:n
\ExplSyntaxOff
\makeatletter
\newcol@{G}[0]{>{\scriptsize}r}
\makeatother

\end{document}
